% Generated by build_notes.py from modules/*.md. Edit those sources.
\documentclass[11pt,oneside]{report}
\usepackage[letterpaper,margin=1in,headheight=15pt]{geometry}
\usepackage{fontspec}
\setmainfont{texgyrepagella-regular.otf}[BoldFont=texgyrepagella-bold.otf,ItalicFont=texgyrepagella-italic.otf,BoldItalicFont=texgyrepagella-bolditalic.otf]
\setsansfont{texgyreheros-regular.otf}[BoldFont=texgyreheros-bold.otf,ItalicFont=texgyreheros-italic.otf]
\setmonofont{lmmono10-regular.otf}[Scale=0.85]
\usepackage{amsmath,amssymb,amsthm,booktabs,longtable,array,multirow,calc}
\usepackage{graphicx,xcolor,tikz,pgfplots}
\pgfplotsset{compat=1.18}
\usetikzlibrary{arrows.meta,positioning,calc,fit,shapes.geometric,patterns,decorations.pathreplacing}
% UIUC official brand colors, verified 2026-10-23.
% https://brand.illinois.edu/visual-identity/color
\definecolor{IlliniOrange}{HTML}{FF5F05}
\definecolor{IlliniBlue}{HTML}{13294B}
\definecolor{Storm}{HTML}{707372}
\definecolor{LightStorm}{HTML}{C6C7C6}
\definecolor{Industrial}{HTML}{1D58A7}
\definecolor{Patina}{HTML}{007E8E}
\tikzset{>=Stealth, every node/.style={font=\small,text=IlliniBlue}, box/.style={draw=IlliniBlue,rounded corners=2pt,align=center,inner sep=6pt,execute at begin node=\hyphenpenalty10000\relax}, orangebox/.style={box,draw=IlliniOrange,very thick}, edge/.style={->,thick,IlliniBlue}, accent/.style={->,very thick,IlliniOrange}}
\pgfplotsset{every axis/.append style={width=13.5cm,height=6.5cm,axis line style={IlliniBlue},tick label style={font=\small,text=IlliniBlue},label style={font=\small,text=IlliniBlue},legend style={font=\small,draw=none,fill=white},grid style={LightStorm},major grid style={LightStorm},axis lines=left,clip=false}, every axis plot/.append style={thick}}

\usepackage{titlesec,fancyhdr,enumitem,caption,float,placeins,fvextra}
\usepackage{needspace}
\usepackage{microtype,xurl}
\usepackage{hyperref}
\hypersetup{colorlinks=true,linkcolor=IlliniBlue,urlcolor=Industrial,citecolor=IlliniBlue,pdftitle={BADM 534: Data Analytics for Marketing Applications},pdfauthor={Nathan Yang}}
\usepackage{bookmark}
\titleformat{\chapter}[display]{\normalfont\huge\bfseries\raggedright\hyphenpenalty=10000\color{IlliniBlue}}{\chaptertitlename\ \thechapter}{8pt}{\huge}
\titlespacing*{\chapter}{0pt}{-18pt}{20pt}
\titleformat{\section}{\Large\bfseries\raggedright\hyphenpenalty=10000\color{IlliniBlue}}{\thesection}{0.7em}{}
\titleformat{\subsection}{\large\bfseries\raggedright\hyphenpenalty=10000\color{IlliniBlue}}{\thesubsection}{0.7em}{}
\setlength{\parindent}{1.3em}
\setlength{\parskip}{3pt}
\linespread{1.06}
\setlist{itemsep=3pt,topsep=5pt}
\setcounter{tocdepth}{1}
\setcounter{secnumdepth}{2}
\pagestyle{fancy}
\fancyhf{}
\fancyhead[L]{\small\sffamily BADM 534}
\fancyhead[R]{\small\sffamily Lecture Notes}
\fancyfoot[C]{\thepage}
\renewcommand{\headrulewidth}{0.4pt}
\newcommand{\tightlist}{\setlength{\itemsep}{0pt}\setlength{\parskip}{0pt}}
\DefineVerbatimEnvironment{Highlighting}{Verbatim}{breaklines,fontsize=\small,commandchars=\\\{\}}
\newtheorem{assumption}{Assumption}[chapter]
\captionsetup{font=small,labelfont=bf,labelsep=period}
\setlength{\emergencystretch}{4em}
\providecommand{\pandocbounded}[1]{#1}

\clubpenalty=10000
\widowpenalty=10000

\begin{document}
\hypersetup{pageanchor=false}
\begin{titlepage}
\color{IlliniBlue}
\vspace*{1.2cm}
{\Huge\bfseries Marketing Analytics\par}
\vspace{.8cm}
{\LARGE Lecture Notes\par}
\vspace{.5cm}
{\large Marketing Decisions and AI-Assisted Coding\par}
\vspace{1.1cm}
\noindent{\color{IlliniOrange}\rule{\textwidth}{2pt}}\par
\vspace{.7cm}
{\large By Nathan Yang\par}
{\large October 23, 2026\par}
\end{titlepage}
\pagenumbering{roman}
\hypersetup{pageanchor=true}
\tableofcontents
\clearpage
\pagenumbering{arabic}
\hypertarget{M01-m1.-marketing-decisions-data-and-ai-assisted-python}{%
\chapter{Marketing Decisions, Data, and AI-Assisted
Python}\label{M01-m1.-marketing-decisions-data-and-ai-assisted-python}}

\noindent A retailer must decide which customers should receive a coupon, how much
inventory to place in each store, and which product descriptions should
appear online. Each decision needs a different analytical object. A
forecast of next week's purchases can inform inventory. A comparison of
purchases with and without a coupon can inform targeting. A summary of
product complaints can identify a service problem. An analysis becomes
useful when its output corresponds to the decision the manager can
actually change.

\hypertarget{M01-from-a-marketing-question-to-an-analytical-object}{%
\section{From a Marketing Question to an Analytical
Object}\label{M01-from-a-marketing-question-to-an-analytical-object}}

\noindent Consider a subscription service deciding whether to offer a renewal
discount. The service observes account tenure, past payments, recent
use, and prior contacts. A descriptive analysis reports renewal rates by
tenure. A predictive analysis estimates the probability that an account
renews under the conditions represented in the training data. A causal
analysis estimates how an offer changes renewal. A decision analysis
combines that change with margins, contact costs, and operational
constraints. These objects can use overlapping data while answering
different questions.

The timing of the decision determines what information is available. If
an account must be selected on Monday morning, a support call made on
Wednesday cannot be a predictor for that selection. A variable can be
strongly associated with an outcome and still be unusable at the
decision time. Write down a forecast origin, an outcome window, and a
feature window before constructing a table. An account-week table and an
account-lifetime table support different forecasts because their rows
and time horizons differ.

For a renewal offer, a useful decision quantity is incremental expected
contribution:

\[
\Delta\pi_i=m_i\{p_i(1)-p_i(0)\}-c_i.
\]

Here \(i\) identifies an account, \(p_i(1)\) and \(p_i(0)\) are its
renewal probabilities with and without the offer, \(m_i\) is
contribution in dollars per retained account over the specified horizon,
and \(c_i\) is the expected incremental cost of making the offer. This
compact expression assumes that the offer's costs can be summarized by
\(c_i\) and that \(m_i\) already reflects relevant variable costs. When
every renewing recipient receives a discount, discount expense depends
on renewal under treatment and should be modeled explicitly.

The difference between the two renewal probabilities matters because
many accounts would renew without assistance. Suppose contribution is
\$100 and contact cost is \$3. Then an account whose renewal probability
rises from 0.90 to 0.91 has incremental expected contribution of \$-2.
An account whose probability rises from 0.40 to 0.50 has contribution of
\$7. Ranking by treated renewal probability chooses the first account;
ranking by incremental contribution chooses the second. Observing only
treated customers would not reveal the untreated probabilities.

\begin{figure}[htbp]
\centering
\begin{tikzpicture}
\begin{axis}[name=left,width=6.3cm,height=6.1cm,xmin=.5,xmax=2.5,ybar,bar width=20pt,ymin=0,ymax=1.05,xtick={1,2},xticklabels={Account A,Account B},ylabel={Renewal probability with offer},ytick={0,0.5,1},xmajorgrids=false,ymajorgrids=true]
\addplot[fill=IlliniBlue,draw=IlliniBlue] coordinates {(1,0.91)(2,0.50)};
\node at (axis cs:1,.98) {0.91};\node at (axis cs:2,.57) {0.50};
\end{axis}
\begin{axis}[at={(left.east)},anchor=west,xshift=2.3cm,width=6.3cm,height=6.1cm,xmin=.5,xmax=2.5,ybar,bar width=20pt,ymin=-4,ymax=9,xtick={1,2},xticklabels={Account A,Account B},ylabel={Incremental contribution (dollars)},ytick={-2,0,7},ymajorgrids=true]
\addplot[fill=IlliniOrange,draw=IlliniOrange] coordinates {(1,-2)(2,7)};
\draw[IlliniBlue] (axis cs:.6,0)--(axis cs:2.4,0);
\end{axis}
\end{tikzpicture}

\caption{High response and high incremental value can identify different customers.}\label{fig:M01_response_vs_value}
\end{figure}

With contribution and contact cost held fixed, the first account's high
baseline renewal probability reverses its ranking under incremental
value. A response model and a targeting policy therefore need separate
evaluation.

\hypertarget{M01-the-unit-of-observation-determines-the-calculation}{%
\section{The Unit of Observation Determines the
Calculation}\label{M01-the-unit-of-observation-determines-the-calculation}}

\noindent A transaction system often stores one row per product line, while the
marketing question concerns orders or customers. Suppose order O1
contains two product lines, O2 contains one, and O3 contains three.
There are six rows, three orders, and perhaps only two customers.
Counting rows as purchases would report six orders. Averaging customer
attributes over product lines would give more weight to customers with
larger baskets.

A key uniquely identifies the intended unit. At the line level, the key
might be the combination of order ID and line number. At the order
level, it is order ID. At the customer level, it is customer ID. A
useful data contract states the key, allowed duplicates, measurement
units, missing-value codes, and date conventions. It also states whether
a negative quantity represents a return, whether a price includes tax,
and whether canceled orders belong in the analysis.

For line \(\ell\) in order \(o\), gross line sales are quantity times
unit price, \(s_{o\ell}=q_{o\ell}p_{o\ell}\). Order sales are
\(S_o=\sum_\ell s_{o\ell}\), and customer-period sales are the sum
across that customer's eligible orders during the period. These are
identities once the population and variables have been defined. They do
not require a behavioral assumption. The substantive choices enter
through eligibility, returns, discounts, and the observation window.

Joining data can change the population or its weights. A many-to-one
join from order lines to a customer table should add customer attributes
without multiplying the original rows. If the customer table
accidentally contains two records for a customer, every matching order
line appears twice. A program can finish without an error while
overstating sales. Verify the relationship between keys before joining,
and compare row counts and totals afterward.

\begin{figure}[htbp]
\centering
% Original synthetic teaching figure; see corresponding module.
\begin{tikzpicture}
\node[box,text width=3.1cm] (o) at (0,0) {One order\\Customer C1\\Revenue: \$40};
\node[box,text width=3.4cm] (one) at (5,1.6) {Unique lookup key\\C1: one customer row};
\node[orangebox,text width=3.4cm] (two) at (5,-1.6) {Duplicate lookup key\\C1: two customer rows};
\node[box,text width=2.8cm] (good) at (10,1.6) {One joined row\\Total: \$40};
\node[orangebox,text width=2.8cm] (bad) at (10,-1.6) {Two joined rows\\Total: \$80};
\draw[edge] (o)--(one); \draw[accent] (o)--(two);
\draw[edge] (one)--(good); \draw[accent] (two)--(bad);
\node[align=center,text width=11cm] at (5,-3.2) {A valid-looking join can double revenue.\\Check lookup-key uniqueness before merging, then reconcile totals.};
\end{tikzpicture}

\caption{A duplicate customer key multiplies transactions in a join.}\label{fig:M01_join_multiplication}
\end{figure}

For a single \$40 order, a customer lookup with two matching records
produces two joined rows and an apparent \$80 total. This error arises
from the structure of the tables. Increasing the sophistication of the
subsequent model cannot repair it.

\hypertarget{M01-python-objects-and-reproducible-calculations}{%
\section{Python Objects and Reproducible
Calculations}\label{M01-python-objects-and-reproducible-calculations}}

\noindent Python stores individual values in objects such as integers,
floating-point numbers, strings, and booleans. A list stores an ordered
collection; a dictionary associates names with values. A function
packages a calculation with explicit inputs and an output. These
distinctions matter when an assistant supplies code: a column of
numeric-looking strings does not behave like a numeric column, and an
empty string does not automatically represent a missing value.

Suppose two orders contain three product lines. Multiplying quantity by
unit price gives line sales; summing within each order gives order
sales. Returning each calculation as a number allows its result to be
checked and reused.

\begin{Verbatim}[breaklines,breakanywhere,fontsize=\small]
lines = [
    {"order": "O1", "quantity": 2, "unit_price": 10.0},
    {"order": "O1", "quantity": 1, "unit_price": 5.0},
    {"order": "O2", "quantity": 3, "unit_price": 4.0},
]

def line_sales(line):
    return line["quantity"] * line["unit_price"]

order_sales = {}
for line in lines:
    key = line["order"]
    order_sales[key] = order_sales.get(key, 0.0) + line_sales(line)

assert order_sales == {"O1": 25.0, "O2": 12.0}
assert sum(order_sales.values()) == 37.0
\end{Verbatim}

The total is \$37 across two orders, so average order value is \$18.50.
Dividing by the three product lines would instead give average line
sales, \$12.33. Neither denominator is inherently wrong; they describe
different objects. The mistake would be assigning one calculation the
other's interpretation.

A DataFrame provides a tabular interface for larger data. In pandas,
selecting columns, filtering rows, creating variables, grouping
observations, and joining tables are distinct operations. Use explicit
intermediate names when those operations change the unit of observation.
For example, \texttt{eligible\_lines}, \texttt{orders}, and
\texttt{customer\_months} communicate more than repeatedly overwriting
\texttt{df}. The
\href{https://pandas.pydata.org/docs/getting_started/intro_tutorials/}{pandas
tutorials} provide syntax references; the analytical specification
should determine which operations are needed.

Notebook execution order can hide dependencies. A cell may work because
an earlier interactive command created a variable that is absent from
the saved notebook. Before describing a notebook as reproducible,
restart the kernel and run all cells in order. Store file paths relative
to the project, record package versions, and distinguish source data
from derived outputs. A random seed helps reproduce a simulation or
split, but it does not guarantee identical results across every software
version or computing platform.

\hypertarget{M01-case-study-a-coupon-campaign-specification}{%
\section{Case Study: A Coupon Campaign
Specification}\label{M01-case-study-a-coupon-campaign-specification}}

\noindent A retailer has a monthly contact budget for 5,000 customers. The
proposed campaign offers a discount on the next purchase. The manager
initially asks for ``a model to find our best customers.'' That request
leaves the action, outcome, and objective unresolved. The
highest-spending customers may receive the offer without changing their
behavior; customers likely to respond may buy products with low margins.

A usable specification would identify the decision date, eligible
customer population, offer, evaluation horizon, and objective. One
possibility is to maximize incremental contribution during the next 30
days, subject to the contact limit. Eligibility might require marketing
consent and an active account. The analysis should define contribution
after discounts, returns, fulfillment, and contact costs. It should
identify which customer information existed before assignment and how
untreated outcomes will be measured.

The initial dataset might include customer-level histories and a
randomized pilot. Without such a pilot or another credible design, a
response model can help describe or predict behavior, but the
incremental contribution calculation remains unresolved. The next
operational step may therefore be a randomized test, rather than
immediate large-scale optimization.

A specification also needs acceptance criteria. The eligible-customer
table must have one row per customer, no post-assignment features, and
the intended date coverage. A campaign summary must reconcile assigned
counts, contact costs, purchases, and contribution. These checks are
small enough to state before asking an AI assistant to implement the
analysis. They protect the meaning of the result while allowing
flexibility in the code.

\hypertarget{M01-ai-assisted-coding-as-a-sequence-of-verifiable-steps}{%
\section{AI-Assisted Coding as a Sequence of Verifiable
Steps}\label{M01-ai-assisted-coding-as-a-sequence-of-verifiable-steps}}

\noindent An effective coding request names the inputs, output, definitions, and
checks. For the transaction example, ask for a function that accepts
product-line records, validates required fields, aggregates sales by
order, and returns both order totals and an exception report. Provide
the three-line example and its expected \$37 total. Ask the assistant to
explain what happens with a return, a missing order ID, or a duplicated
line.

Generate one analytical transformation at a time when the next step
depends on a definition. Inspect the proposed code before executing it,
especially commands that install software, delete files, upload data, or
invoke an external API. Read error messages as evidence about the code's
assumptions. A type error may reveal numeric values stored as strings; a
merge error may reveal an incorrect key. Repairing the error without
understanding its cause can preserve the underlying analytical mistake.

Verification should use both identities and meaningful counterexamples.
Order totals must sum to the line total for the same eligible rows. A
single known customer should produce a manually calculable summary. An
empty input should have an explicit result. A missing customer lookup
should not silently remove sales. These checks constrain what the
implementation is allowed to do without forcing one particular coding
style.

Keep a short record of consequential AI contributions: the task
requested, important changes to the plan or code, and the checks
performed. A saved conversation alone does not establish correctness.
The record should make it possible to explain why a transformation was
accepted.

\hypertarget{M01-interpreting-and-communicating-a-result}{%
\section{Interpreting and Communicating a
Result}\label{M01-interpreting-and-communicating-a-result}}

\noindent A useful recommendation names the decision, states the relevant estimate
with units and uncertainty, and explains how it changes the action.
``The model has high accuracy'' leaves the manager's decision
unresolved. ``At a \$3 contact cost and the stated contribution
assumptions, contact the customers with positive estimated incremental
contribution, subject to capacity'' connects the analysis to the
objective. The recommendation still needs evidence about how the
probabilities were estimated and whether the policy was evaluated on new
customers or periods.

Results should distinguish observation from assumption. A transaction
total is calculated from recorded eligible rows. A renewal probability
is an estimate. A contribution margin may be supplied by finance. An
untreated outcome for an offered customer is counterfactual. Listing
these roles prevents a spreadsheet or notebook from giving every input
the appearance of equal certainty.

If one customer places four orders and another places one, an
order-weighted mean gives the first customer four times the weight. If
both customers should count equally, calculate each customer's mean
first. If an AI-generated program reports six purchases from six product
lines, inspect whether purchase was defined as a line or an order. If an
offer has a high response rate, ask for the response rate under no offer
before interpreting the response as an effect.

\hypertarget{M01-readings}{%
\section{Readings}\label{M01-readings}}

\noindent The
\href{https://github.com/python-marketing-research/python-marketing-research-1ed}{Python
marketing research companion repository} provides examples tied to
marketing applications. Its notebooks are reference material and require
inspection of dependencies and data paths before adaptation. The
\href{https://pandas.pydata.org/docs/getting_started/intro_tutorials/}{pandas
tutorials} support tabular implementation.

\hypertarget{M01-mhealth-from-food-entries-to-a-decision-table}{%
\section{mHealth: From Food Entries to a Decision
Table}\label{M01-mhealth-from-food-entries-to-a-decision-table}}

\noindent An mHealth product manager must decide which food categories to suggest
when a customer returns to the app. Recent logs reveal repeated
categories, variation across days, and gaps in observation. The
immediate analytical task is to construct a reliable customer history. A
recommendation based on duplicated daily totals or an accidental change
in the customer population can look precise while describing the wrong
behavior.

The released files contain actual 2016 records. User identifiers have
been replaced, birthdays have been replaced with age on January 1, 2016,
and weight fields contain documented random perturbations. The food file
has 13,874,280 entries from 5,908 users. The separate daily file covers
a smaller set of users and dates. Food categories, descriptions, and
calorie fields retain their source values, apart from standardizing
provider references to mHealth. The
\href{../03_Data_and_Code/mHealth/DATA_DICTIONARY.md}{data dictionary}
defines the available fields and unresolved coding conventions. In
particular, \texttt{Weightchange}, \texttt{Goalplan}, and
\texttt{Gender} have undocumented code meanings. None establishes a
treatment, a behavioral goal, or a validated sex variable for an energy
equation.

The working cohort retains every entry for users whose first four
hexadecimal pseudonym digits have a value below 4,096. This fixed rule
reduces memory requirements and preserves within-user histories. Its
realized sample contains 917,153 entries, 370 users, and 76,804 observed
user-days. Ages at the reference date range from 17 to 76. The cohort is
a subset of app users who appear in the released file; it is not a
representative population sample. The adult nutrition sensitivity
analysis in M2 uses a separate age restriction.

A food-entry row records an item, while \texttt{Foodcalories} repeats
the source daily total. On the first eligible user-day selected by the
code below, five entries on July 4 have item calories summing to
1,132.7852. Each row repeats a daily total of 1,132.79. Summing that
repeated field gives 5,663.95, exactly five times the reported daily
total. The small difference between the item sum and one copy of the
daily total is compatible with rounding in this case. Larger
discrepancies elsewhere require investigation.

The daily key is \texttt{(Userid,\ date)}. A left join preserves all
76,804 cohort days, with 37,634 matched days and 39,170 unmatched days
in the separate daily file. Replacing it with an inner join would remove
more than half the cohort days. An unmatched row is therefore a
sample-definition decision, even if the software reports no error.

\begin{figure}[htbp]
\centering
\begin{tikzpicture}
\node[box,text width=4cm] (start) at (0,0) {Observed food-log days\\76,804};
\node[orangebox,text width=4.5cm] (left) at (-3,-2) {Left Join: 76,804 Days\\37,634 matched\\39,170 unmatched};
\node[box,text width=4.5cm] (inner) at (3,-2) {Inner Join: 37,634 Days\\37,634 matched\\39,170 removed};
\draw[accent] (start) -- (left);
\draw[edge] (start) -- (inner);
\node[align=center,text width=12cm] at (0,-4) {Same source keys; different analysis populations.\\An unmatched day is still an observed food-log day.};
\end{tikzpicture}

\caption{A left join preserves observed days while an inner join changes the analysis population.}\label{fig:M01_mhealth_join}
\end{figure}

The two branches retain the same matched records. Only the left-join
branch preserves the unmatched days needed to describe the food-log
cohort. Neither branch supplies unobserved calendar days.

\hypertarget{M01-loading-data-and-constructing-daily-measures}{%
\subsection{Loading Data and Constructing Daily
Measures}\label{M01-loading-data-and-constructing-daily-measures}}

\noindent Use Python 3.13 with NumPy 2.1.3, pandas 2.2.3, PyArrow 19.0.0, SciPy
1.15.3, scikit-learn 1.6.1, and statsmodels 0.14.4 to reproduce the
saved calculations. From the repository root, install the pinned
packages into an environment, verify the release files, and run the
analyses in dependency order:

\begin{Verbatim}[breaklines,breakanywhere,fontsize=\small]
python -m pip install -r 03_Data_and_Code/mHealth/recurring_examples/requirements.txt
python 03_Data_and_Code/mHealth/download_data.py
python 03_Data_and_Code/mHealth/recurring_examples/run_all.py
\end{Verbatim}

The downloader checks the published SHA-256 hashes. The first run reads
the compressed food file in chunks and saves a local Parquet cache;
subsequent runs reuse it. The cache is excluded from Git. Delete the
\texttt{recurring\_examples/cache} directory to force a fresh
derivation. \texttt{common.py} supplies the same sample and measurement
functions to every module. Paths are relative to the script location, so
the saved programs can also be run from another working directory.

For each observed day, category breadth counts distinct recorded
categories, assigning missing labels to \texttt{Unclassified}.
Calorie-share variety uses the sum of included item calories as its
denominator. A negative or missing item calorie value invalidates that
day's share calculation; a nonpositive item sum also makes shares
undefined. These conditions affect calorie-based variety, while the
day's existence can still be used to measure logging. The first reported
daily total is retained only after checking that its nonmissing values
agree within the day.

\FloatBarrier
\Needspace{8\baselineskip}
Python source:
\href{../03_Data_and_Code/mHealth/recurring_examples/common.py}{common.py}.
From its directory, run \texttt{python\ common.py}.

\begin{Verbatim}[breaklines,breakanywhere,fontsize=\small]
"""Shared input and measurement functions for the eight mHealth examples."""

from pathlib import Path
import hashlib
import json

import numpy as np
import pandas as pd

ROOT = Path(__file__).resolve().parent
DATA = ROOT.parent / "data"
CACHE = ROOT / "cache"
OUTPUT = ROOT / "results"
SEED = 534
COHORT_THRESHOLD = 4096


def save_result(module, result):
    OUTPUT.mkdir(exist_ok=True)
    (OUTPUT / f"{module}.json").write_text(
        json.dumps(result, indent=2, allow_nan=False) + "\n"
    )


def load_food():
    """Keep all rows for users whose first four pseudonym hex digits are <4096."""
    CACHE.mkdir(exist_ok=True)
    cached = CACHE / "food.parquet"
    metadata = CACHE / "source.json"
    manifest = json.loads((ROOT.parent / "manifest.json").read_text())
    signature = {
        "sha256": manifest["files"]["food_entries.csv.gz"]["sha256"],
        "threshold": COHORT_THRESHOLD,
        "preparation_version": 2,
    }
    if (
        cached.exists()
        and metadata.exists()
        and json.loads(metadata.read_text()) == signature
    ):
        return pd.read_parquet(cached)
    columns = [
        "Userid",
        "Fancydate",
        "Name",
        "Category",
        "Calories",
        "Foodcalories",
        "Budgetcalories",
        "Age_at_2016_01_01",
    ]
    parts = []
    for chunk in pd.read_csv(
        DATA / "food_entries.csv.gz", usecols=columns, chunksize=250000
    ):
        code = chunk.Userid.str.slice(1, 5).map(lambda value: int(value, 16))
        parts.append(chunk.loc[code < COHORT_THRESHOLD].copy())
    food = pd.concat(parts, ignore_index=True)
    food["date"] = pd.to_datetime(food.Fancydate)
    food["category"] = food.Category.fillna("Unclassified")
    food["description"] = (
        food.Name.fillna("")
        .str.casefold()
        .str.replace(r"\s+", " ", regex=True)
        .str.strip()
    )
    food.to_parquet(cached, index=False)
    metadata.write_text(json.dumps(signature, indent=2) + "\n")
    return food


def make_daily(food):
    """One row per observed user-day; missing calendar days are not created."""
    keys = ["Userid", "date"]
    grouped = food.groupby(keys, observed=True)
    daily = grouped.agg(
        entries=("Calories", "size"),
        category_count=("category", "nunique"),
        item_calories=("Calories", "sum"),
        food_calories=("Foodcalories", "first"),
        budget_calories=("Budgetcalories", "first"),
        daily_total_versions=("Foodcalories", "nunique"),
        budget_versions=("Budgetcalories", "nunique"),
    )
    daily.loc[daily.daily_total_versions != 1, "food_calories"] = np.nan
    daily.loc[daily.budget_versions != 1, "budget_calories"] = np.nan
    invalid = food.Calories.isna() | food.Calories.lt(0)
    daily["invalid_item_calories"] = invalid.groupby([food.Userid, food.date]).sum()
    categories = food.groupby(keys + ["category"], observed=True).Calories.sum()
    category_totals = categories.groupby(level=keys).transform("sum")
    concentration = (
        (categories / category_totals.where(category_totals > 0))
        .pow(2)
        .groupby(level=keys)
        .sum(min_count=1)
    )
    daily["variety"] = (1 - concentration).where(
        (daily.item_calories > 0) & (daily.invalid_item_calories == 0)
    )
    daily = daily.reset_index().sort_values(keys).reset_index(drop=True)
    assert not daily.duplicated(keys).any()
    assert daily.variety.dropna().between(-1e-10, 1 + 1e-10).all()
    return daily


def make_baskets(food):
    """Distinct observed categories; the missing-category label is not a candidate."""
    selected = food.loc[
        food.category != "Unclassified", ["Userid", "date", "category"]
    ].drop_duplicates()
    return (
        selected.groupby(["Userid", "date"], observed=True)
        .category.agg(lambda values: tuple(sorted(values)))
        .reset_index(name="basket")
    )


def description_holdout(text):
    """A fixed description-level split independent of labels or outcomes."""
    return int(hashlib.sha256(text.encode()).hexdigest()[:8], 16) % 5 == 0


def load_inputs():
    food = load_food()
    return food, make_daily(food), make_baskets(food)


if __name__ == "__main__":
    food, daily, baskets = load_inputs()
    daily.to_parquet(CACHE / "daily.parquet", index=False)
    baskets.to_parquet(CACHE / "baskets.parquet", index=False)
    result = {
        "rows": len(food),
        "users": int(food.Userid.nunique()),
        "user_days": len(daily),
        "categories": int(food.category.nunique()),
        "missing_category_rows": int(food.Category.isna().sum()),
        "invalid_item_calorie_rows": int(
            (food.Calories.isna() | food.Calories.lt(0)).sum()
        ),
        "days_with_inconsistent_daily_total": int(
            daily.daily_total_versions.ne(1).sum()
        ),
        "days_with_inconsistent_budget": int(daily.budget_versions.ne(1).sum()),
        "nonpositive_recorded_daily_totals": int(daily.food_calories.le(0).sum()),
        "variety_eligible_days": int(daily.variety.notna().sum()),
        "item_total_matches_recorded_total_within_one": int(
            (daily.item_calories - daily.food_calories).abs().le(1).sum()
        ),
    }
    save_result("cohort", result)
    print(json.dumps(result, indent=2))
\end{Verbatim}

\hypertarget{M01-join-validation-and-a-hand-checkable-day}{%
\subsection{Join Validation and a Hand-Checkable
Day}\label{M01-join-validation-and-a-hand-checkable-day}}

\noindent The join declares a one-to-one relationship and records unmatched keys.
The small day-level example checks both candidate aggregations against
the repeated total. Keeping these checks beside the join prevents a
later model from silently inheriting a different denominator.

\FloatBarrier
\Needspace{8\baselineskip}
Python source:
\href{../03_Data_and_Code/mHealth/recurring_examples/m01.py}{m01.py}.
From its directory, run \texttt{python\ m01.py}.

\begin{Verbatim}[breaklines,breakanywhere,fontsize=\small]
"""Audit the cohort, repeated totals, and the separate daily file."""

import json
import pandas as pd
from common import DATA, load_inputs, save_result

food, daily, baskets = load_inputs()
other = pd.read_csv(DATA / "daily_records.csv.gz", usecols=["Userid", "Fancydate"])
other["date"] = pd.to_datetime(other.Fancydate)
assert not other.duplicated(["Userid", "date"]).any()
joined = daily.merge(
    other[["Userid", "date"]],
    how="left",
    on=["Userid", "date"],
    validate="one_to_one",
    indicator=True,
)
assert len(joined) == len(daily)
example = daily.loc[(daily.entries.between(3, 6)) & daily.food_calories.gt(0)].iloc[0]
rows = food.loc[food.Userid.eq(example.Userid) & food.date.eq(example.date)]
result = {
    "rows": len(food),
    "users": int(food.Userid.nunique()),
    "days": len(daily),
    "matched_days": int(joined._merge.eq("both").sum()),
    "unmatched_days": int(joined._merge.eq("left_only").sum()),
    "example_date": str(example.date.date()),
    "example_entries": len(rows),
    "example_item_sum": float(rows.Calories.sum()),
    "example_daily_total": float(example.food_calories),
    "example_wrong_sum": float(rows.Foodcalories.sum()),
}
save_result("M01", result)
print(json.dumps(result, indent=2))
\end{Verbatim}

\hypertarget{M01-data-cleaning-as-a-sequence-of-decisions}{%
\subsection{Data Cleaning as a Sequence of
Decisions}\label{M01-data-cleaning-as-a-sequence-of-decisions}}

\noindent Structural checks stop calculations that violate their stated
observation unit. A duplicate key in a daily lookup can multiply food
records, so it must be resolved before joining. Review flags serve a
different purpose: they identify observations whose meaning needs
investigation. The cohort contains 52,838 rows equal on the selected
fields of user, day, normalized description, category, and calories.
Those fields omit event timestamps and meal identifiers. Repeated
servings can therefore produce equal rows. Automatically dropping them
would impose an unsupported definition of duplication.

The audit finds 206 food rows with missing categories, one negative item
calorie value, and 23 days with a nonpositive sum of item calories. The
calorie-share calculation consequently excludes 24 days. By contrast,
319 days have a nonpositive reported daily total. These counts differ
because the item sum and reported total are separate measurements. Their
discrepancy exceeds one calorie unit on 23,729 days. These flags can
overlap and must not be added as if they represented disjoint
exclusions.

A value outside a chosen review range is evidence for inspection, not
proof of an error. The reported daily calorie median is 1,475.63, and
the range is 0 to 9,190. A review flag below 100 or above 5,000 marks
651 days. Those thresholds are imposed sensitivity choices. They do not
establish nutritional adequacy or identify incorrect reports. Preserve
the released values, add flags, document the rule for each analysis, and
report the resulting sample. M2 compares nutrition-informed
alternatives; M3 repeats a regression after a stated restriction.

A useful AI coding request is: ``Create separate structural checks and
review flags. Stop on a duplicate daily key, preserve repeated food
entries, count overlapping flags separately, and return the retained
sample for each calculation.'' Verify the proposed implementation with a
two-entry day whose repeated total is 100 and whose calorie shares are
0.9 and 0.1. It must produce a daily total of 100 and variety of 0.18. A
gap from January 1 to January 3 must remain a missing January 2
observation. The executable counterexamples are in
\texttt{test\_logic.py}.

\FloatBarrier
\Needspace{8\baselineskip}
Python source:
\href{../03_Data_and_Code/mHealth/recurring_examples/audit.py}{audit.py}.
From its directory, run \texttt{python\ audit.py}.

\begin{Verbatim}[breaklines,breakanywhere,fontsize=\small]
"""Separate structural failures, review flags, and analysis-specific exclusions."""

import json
import numpy as np
import pandas as pd
from common import DATA, load_inputs, save_result


def require_unique(frame, keys):
    if frame[keys].isna().any().any():
        raise ValueError("Missing join key")
    if frame.duplicated(keys).any():
        raise ValueError("Duplicate key at the declared observation unit")


def daily_flags(daily):
    flags = pd.DataFrame(index=daily.index)
    flags["invalid_item"] = daily.invalid_item_calories.gt(0)
    flags["nonpositive_item_sum"] = daily.item_calories.le(0)
    flags["nonpositive_daily_total"] = daily.food_calories.le(0)
    flags["missing_daily_total"] = daily.food_calories.isna()
    flags["total_discrepancy_over_one"] = (
        (daily.food_calories - daily.item_calories).abs().gt(1)
    )
    flags["daily_total_below_100_or_above_5000"] = daily.food_calories.lt(
        100
    ) | daily.food_calories.gt(5000)
    return flags


def run_audit():
    food, daily, baskets = load_inputs()
    assert food.Userid.notna().all() and food.date.notna().all()
    assert food.date.between("2016-01-01", "2016-12-31").all()
    require_unique(daily, ["Userid", "date"])
    # Repeated entries can be legitimate. Flag equality on these fields without deletion.
    fields = ["Userid", "date", "description", "category", "Calories"]
    duplicate_flags = food.duplicated(fields, keep=False)
    flags = daily_flags(daily)
    health = pd.read_csv(
        DATA / "daily_records.csv.gz",
        usecols=["Userid", "Fancydate", "Age_at_2016_01_01", "Currentweight"],
    )
    health = health.loc[health.Userid.isin(food.Userid.unique())]
    require_unique(health, ["Userid", "Fancydate"])
    age_versions = health.groupby("Userid").Age_at_2016_01_01.nunique(dropna=False)
    assert age_versions.eq(1).all()
    result = {
        "food_rows": len(food),
        "daily_rows": len(daily),
        "missing_category_rows": int(food.Category.isna().sum()),
        "rows_equal_on_selected_entry_fields": int(duplicate_flags.sum()),
        "missing_item_calorie_rows": int(food.Calories.isna().sum()),
        "negative_item_calorie_rows": int(food.Calories.lt(0).sum()),
        "flags_by_user_day": {key: int(value) for key, value in flags.sum().items()},
        "recorded_calorie_quantiles": {
            str(q): float(value)
            for q, value in daily.food_calories.quantile(
                [0, 0.01, 0.5, 0.99, 1]
            ).items()
        },
        "matched_daily_file_rows": len(health),
        "fixed_age_min": int(health.Age_at_2016_01_01.min()),
        "fixed_age_max": int(health.Age_at_2016_01_01.max()),
        "nonpositive_weight_rows": int(health.Currentweight.le(0).sum()),
    }
    save_result("audit", result)
    print(json.dumps(result, indent=2))


if __name__ == "__main__":
    run_audit()
\end{Verbatim}

\hypertarget{M01-python-calculations-for-the-synthetic-cases}{%
\section{Python Calculations for the Synthetic
Cases}\label{M01-python-calculations-for-the-synthetic-cases}}

\noindent The inputs below are imposed teaching values. Running the program
creates the labeled synthetic CSV files and prints the corresponding
calculations. They are separate from the mHealth estimates. Keep the
script beside \texttt{example\_helpers.py} in
\texttt{03\_Data\_and\_Code/Lecture\_Examples}.

\FloatBarrier
\Needspace{8\baselineskip}
Python source:
\href{../03_Data_and_Code/Lecture_Examples/example_helpers.py}{example\_helpers.py}.
From its directory, run \texttt{python\ example\_helpers.py}.

\begin{Verbatim}[breaklines,breakanywhere,fontsize=\small]
"""File output and arithmetic helpers for synthetic lecture examples."""

import csv
import itertools
import json
import math
from pathlib import Path

ROOT = Path(__file__).resolve().parent


def write_csv(name, fields, rows):
    with (ROOT / name).open("w", newline="", encoding="utf-8") as f:
        out = csv.writer(f)
        out.writerow(fields)
        out.writerows(rows)


def ols_slope(points):
    xbar = sum(x for x, y in points) / len(points)
    ybar = sum(y for x, y in points) / len(points)
    slope = sum((x - xbar) * (y - ybar) for x, y in points) / sum(
        (x - xbar) ** 2 for x, y in points
    )
    return {"intercept": ybar - slope * xbar, "slope": slope}


def shares(utilities):
    weights = {key: math.exp(val) for key, val in utilities.items()}
    return {key: weight / sum(weights.values()) for key, weight in weights.items()}
\end{Verbatim}

\FloatBarrier
\Needspace{8\baselineskip}
Python source:
\href{../03_Data_and_Code/Lecture_Examples/m01_synthetic.py}{m01\_synthetic.py}.
From its directory, run \texttt{python\ m01\_synthetic.py}.

\begin{Verbatim}[breaklines,breakanywhere,fontsize=\small]
"""Synthetic numerical examples for M1; run from any working directory."""

import itertools
import json
import math
from example_helpers import write_csv, ols_slope, shares

r = {}
lines = [("O1", "C1", 2, 10), ("O1", "C1", 1, 5), ("O2", "C2", 3, 4)]
write_csv(
    "M01_order_lines_synthetic.csv",
    ["order_id", "customer_id", "quantity", "unit_price_dollars"],
    lines,
)
orders = {}
for order, customer, qty, price in lines:
    orders[order] = orders.get(order, 0) + qty * price
r["M01"] = {
    "revenue": sum(orders.values()),
    "order_count": len(orders),
    "average_order_value": sum(orders.values()) / len(orders),
    "average_line_revenue": sum(orders.values()) / len(lines),
    "renewal_incremental_value_A": 100 * (0.91 - 0.90) - 3,
    "renewal_incremental_value_B": 100 * (0.50 - 0.40) - 3,
}
assert orders == {"O1": 25, "O2": 12}
print(json.dumps(r, indent=2))
\end{Verbatim}

\hypertarget{M01-mhealth-an-auditable-monthly-business-review}{%
\section{mHealth: An Auditable Monthly Business
Review}\label{M01-mhealth-an-auditable-monthly-business-review}}

\noindent An mHealth manager needs to decide whether a decline in food entries
warrants attention to customer engagement. Entry counts combine the
number of eligible customers, the length of the reporting month, the
frequency of logging, and the intensity of logging on observed days. A
decline can arise from any of these components. The comparison becomes
interpretable when the eligible population is fixed before the reporting
periods and every component has an explicit denominator.

Define the population as users in the teaching cohort with at least one
observed day before May 1, 2016. All members remain in the May and June
denominators, including those with no logs in a month. A missing day
contributes zero to the indicator for an observed log; its food
consumption remains unobserved. For month \(t\), the number of entries
satisfies the accounting identity

\[
L_t=N D_t p_t a_t,
\]

where \(N\) is the fixed number of eligible users, \(D_t\) is calendar
days in the month, \(p_t\) is observed user-days divided by \(ND_t\),
and \(a_t\) is entries per observed user-day. The identity separates
calendar exposure from logging behavior. It uses counts of records, with
no assumption that an entry represents a purchase or a complete meal.

\begin{table}[H]
\centering\small
\begin{tabular}{@{}lrr@{}}
\toprule
Quantity & May & June \\
\midrule

Eligible Users & 340 & 340 \\
Calendar User-Days & 10,540 & 10,200 \\
Observed User-Days & 7,082 & 6,331 \\
Observed Logging Fraction & 0.6719 & 0.6207 \\
Entries per Observed Day & 11.7995 & 11.6910 \\
Food Entries & 83,564 & 74,016 \\
\bottomrule
\end{tabular}
\end{table}

The table identifies a 5.12 percentage-point decline in the logging
fraction, alongside a shorter month and a modest decline in entries per
observed day. Changing the denominator to users who happened to log
during June would remove inactive members from the comparison. The
fixed-population result supports investigation of engagement, recording
practices, and seasonality. It supplies no estimate of a campaign effect
or change in nutrition.

An AI-assisted business review should retrieve numerical claims from a
computed evidence file. The file below records the source checksum,
membership rule, observation unit, month, numerator, and denominator. A
useful coding request is: ``Construct the fixed cohort before May, count
distinct user-dates, retain members with zero logs, and verify the
accounting identity separately for each month.'' The independent
calculation must finish before an assistant writes the managerial
explanation.

The \href{https://github.com/takechanman1228/claude-ecom}{claude-ecom
reference} separates numerical computation from a narrative layer for
ecommerce reviews. The mHealth calculation applies that separation to
logging data. Its descriptive evidence file also supplies the
structured-claim checks in M8.

\FloatBarrier
\Needspace{8\baselineskip}
Python source:
\href{../03_Data_and_Code/mHealth/recurring_examples/reporting.py}{reporting.py}.
From its directory, run \texttt{python\ reporting.py}.

\begin{Verbatim}[breaklines,breakanywhere,fontsize=\small]
"""Compute an auditable, fixed-cohort mHealth business review locally."""

import json
import pandas as pd
from common import ROOT, load_food, save_result


def monthly_review(food):
    # Fix membership before either reporting month; retain zero-log users.
    members = sorted(food.loc[food.date < "2016-05-01", "Userid"].unique())
    rows = []
    for month in ["2016-05", "2016-06"]:
        period = pd.Period(month, freq="M")
        part = food.loc[
            food.Userid.isin(members) & food.date.dt.to_period("M").eq(period)
        ]
        days = part[["Userid", "date"]].drop_duplicates()
        denominator = len(members) * period.days_in_month
        assert len(days) <= denominator
        rows.append(
            {
                "period": month,
                "eligible_users": len(members),
                "calendar_user_days": denominator,
                "observed_user_days": len(days),
                "users_with_logs": int(days.Userid.nunique()),
                "food_entries": len(part),
                "logging_fraction": len(days) / denominator,
                "entries_per_observed_day": len(part) / len(days),
            }
        )
    manifest = json.loads((ROOT.parent / "manifest.json").read_text())
    return {
        "source_sha256": manifest["files"]["food_entries.csv.gz"]["sha256"],
        "membership": "At least one observed day before 2016-05-01",
        "unit": "eligible user-calendar-day",
        "metric": "logging_fraction",
        "interpretation": "descriptive",
        "periods": rows,
    }


if __name__ == "__main__":
    result = monthly_review(load_food())
    save_result("reporting", result)
    print(json.dumps(result, indent=2))
\end{Verbatim}

\hypertarget{M02-m2.-customer-data-measurement-and-visualization}{%
\chapter{Customer Data, Measurement, and
Visualization}\label{M02-m2.-customer-data-measurement-and-visualization}}

\noindent A grocery retailer sees larger online baskets and asks whether moving
customers to its app will raise sales. The raw comparison combines the
shopping channel with who uses it, why they shop, and which trips are
observed. A household may order online for a planned weekly restock and
visit a store for a forgotten ingredient. Measurement and visualization
determine whether the analysis reveals that distinction or conceals it.

\hypertarget{M02-define-the-population-before-the-metric}{%
\section{Define the Population Before the
Metric}\label{M02-define-the-population-before-the-metric}}

\noindent A customer dataset is a record of an observation process. Loyalty
accounts exclude unidentified shoppers. Online logs record sessions that
may contain repeated visits by the same person. Orders exclude browsing
without purchase. App events reflect tracking rules, consent, device
settings, and instrumentation changes. Before calculating a metric,
identify the target population and the mechanism by which its members
enter the table.

Suppose monthly revenue per customer is intended to describe all
eligible accounts. Dividing total revenue by purchasing accounts
produces revenue per purchasing account. If inactivity changes, that
measure can increase while revenue per eligible account falls. The
denominator should therefore be written into the metric name or its
definition. For customer acquisition cost, specify whether the
denominator is new registered accounts, first-time purchasers, or
customers with positive contribution over a defined horizon.

A conversion rate requires an eligible exposure count. If \(b\) eligible
sessions produce at least one purchase among \(n\) eligible sessions,
session conversion is \(b/n\). Customer conversion instead uses
customers in both numerator and denominator. A customer with ten
sessions should not receive ten times the weight in a customer-level
rate. When the same customer contributes repeated observations,
statistical uncertainty must also account for that dependence.

Returns require a substantive rule. Gross sales, net sales after
returns, and contribution after variable costs answer different
questions. A comparison that deducts online returns but observes only
gross offline sales mixes measurement rules. Currency, tax, fulfillment
fees, and discounts should be aligned across channels before
interpreting differences.

\hypertarget{M02-build-a-customer-table-without-losing-customers}{%
\section{Build a Customer Table Without Losing
Customers}\label{M02-build-a-customer-table-without-losing-customers}}

\noindent Begin with a customer roster that defines eligibility. Aggregate
eligible transactions to customer-period totals, then left-join those
totals to the roster. A customer with no transactions can receive zero
sales if the transaction system fully captures activity during the
observation period. Zero should not replace a value that is missing
because the customer was unobserved or the data feed failed.

A join should have a documented relationship: one-to-one, many-to-one,
or one-to-many. Validate key uniqueness in the lookup table and inspect
unmatched records. A match rate of 99\% can still conceal a material
problem if the unmatched 1\% contains the highest-value customers or one
entire channel. Report both counts and relevant value totals for
unmatched observations.

Dates need similar attention. Convert timestamps to a consistent
timezone before assigning day or week. Determine whether the date
represents ordering, shipment, payment, or return. A promotion announced
late in a period may affect orders immediately but recognized revenue
later. Aligning metrics on an inappropriate date can create an apparent
lag that comes from accounting rather than customer behavior.

Outliers should be investigated before removal. A \$10,000 order could
be a business customer, a data-entry error, or a legitimate event. A
rule that trims all orders above a percentile changes the estimand to a
restricted population. State that restriction and compare conclusions
with and without the rule when it can affect the decision. Robust
summaries such as medians and quantiles often reveal more than one mean
with an undocumented exclusion.

\hypertarget{M02-case-study-a-reversal-caused-by-customer-mix}{%
\section{Case Study: A Reversal Caused by Customer
Mix}\label{M02-case-study-a-reversal-caused-by-customer-mix}}

\noindent A marketing manager compares campaigns A and B using conversion. Both
campaigns reach new and returning customers, but their audience
composition differs. Assume the following customer counts and outcomes,
with each customer belonging to one segment and appearing once.

\begin{table}[H]
\centering\small
\begin{tabular}{@{}
  >{\raggedright\arraybackslash}p{(\columnwidth - 10\tabcolsep) * \real{0.13}}
  >{\raggedleft\arraybackslash}p{(\columnwidth - 10\tabcolsep) * \real{0.17}}
  >{\raggedleft\arraybackslash}p{(\columnwidth - 10\tabcolsep) * \real{0.17}}
  >{\raggedleft\arraybackslash}p{(\columnwidth - 10\tabcolsep) * \real{0.17}}
  >{\raggedleft\arraybackslash}p{(\columnwidth - 10\tabcolsep) * \real{0.17}}
  >{\raggedleft\arraybackslash}p{(\columnwidth - 10\tabcolsep) * \real{0.19}}@{}}
\toprule
Campaign & New Customers & New Conversions & Returning Customers &
Returning Conversions & Overall Conversion \\
\midrule

A & 100 & 9 & 100 & 80 & 89/200 = 44.5\% \\
B & 200 & 20 & 100 & 81 & 101/300 = 33.7\% \\
\bottomrule
\end{tabular}
\end{table}

Campaign B has higher conversion among new customers, 10\% versus 9\%,
and among returning customers, 81\% versus 80\%. Campaign A nevertheless
has the higher aggregate rate because half its audience consists of
returning customers, compared with one-third for B. Returning customers
convert much more frequently in both campaigns.

The aggregate rate is a weighted average of segment rates:

\[
\bar p_j=\sum_g w_{jg}p_{jg},\qquad \sum_g w_{jg}=1.
\]

Here \(j\) indexes campaigns, \(g\) indexes customer segments,
\(p_{jg}\) is the conversion fraction within a segment, and \(w_{jg}\)
is that segment's share of the campaign audience. The weights are part
of the comparison. Holding them at 50\% new and 50\% returning gives
standardized rates of 44.5\% for A and 45.5\% for B. This adjustment
describes performance for the chosen common mix. It establishes a causal
campaign effect only if the design supports that interpretation.

\begin{figure}[htbp]
\centering
% Original synthetic teaching figure; see corresponding module.
\begin{tikzpicture}
\begin{axis}[ybar,bar width=15pt,width=13cm,height=7cm,ymin=0,ymax=100,xmin=.4,xmax=4.6,xtick={1,2,3,4},xticklabels={New,Returning,Observed mix,Common mix},ylabel={Conversion rate (percent)},ytick={0,20,40,60,80,100},ymajorgrids=true,legend style={at={(.5,1.13)},anchor=south,legend columns=2},nodes near coords,nodes near coords style={font=\scriptsize},every node near coord/.append style={/pgf/number format/fixed,/pgf/number format/precision=1}]
\addplot[fill=IlliniBlue,draw=IlliniBlue] coordinates {(1,9)(2,80)(3,44.5)(4,44.5)};
\addplot[fill=IlliniOrange,draw=IlliniOrange] coordinates {(1,10)(2,81)(3,33.6667)(4,45.5)};
\legend{Campaign A,Campaign B}
\end{axis}
\end{tikzpicture}

\caption{Audience composition reverses the aggregate campaign ranking.}\label{fig:M02_mix_reversal}
\end{figure}

B's small advantages within both segments coexist with A's aggregate
advantage. A manager deciding where to send the next campaign needs to
specify the future audience mix and whether customers were assigned
comparably. An aggregate bar chart alone cannot answer that question.

\hypertarget{M02-cohorts-retention-and-customer-age}{%
\section{Cohorts, Retention, and Customer
Age}\label{M02-cohorts-retention-and-customer-age}}

\noindent A cohort groups customers by a common starting event, such as first
purchase or subscription activation. Retention at customer age \(k\) is
the fraction of the original eligible cohort meeting a stated activity
criterion at that age. A cohort table separates calendar time from
customer age: an October signup observed one month later is not
comparable with an April signup observed seven months later unless the
question explicitly concerns those different ages.

For cohort \(c\), define \(N_c\) as the number of eligible customers at
entry and \(A_{c,k}\) as the number active at age \(k\). The descriptive
retention rate is \(R_{c,k}=A_{c,k}/N_c\). A later cohort may not yet
have reached age \(k\); that cell is unobserved, not zero. Omitting
inactive customers from the denominator would turn retention into a
different conditional measure.

\begin{table}[H]
\centering\small
\begin{tabular}{@{}
  >{\raggedright\arraybackslash}p{(\columnwidth - 8\tabcolsep) * \real{0.15}}
  >{\raggedleft\arraybackslash}p{(\columnwidth - 8\tabcolsep) * \real{0.21}}
  >{\raggedleft\arraybackslash}p{(\columnwidth - 8\tabcolsep) * \real{0.21}}
  >{\raggedleft\arraybackslash}p{(\columnwidth - 8\tabcolsep) * \real{0.21}}
  >{\raggedleft\arraybackslash}p{(\columnwidth - 8\tabcolsep) * \real{0.22}}@{}}
\toprule
Signup Cohort & Entry Customers & Active at Month 1 & Active at Month 2
& Active at Month 3 \\
\midrule

January & 100 & 60 & 45 & 40 \\
February & 100 & 70 & 50 & Not yet observed \\
March & 100 & 80 & Not yet observed & Not yet observed \\
\bottomrule
\end{tabular}
\end{table}

Month-one retention rises from 60\% for the January cohort to 80\% for
the March cohort. March customers have not yet reached the later
observation windows, so their long-run retention remains unobserved.
Comparing March's month one with January's month three would combine
cohort quality with customer age. The right comparison depends on
whether the decision concerns onboarding, long-term value, or a
calendar-time change affecting all existing customers.

\begin{figure}[htbp]
\centering
% Original synthetic teaching figure; see corresponding module.
\begin{tikzpicture}[x=2.25cm,y=1.05cm]
\foreach \c/\lab in {0/At Entry,1/Month 1,2/Month 2,3/Month 3} {\node[font=\small\bfseries] at (\c,1) {\lab};}
\foreach \r/\lab in {0/January,-1/February,-2/March} {\node[anchor=east] at (-.6,\r) {\lab};}
\foreach \x/\y/\v in {0/0/100,1/0/60,2/0/45,3/0/40,0/-1/100,1/-1/70,2/-1/50,0/-2/100,1/-2/80} {\draw[IlliniBlue] (\x-.47,\y-.4) rectangle (\x+.47,\y+.4);\node at (\x,\y) {\v\%};}
\foreach \x/\y in {3/-1,2/-2,3/-2} {\fill[LightStorm] (\x-.47,\y-.4) rectangle (\x+.47,\y+.4);\node[align=center,font=\scriptsize] at (\x,\y) {Not yet\\observed};}
\node[align=center,text width=11cm] at (1.3,-3.2) {Each cohort begins with 100 customers.\\Only observed customer ages enter the retention denominator.};
\end{tikzpicture}

\caption{A cohort’s unobserved future is different from zero retention.}\label{fig:M02_cohort_triangle}
\end{figure}

Unobserved future periods require forecasts from an explicit model.
Coding those periods as zero would treat unavailable outcomes as
observed failures and bias retention downward.

\hypertarget{M02-visual-encodings-should-match-the-question}{%
\section{Visual Encodings Should Match the
Question}\label{M02-visual-encodings-should-match-the-question}}

\noindent Use position on a common scale when precise comparison matters. A dot
plot can compare segment conversion rates more clearly than area-sized
bubbles. A line chart is appropriate when the ordering of time is
meaningful. A histogram or empirical cumulative distribution reveals
skewness that a mean hides. A scatterplot supports examination of a
relationship, with transparency or aggregation when observations
overlap.

The scale changes the interpretation. A logarithmic sales axis
emphasizes proportional differences; a linear axis emphasizes dollar
differences. A truncated bar-chart axis exaggerates differences because
bar length encodes magnitude from its baseline. A carefully labeled
narrow axis can be useful for a line or dot plot when the scientific
question concerns small differences. The choice should preserve the
intended comparison and disclose the scale.

Uncertainty should be shown for an estimated comparison when it affects
a decision. For an independent Bernoulli sample, a common large-sample
standard error for a proportion is \(\sqrt{\hat p(1-\hat p)/n}\). Near
zero or one, or with small samples, use an interval appropriate to
binomial data. With repeated observations per customer, treating every
session as independent generally understates uncertainty. Resample or
cluster at the unit that captures the relevant dependence.

Color can encode categories or highlight the comparison that matters.
Use line styles, markers, and labels alongside color, and keep small
text dark enough for contrast. A figure should remain interpretable when
printed in grayscale or described to a student who cannot see it.
Explain the axes, units, sample, and principal comparison in its caption
or accompanying text.

\hypertarget{M02-case-study-online-grocery-baskets-and-repeated-purchases}{%
\section{Case Study: Online Grocery Baskets and Repeated
Purchases}\label{M02-case-study-online-grocery-baskets-and-repeated-purchases}}

\noindent Chintala, Liaukonytė, and Yang's
\href{https://doi.org/10.1287/mksc.2022.0292}{grocery-shopping study}
compares online and offline baskets using household shopping records. It
studies basket composition and repeated purchasing, with attention to
differences in shopping-trip types. Larger baskets need not contain a
wider range of categories, and a household's choice of channel can be
related to the purpose of its trip.

Suppose one household buys the sets \{milk, bread, apples, yogurt\} and
\{milk, bread, apples, coffee\} on successive trips. Each basket
contains four distinct items. Their Jaccard similarity is

\[
J(A,B)=\frac{|A\cap B|}{|A\cup B|}=\frac{3}{5}=0.60.
\]

The numerator counts shared distinct items and the denominator counts
all distinct items appearing in either trip. A different pair of
four-item baskets with one shared item has similarity
\(1/7\approx0.143\). Basket size is identical in the two examples, while
repeated purchasing differs substantially. Whether quantities or
spending weights should enter the measure depends on the decision;
set-based Jaccard treats one unit and ten units of an item alike.

A defensible descriptive study would define eligible grocery trips,
compare channels within households when appropriate, and examine basket
size, distinct categories, repeat-item overlap, and spending. It would
inspect whether findings are driven by large restocking trips, small
convenience trips, or particular retailers. Channel labels should not
silently determine the definition of a ``typical'' trip when the
objective is a comparable classification. Clustering can classify trips
using their observed composition.

A retail recommendation might distinguish discovery from replenishment.
High repetition can support convenient reordering while making it harder
for unfamiliar products to receive attention. A proposed interface
change should therefore specify both the desired outcome and the
possible tradeoff. Its effect requires evaluation under an appropriate
design, even when the descriptive pattern provides a strong reason to
test it.

\hypertarget{M02-variety-depends-on-both-breadth-and-concentration}{%
\subsection{Variety Depends on Both Breadth and
Concentration}\label{M02-variety-depends-on-both-breadth-and-concentration}}

\noindent Two customers can purchase the same number of categories while
allocating spending very differently. One divides a basket equally
between two categories; another assigns 90\% to one and 10\% to the
other. A category count assigns both customers breadth two. To measure
how evenly spending is distributed, define \(V=1-\sum_c s_c^2\), where
\(s_c\) is category \(c\)'s share of the customer's spending within a
stated basket or period. The shares must use a common denominator and
sum to one. The two baskets have \(V=0.50\) and \(V=0.18\),
respectively.

\href{https://doi.org/10.1086/720453}{Nielsen and coauthors (2022)} use
item counts and a reverse-coded concentration measure based on calorie
shares to study food-consumption variety. The underlying unit matters:
spending shares, calorie shares, and purchase counts can rank the same
baskets differently. A wider category definition also changes measured
variety without changing consumption. Report the category mapping,
observation window, and share denominator before interpreting a trend or
segment difference. M7, Section 7.12, connects these measurements to
recommendation design.

\hypertarget{M02-from-evidence-to-a-management-presentation}{%
\section{From Evidence to a Management
Presentation}\label{M02-from-evidence-to-a-management-presentation}}

\noindent The Pyramid Principle organizes a recommendation around a
decision-relevant answer, followed by a small number of supporting
reasons and the evidence for each reason. This structure works when the
answer's strength matches the analysis. A descriptive study might
recommend a test or a closer diagnostic comparison, rather than
immediate deployment.

For the campaign example, a clear opening is: ``Compare the campaigns at
a common customer mix before reallocating the budget.'' The supporting
reasons are that returning customers convert more often, the campaigns
reached different mixes, and the within-segment ranking differs from the
aggregate ranking. The supporting evidence should quantify each reason
and connect it to the allocation decision.

An effective chart title can state the comparison, such as ``Campaign
A's aggregate advantage reflects its returning-customer share.'' The
caption should define the conversion denominator and identify the data
as synthetic.

\hypertarget{M02-ai-assisted-workflow-and-practice}{%
\section{AI-Assisted Workflow and
Practice}\label{M02-ai-assisted-workflow-and-practice}}

\noindent Ask an assistant to construct a customer-period table from a specified
roster and transaction table, preserving customers with no observed
purchases. Require a join audit, a missing-data report, and a
reconciliation of total sales. Next ask it to calculate the campaign
table above and standardized rates under a 50/50 audience mix. Verify
those small numbers before applying the same functions to a larger file.

For a chart request, specify the question, observation unit, axis units,
ordering, and uncertainty representation. Ask the assistant to describe
what visual pattern would change the recommendation. Compare its
description with the plotted data. If it labels an observational channel
difference as an effect, revise the interpretation and identify the
design that would be needed to support the stronger claim.

For practice, change campaign B's audience to 100 new and 100 returning
customers while keeping its segment conversion rates fixed. Its
aggregate rate becomes 45.5\%, and the reversal disappears. Next
calculate month-two retention for the January and February cohorts: 45\%
and 50\%. March has no observed month-two rate. Finally, compare a
four-item basket with itself: Jaccard similarity is one, regardless of
whether the quantities have changed. Each exercise tests a definition
that should remain explicit in an automated analysis.

\hypertarget{M02-readings}{%
\section{Readings}\label{M02-readings}}

\noindent \href{https://clauswilke.com/dataviz/}{Wilke's visualization text}
offers a visual reference for distributions, comparisons, and
uncertainty. The
\href{https://pandas.pydata.org/docs/user_guide/groupby.html}{pandas
grouping documentation} supports implementation.

\hypertarget{M02-mhealth-recorded-variety-and-missing-days}{%
\section{mHealth: Recorded Variety and Missing
Days}\label{M02-mhealth-recorded-variety-and-missing-days}}

\noindent The manager wants to know whether customers with broad recorded diets
need different suggestions from customers who repeatedly log a few
categories. Category count measures breadth. A calorie-share measure
also distinguishes an even distribution from a day dominated by one
category. For user \(i\), day \(t\), and recorded category \(c\), define
\(k_{ict}\) as the sum of item calories in that category. On days with
nonnegative observed item values and a positive item total, define

\[
s_{ict}=\frac{k_{ict}}{\sum_h k_{iht}},\qquad
V_{it}=1-\sum_c s_{ict}^2.
\]

The shares sum to one across recorded categories, including
\texttt{Unclassified} where necessary. \(V_{it}\) approaches zero when
one category accounts for almost all recorded item calories. It rises
when calorie shares become more evenly distributed across categories.
Multiplying every item calorie value by the same positive constant
leaves these shares unchanged. This scale invariance helps distinguish
composition from recorded quantity, although omitted foods can still
change composition.

Among the 76,780 eligible user-days, mean variety is 0.7834 when each
day has equal weight and 0.7855 when each user first contributes a
personal average. The small difference in these means does not make the
weighting choices interchangeable. A day-weighted comparison gives
frequent loggers more influence. A user-weighted comparison answers a
question about the average observed user.

\begin{table}[H]
\centering\small
\begin{tabular}{@{}lrr@{}}
\toprule
Recorded Categories & Eligible Days & Mean Variety \\
\midrule

1--5 & 11,861 & 0.5306 \\
6--10 & 30,258 & 0.7895 \\
11--20 & 32,492 & 0.8618 \\
21 or More & 2,169 & 0.9046 \\
\bottomrule
\end{tabular}
\end{table}

Breadth and calorie-share variety rise together in these groups, partly
because adding categories expands the possible distribution of shares.
This is a descriptive relation among recorded days. It does not show
that the app caused broader consumption or that greater category variety
improved health. Source categories are not a validated diet-quality
index.

The complete 2016 calendar contains 135,420 potential user-days for the
selected users. Only 76,804 have food entries; 58,616 are unobserved. A
zero-filled calorie panel would interpret every absent record as no
consumption. For a logging outcome, zero can instead mean that no entry
was observed on a specified calendar day. The variable's definition
determines whether filling a missing cell with zero is appropriate.
These users are identified from a full-year file, so the calendar is a
descriptive grid rather than a verified account-enrollment roster.

The item sum is within one calorie unit of the reported daily total on
53,075 observed days. Using the reported total as the denominator of
category shares would fail to make shares sum to one on many other days.
The calculation below preserves both measures, constructs the
observed-day summary, and builds the calendar separately.

\FloatBarrier
\Needspace{8\baselineskip}
Python source:
\href{../03_Data_and_Code/mHealth/recurring_examples/m02.py}{m02.py}.
From its directory, run \texttt{python\ m02.py}.

\begin{Verbatim}[breaklines,breakanywhere,fontsize=\small]
"""Measure recorded variety and separate observed days from calendar days."""

import json
import pandas as pd
from common import OUTPUT, load_inputs, save_result

food, daily, baskets = load_inputs()
eligible = daily.loc[daily.variety.notna()].copy()
eligible["breadth_group"] = pd.cut(
    eligible.category_count,
    [0, 5, 10, 20, 1000],
    labels=["1-5", "6-10", "11-20", "21+"],
)
table = eligible.groupby("breadth_group", observed=True).agg(
    days=("variety", "size"),
    mean_variety=("variety", "mean"),
    mean_recorded_calories=("food_calories", "mean"),
)
OUTPUT.mkdir(exist_ok=True)
table.to_csv(OUTPUT / "M02_breadth.csv")
assert int(table.days.sum()) == len(eligible)
calendar = pd.MultiIndex.from_product(
    [sorted(food.Userid.unique()), pd.date_range("2016-01-01", "2016-12-31")],
    names=["Userid", "date"],
)
panel = daily.set_index(["Userid", "date"]).reindex(calendar)
observed = panel.entries.notna()
result = {
    "eligible_days": len(eligible),
    "excluded_days": len(daily) - len(eligible),
    "calendar_user_days": len(panel),
    "observed_days": int(observed.sum()),
    "unobserved_days": int((~observed).sum()),
    "mean_variety_day_weighted": float(eligible.variety.mean()),
    "mean_variety_user_weighted": float(
        eligible.groupby("Userid").variety.mean().mean()
    ),
    "daily_total_equal_item_sum_within_one": int(
        (daily.food_calories - daily.item_calories).abs().le(1).sum()
    ),
    "daily_total_comparison_days": len(daily),
    "breadth": table.reset_index().to_dict(orient="records"),
}
save_result("M02", result)
print(json.dumps(result, indent=2))
\end{Verbatim}

\hypertarget{M02-mhealth-nutrition-knowledge-and-sample-restrictions}{%
\section{mHealth: Nutrition Knowledge and Sample
Restrictions}\label{M02-mhealth-nutrition-knowledge-and-sample-restrictions}}

\noindent An mHealth analyst evaluating suggestions about diet must determine
whether the records measure food selection, complete daily intake, or
usual intake over a longer period. A day with one logged snack supplies
evidence about that snack. It does not establish the person's total
intake for the day. The National Cancer Institute's
\href{https://epi.grants.cancer.gov/dietary-assessment-primer/profiles/record/}{food-record
guidance} identifies recording burden and reactivity as features of the
measurement process. Its
\href{https://epi.grants.cancer.gov/dietary-assessment-primer/learn/usual.html}{usual-intake
guidance} distinguishes day-to-day variation from a person's longer-run
average. Repeated logs reduce some sampling variation while selective
omissions can persist.

Nutrition studies often collect information that makes filtering more
defensible. The
\href{https://wwwn.cdc.gov/Nchs/Data/Nhanes/Public/2021/DataFiles/DR1TOT_L.htm}{NHANES
dietary codebook} documents recall reliability through
\texttt{DR1DRSTZ}. mHealth has no corresponding completeness indicator.
Several food entries or a positive calorie total cannot be substituted
for that validation. NHANES also permits reliable records involving
fasting or only water, so a zero total cannot universally be classified
as erroneous. A ratio requiring a positive denominator nevertheless
excludes zero totals by mathematical necessity.

For an adult descriptive analysis, the base sample requires age of at
least 18 on January 1, positive finite reported daily calories, and a
defined calorie-share variety measure. The fixed-age restriction removes
users who were 17 at the reference date even if they turned 18 later in
the year. It retains older adults; age 65 is not a general nutritional
validity boundary. The combined restrictions retain 75,855 days from 367
users. They define the population and the calculations; they do not
certify complete intake.

The National Cancer Institute's
\href{https://epi.grants.cancer.gov/dietary-assessment-primer/learn/outliers.html}{outlier
guidance} supports examining extreme observations and comparing analyses
under alternative exclusions. Short-term intake can genuinely vary
substantially. A population percentile screen is therefore a sensitivity
analysis. The 1st and 99th percentiles of positive adult daily totals in
January through June are 209.01 and 3,534.11. Fixing those cutoffs using
the earlier period prevents later observations from choosing their own
reference range. A second screen uses an explicitly imposed 100--5,000
range; a third requires item totals and reported totals to agree within
one calorie unit.

\begin{table}[H]
\centering\small
\begin{tabular}{@{}
  >{\raggedright\arraybackslash}p{(\columnwidth - 6\tabcolsep) * \real{0.20}}
  >{\raggedleft\arraybackslash}p{(\columnwidth - 6\tabcolsep) * \real{0.26}}
  >{\raggedleft\arraybackslash}p{(\columnwidth - 6\tabcolsep) * \real{0.26}}
  >{\raggedleft\arraybackslash}p{(\columnwidth - 6\tabcolsep) * \real{0.28}}@{}}
\toprule
Rule Applied to the Adult Base & Days Retained & Mean Reported Calories
& Mean Variety \\
\midrule

Adult Positive Base & 75,855 & 1,541.82 & 0.7837 \\
Imposed 100--5,000 Range & 75,523 & 1,543.04 & 0.7863 \\
Earlier-Period 1st--99th Percentiles & 74,224 & 1,531.07 & 0.7916 \\
Item/Day Agreement Within One & 52,707 & 1,493.19 & 0.7671 \\
\bottomrule
\end{tabular}
\end{table}

Each row applies one screen independently to the same base. The
agreement rule removes 23,148 days, far more than the percentile rule's
1,631. Its lower mean variety could reflect the kinds of records
retained, including differences in logging intensity. Agreement is
useful evidence about accounting consistency, but it does not prove
nutritional completeness. A preference for cleaner-looking totals can
therefore change the population being described.

\begin{figure}[htbp]
\centering
\begin{tikzpicture}
\begin{axis}[xlabel={Adult Base Days Retained (Percent)},ylabel={Mean Recorded Variety},xmin=65,xmax=104,ymin=.76,ymax=.803,xtick={70,80,90,100},ytick={.765,.775,.785,.795},yticklabel style={/pgf/number format/fixed,/pgf/number format/precision=3},grid=major]
\addplot[only marks,mark=*,mark size=3pt,IlliniBlue] coordinates {(100,.7837115)};
\addplot[only marks,mark=square*,mark size=3pt,Patina] coordinates {(99.5623,.7862868)};
\addplot[only marks,mark=triangle*,mark size=3.5pt,Industrial] coordinates {(97.8498,.7916456)};
\addplot[only marks,mark=diamond*,mark size=4pt,IlliniOrange] coordinates {(69.4839,.7671419)};
\node[anchor=north east,font=\footnotesize] at (axis cs:100,.7828) {Adult positive base};
\node[anchor=east,font=\footnotesize] at (axis cs:98.7,.7863) {100--5,000 range};
\node[anchor=south east,font=\footnotesize] at (axis cs:98,.793) {Earlier-period percentiles};
\node[anchor=south west,font=\footnotesize] at (axis cs:69.8,.768) {Item/day agreement};
\end{axis}
\end{tikzpicture}

\caption{Alternative screens retain different samples and change measured variety.}\label{fig:M02_mhealth_screening}
\end{figure}

The horizontal movement measures retention of adult base days; the
vertical movement measures their mean recorded variety. The
accounting-consistency restriction produces the largest movement in both
dimensions. Treating it as a harmless technical correction would obscure
a substantive change in the sample.

\hypertarget{M02-energy-plausibility-requires-physiological-inputs}{%
\subsection{Energy Plausibility Requires Physiological
Inputs}\label{M02-energy-plausibility-requires-physiological-inputs}}

\noindent Energy-plausibility methods compare reported intake with expenditure or
a model of requirements.
\href{https://doi.org/10.3389/fnut.2017.00045}{Banna, McCrory,
Fialkowski, and Boushey (2017)} compare fixed cutoffs with
individualized methods.
\href{https://doi.org/10.1038/sj.ijo.0801376}{Black (2000)}, especially
Table 8, updates the Goldberg approach. For an individual under the
method's energy-balance and measurement conditions, bounds for mean
reported energy intake divided by basal metabolic rate are

\[
L,U=\operatorname{PAL}\exp\left(\mp 2\sqrt{\frac{(0.23)^2}{d}+(0.085)^2+(0.15)^2}\right).
\]

Here \(d\) is the number of assessment days and PAL is total energy
expenditure divided by basal metabolic rate. The coefficients represent
the variation assumptions for intake, predicted basal metabolic rate,
and activity. The factor 2 gives the paper's approximate 95\% limits for
an individual; the sample-size term is one. These are conditional
screening bounds, not an individual's recommended intake. Appropriate
PAL and basal-metabolic-rate inputs are essential.

With an imposed basal metabolic rate of 1,500 kcal/day and PAL of 1.55,
the bounds are approximately 1,308--4,131 kcal/day for one assessment
day and 1,580--3,421 for seven days. Raising PAL to 1.80 raises the
seven-day bounds to 1,835--3,973. A hypothetical seven-day mean of 1,600
would fall within the first interval and below the second. The same
intake can receive a different flag because of the activity assumption.
Additional days reduce the intake-variation component while the other
variance components remain.

\begin{figure}[htbp]
\centering
\begin{tikzpicture}
\begin{axis}[title={Hypothetical BMR: 1,500 kcal/day},xlabel={Mean Reported Energy (kcal/day)},xmin=1000,xmax=5100,ymin=.4,ymax=4.8,ytick={1,2,3,4},yticklabels={7 Days; PAL 1.80,7 Days; PAL 1.55,1 Day; PAL 1.80,1 Day; PAL 1.55},width=11cm,xmajorgrids=true]
\addplot[IlliniBlue,very thick,mark=|,mark size=7pt] coordinates {(1308.43,4)(4131.38,4)};
\addplot[IlliniOrange,very thick,dashed,mark=|,mark size=7pt] coordinates {(1519.47,3)(4797.73,3)};
\addplot[IlliniBlue,very thick,mark=|,mark size=7pt] coordinates {(1580.20,2)(3420.86,2)};
\addplot[IlliniOrange,very thick,dashed,mark=|,mark size=7pt] coordinates {(1835.07,1)(3972.61,1)};
\addplot[Storm,densely dotted,thick] coordinates {(1600,.5)(1600,4.6)};
\node[anchor=south west,font=\footnotesize] at (axis cs:1650,4.45) {Hypothetical mean: 1,600};
\end{axis}
\end{tikzpicture}

\caption{Assessment duration narrows an energy-plausibility interval while assumed activity shifts it.}\label{fig:M02_mhealth_goldberg}
\end{figure}

The code calculates these hypothetical bounds separately from the
real-data screens. mHealth does not supply validated activity, basal
metabolic rate, complete assessment days, or a documented mapping of
\texttt{Gender} to the sex variable required by common prediction
equations. No participant is assigned an individualized Goldberg label.
Weight-loss goals and perturbed weights also complicate an
energy-balance interpretation. Applying the equations to invented
physiological inputs would turn an assumption into an apparent
observation.

An AI coding request should therefore name the assessment instrument,
population, units, necessary inputs, and output flags. Require the
program to preserve raw values, distinguish overlapping exclusions, and
return the analysis with and without the screen. Ask it to explain why
one-day outliers, genuine fasting, and incomplete logging require
different interpretations. The
\href{../01_Resources/12_Nutrition_and_Dietary_Data_Quality/README.md}{nutrition
resource guide} supplies the original methods and downloadable readings.

\FloatBarrier
\Needspace{8\baselineskip}
Python source:
\href{../03_Data_and_Code/mHealth/recurring_examples/nutrition.py}{nutrition.py}.
From its directory, run \texttt{python\ nutrition.py}.

\begin{Verbatim}[breaklines,breakanywhere,fontsize=\small]
"""Domain-informed sample definitions and explicitly hypothetical energy screens."""

import json
import math
import numpy as np
from common import load_inputs, save_result


def goldberg_bounds(days, pal, cv_intake=23, cv_bmr=8.5, cv_pal=15):
    """Individual EI:BMR bounds with Black (2000) coefficients, z=2, n=1."""
    if days < 1 or pal <= 0:
        raise ValueError("Positive assessment duration and PAL are required")
    variation = math.sqrt(cv_intake**2 / days + cv_bmr**2 + cv_pal**2) / 100
    return pal * math.exp(-2 * variation), pal * math.exp(2 * variation)


def run_analysis():
    food, daily, baskets = load_inputs()
    age_counts = food.groupby("Userid").Age_at_2016_01_01.nunique(dropna=False)
    assert age_counts.eq(1).all()
    ages = food.groupby("Userid").Age_at_2016_01_01.first()
    daily["age"] = daily.Userid.map(ages)
    # Nonnegative category calories and a positive denominator define valid shares.
    base = daily.loc[
        daily.variety.notna()
        & daily.food_calories.gt(0)
        & np.isfinite(daily.food_calories)
        & daily.age.ge(18)
    ].copy()
    training = base.loc[base.date.lt("2016-07-01")]
    lower, upper = training.food_calories.quantile([0.01, 0.99])
    masks = {
        "adult_positive_base": base.food_calories.gt(0),
        "fixed_100_to_5000_sensitivity": base.food_calories.between(100, 5000),
        "prior_period_1st_to_99th_percentile": base.food_calories.between(lower, upper),
        "daily_item_agreement_within_one": (base.food_calories - base.item_calories)
        .abs()
        .le(1),
    }
    summaries = []
    for name, mask in masks.items():
        selected = base.loc[mask]
        summaries.append(
            {
                "rule": name,
                "days": len(selected),
                "users": int(selected.Userid.nunique()),
                "removed_from_base": len(base) - len(selected),
                "mean_recorded_calories": float(selected.food_calories.mean()),
                "mean_variety": float(selected.variety.mean()),
            }
        )
    # Hypothetical complete records, stable energy balance, and predicted BMR.
    # These BMR/PAL inputs are imposed; no mHealth participant is classified.
    scenario = []
    for days in [1, 7]:
        for pal in [1.55, 1.80]:
            lo, hi = goldberg_bounds(days, pal)
            scenario.append(
                {
                    "days": days,
                    "assumed_pal": pal,
                    "bmr_kcal": 1500,
                    "lower_ratio": lo,
                    "upper_ratio": hi,
                    "lower_kcal": 1500 * lo,
                    "upper_kcal": 1500 * hi,
                    "illustrative_1600_kcal_flag": 1600 / 1500 < lo,
                }
            )
    result = {
        "source_observed_days": len(daily),
        "base_days": len(base),
        "base_removed_days": len(daily) - len(base),
        "age_min": float(ages.min()),
        "age_max": float(ages.max()),
        "training_days": len(training),
        "quantile_lower": float(lower),
        "quantile_upper": float(upper),
        "sensitivity": summaries,
        "hypothetical_goldberg": scenario,
    }
    save_result("nutrition", result)
    print(json.dumps(result, indent=2))


if __name__ == "__main__":
    run_analysis()
\end{Verbatim}

\hypertarget{M02-python-calculations-for-the-synthetic-cases}{%
\section{Python Calculations for the Synthetic
Cases}\label{M02-python-calculations-for-the-synthetic-cases}}

\noindent The inputs below are imposed teaching values. Running the program
creates the labeled synthetic CSV files and prints the corresponding
calculations. They are separate from the mHealth estimates. Keep the
script beside \texttt{example\_helpers.py} in
\texttt{03\_Data\_and\_Code/Lecture\_Examples}.

\FloatBarrier
\Needspace{8\baselineskip}
Python source:
\href{../03_Data_and_Code/Lecture_Examples/m02_synthetic.py}{m02\_synthetic.py}.
From its directory, run \texttt{python\ m02\_synthetic.py}.

\begin{Verbatim}[breaklines,breakanywhere,fontsize=\small]
"""Synthetic numerical examples for M2; run from any working directory."""

import itertools
import json
import math
from example_helpers import write_csv, ols_slope, shares

r = {}
campaigns = [
    ("A", "new", 100, 9),
    ("A", "returning", 100, 80),
    ("B", "new", 200, 20),
    ("B", "returning", 100, 81),
]
write_csv(
    "M02_campaigns_synthetic.csv",
    ["campaign", "customer_type", "assigned_customers", "conversions"],
    campaigns,
)
r["M02"] = {
    "observed_A": 89 / 200,
    "observed_B": 101 / 300,
    "common_mix_A": (0.09 + 0.80) / 2,
    "common_mix_B": (0.10 + 0.81) / 2,
    "basket_jaccard_high": 3 / 5,
    "basket_jaccard_low": 1 / 7,
}
assert r["M02"]["observed_A"] > r["M02"]["observed_B"]
assert r["M02"]["common_mix_B"] > r["M02"]["common_mix_A"]
variety_shares = {"equal": [0.5, 0.5], "concentrated": [0.9, 0.1]}
r["M02"]["variety"] = {
    name: 1 - sum(s * s for s in vals) for name, vals in variety_shares.items()
}
write_csv(
    "M02_variety_synthetic.csv",
    ["basket", "category", "spending_share"],
    [
        (name, f"C{i+1}", share)
        for name, vals in variety_shares.items()
        for i, share in enumerate(vals)
    ],
)
assert math.isclose(r["M02"]["variety"]["equal"], 0.5)
assert math.isclose(r["M02"]["variety"]["concentrated"], 0.18)
print(json.dumps(r, indent=2))
\end{Verbatim}

\hypertarget{M02-mhealth-recency-and-frequency-at-a-fixed-date}{%
\section{mHealth: Recency and Frequency at a Fixed
Date}\label{M02-mhealth-recency-and-frequency-at-a-fixed-date}}

\noindent A manager considering a reminder campaign needs to distinguish a
customer who logs frequently from one who logged recently. Both
customers may have the same total number of records. Recency measures
elapsed time since the last observed activity, while frequency measures
repeated activity within a specified window. Their interpretation
depends on the date at which the manager must make the decision.

For every user with a recorded day by June 30, define recency as the
number of days from the most recent recorded day to June 30. Define
frequency as distinct recorded days during June 1--30. A customer who
last logged on June 30 has recency zero. A customer with no June logs
has frequency zero and remains eligible because an earlier record
establishes membership. July activity is an outcome and enters neither
feature.

The 358 eligible users are grouped using prespecified recency boundaries
of 7 and 30 days and a frequency boundary of 15 days. These cutoffs
define a descriptive comparison; they have not been optimized for a
campaign. The empty high-frequency, high-recency cell follows from the
definitions: a customer whose last log was more than 30 days ago has
zero June logging days.

\begin{table}[H]
\centering\small
\begin{tabular}{@{}
  >{\raggedright\arraybackslash}p{(\columnwidth - 8\tabcolsep) * \real{0.16}}
  >{\raggedright\arraybackslash}p{(\columnwidth - 8\tabcolsep) * \real{0.16}}
  >{\raggedleft\arraybackslash}p{(\columnwidth - 8\tabcolsep) * \real{0.22}}
  >{\raggedleft\arraybackslash}p{(\columnwidth - 8\tabcolsep) * \real{0.22}}
  >{\raggedleft\arraybackslash}p{(\columnwidth - 8\tabcolsep) * \real{0.24}}@{}}
\toprule
Recency on June 30 & June Logging Days & Users & Any July Log & Mean
July Logging Days \\
\midrule

0--7 Days & 0--14 & 68 & 92.6\% & 13.90 \\
0--7 Days & 15--30 & 220 & 96.8\% & 22.88 \\
8--30 Days & 0--14 & 34 & 67.6\% & 7.26 \\
8--30 Days & 15--30 & 6 & 83.3\% & 9.67 \\
31+ Days & 0--14 & 30 & 50.0\% & 5.70 \\
\bottomrule
\end{tabular}
\end{table}

Recent frequent loggers have a higher July logging rate in this cohort.
The six-user cell warrants particular care because a single outcome
changes its rate substantially. These rates describe future activity
among known users; they do not identify which users would respond most
to a reminder. Selecting recipients requires an incremental-effect
comparison and a cost per contact.

\href{https://github.com/joaolcorreia/RFM-analysis}{Joao Correia's RFM
example} combines recency, frequency, and monetary value in transaction
data. mHealth supports the first two dimensions. Calories, entry counts,
and body weight have units different from spending and cannot supply
monetary value. A commercial RFM table would need an actual monetary
variable, a transaction definition, and a consistent observation window.

A coding assistant should receive the snapshot date and the eligible
population explicitly. Ask it to verify that adding July rows leaves
both June features unchanged, that repeated food entries on one day
count only once toward frequency, and that zero-log members remain in
the outcome denominator. Quantile bins require extra care when many
customers share the same integer count; the fixed boundaries here avoid
assigning tied counts to different groups.

\FloatBarrier
\Needspace{8\baselineskip}
Python source:
\href{../03_Data_and_Code/mHealth/recurring_examples/recency.py}{recency.py}.
From its directory, run \texttt{python\ recency.py}.

\begin{Verbatim}[breaklines,breakanywhere,fontsize=\small]
"""Describe later logging using features available at a fixed snapshot."""

import json
import pandas as pd
from common import OUTPUT, load_food, make_daily, save_result


def snapshot(daily, cutoff):
    cutoff = pd.Timestamp(cutoff)
    history = daily.loc[daily.date <= cutoff]
    last = history.groupby("Userid").date.max()
    recent = history.loc[history.date > cutoff - pd.Timedelta(days=30)]
    counts = recent.groupby("Userid").date.nunique().reindex(last.index, fill_value=0)
    return pd.DataFrame(
        {"recency_days": (cutoff - last).dt.days, "frequency_30d": counts}
    )


if __name__ == "__main__":
    daily = make_daily(load_food())
    customers = snapshot(daily, "2016-06-30")
    customers["recency_group"] = pd.cut(
        customers.recency_days,
        [-1, 7, 30, float("inf")],
        labels=["0-7 days", "8-30 days", "31+ days"],
    )
    customers["frequency_group"] = customers.frequency_30d.ge(15).map(
        {True: "15-30 days", False: "0-14 days"}
    )
    july = daily.loc[daily.date.between("2016-07-01", "2016-07-31")]
    customers["july_days"] = (
        july.groupby("Userid").date.nunique().reindex(customers.index, fill_value=0)
    )
    customers["july_any_log"] = customers.july_days.gt(0)
    assert customers.frequency_30d.between(0, 30).all()
    summary = (
        customers.groupby(["recency_group", "frequency_group"], observed=True)
        .agg(
            users=("july_days", "size"),
            july_users=("july_any_log", "sum"),
            july_logging_probability=("july_any_log", "mean"),
            mean_july_days=("july_days", "mean"),
        )
        .reset_index()
    )
    result = {
        "snapshot": "2016-06-30",
        "users": len(customers),
        "groups": summary.to_dict(orient="records"),
    }
    save_result("recency", result)
    summary.to_csv(OUTPUT / "recency.csv", index=False)
    print(json.dumps(result, indent=2))
\end{Verbatim}

\hypertarget{M03-m3.-regression-interpretation-and-marketing-decisions}{%
\chapter{Regression, Interpretation, and Marketing
Decisions}\label{M03-m3.-regression-interpretation-and-marketing-decisions}}

\noindent A retailer observes that stores with larger advertising budgets sell
more. The manager wants to know whether another dollar of advertising
would be profitable. The observed relationship combines the response to
advertising with differences in store size, local demand, seasonality,
and the firm's reasons for choosing the budget. Regression provides a
way to summarize conditional relationships. Its value depends on which
variation identifies the comparison and whether that comparison answers
the manager's question.

\hypertarget{M03-a-conditional-mean-and-a-regression-line}{%
\section{A Conditional Mean and a Regression
Line}\label{M03-a-conditional-mean-and-a-regression-line}}

\noindent Let \(Y_i\) be weekly store sales measured in thousands of dollars and
\(X_i\) be weekly advertising expenditure in thousands of dollars. A
simple linear model writes

\[
Y_i=\beta_0+\beta_1X_i+\varepsilon_i.
\]

The intercept \(\beta_0\) is the fitted sales level at zero advertising;
its practical interpretation is weak when zero lies far outside the
observed range. The slope \(\beta_1\) is the change in fitted sales, in
thousands of dollars, associated with an additional thousand dollars of
advertising. The residual captures the difference between observed sales
and the fitted line. Linearity is an approximation to the relationship
over the relevant range, not a claim that every customer's behavior is
linear.

Ordinary least squares chooses coefficients to minimize the sum of
squared residuals. In a simple regression with an intercept, its slope
is

\[
\hat\beta_1=\frac{\sum_i(X_i-\bar X)(Y_i-\bar Y)}{\sum_i(X_i-\bar X)^2}.
\]

The numerator measures whether high-advertising observations tend to
have high sales. The denominator measures how much advertising varies.
When advertising hardly varies, the data provide little information
about a slope, even if the sample contains many rows. A large sample of
nearly identical budgets does not substitute for informative variation.

For three observations \((X,Y)=(1,3),(2,5),(3,7)\), the fitted line is
\(\hat Y=1+2X\). At \(X=2.5\), it predicts sales of \$6,000. If the
sales contribution margin is 30\%, an additional \$1,000 of advertising
is associated with \$600 of contribution from \$2,000 of extra sales.
Treating that association as a return to intervention requires an
identification argument. Even under a causal interpretation, the
incremental contribution is less than the additional advertising expense
in this numerical example.

\hypertarget{M03-multiple-regression-and-the-comparison-it-uses}{%
\section{Multiple Regression and the Comparison It
Uses}\label{M03-multiple-regression-and-the-comparison-it-uses}}

\noindent A multiple regression can include observed store characteristics:

\[
Y_i=\beta_0+\beta_1X_i+\gamma Z_i+\varepsilon_i,
\]

where \(Z_i\) might be floor area measured in thousands of square feet.
The advertising coefficient summarizes variation in sales associated
with advertising after removing the linear association with floor area.
A useful way to understand this is to regress advertising on the other
controls, regress sales on the same controls, and regress the resulting
sales residuals on the advertising residuals. The Frisch-Waugh-Lovell
result gives the same advertising coefficient as the full regression
when the residualization includes the same intercept and controls.

``Controlling for'' a variable therefore changes the comparison. It does
not automatically remove all confounding. If managers observe local
demand information that the analyst lacks, their advertising choices may
remain related to unobserved sales determinants. Controls measured after
the intervention can also remove part of the effect or introduce bias.
Select variables using the decision timing and a substantive account of
how the data were generated.

When an omitted variable affects both the predictor and the outcome, the
simple slope can be misleading. In the linear population model above,
suppose the regression disturbance is mean-independent of \(X\) and
\(Z\). Omitting \(Z\) yields a simple-regression slope equal to

\[
\beta_1+\gamma\frac{\operatorname{Cov}(X,Z)}{\operatorname{Var}(X)}.
\]

The extra term has a substantive interpretation. If larger stores sell
more, so \(\gamma>0\), and larger stores spend more on advertising, so
the covariance is positive, the simple advertising slope exceeds the
conditional slope. The equation is a model-based decomposition; its
interpretation depends on the stated linear specification and exogeneity
condition.

\hypertarget{M03-case-study-attractive-markets-and-apparent-competitive-benefits}{%
\section{Case Study: Attractive Markets and Apparent Competitive
Benefits}\label{M03-case-study-attractive-markets-and-apparent-competitive-benefits}}

\noindent \href{https://doi.org/10.3982/QE478}{Igami and Yang (2016)} study entry
and competition among hamburger chains in Canada. Their setting
motivates a general identification problem: attractive markets can
support both more competitors and better outcomes for a focal firm. A
positive relationship between rival presence and performance can
therefore coexist with a negative competitive effect.

Within low-demand markets, suppose stores facing one, two, and three
rivals have sales of 9, 8, and 7 units. Within high-demand markets,
stores facing four, five, and six rivals have sales of 16, 15, and 14
units. Within each demand group, one additional rival is associated with
one unit less sales. Across the six observations, the pooled slope is
approximately \(+1.57\): high-demand markets have more rivals and higher
sales.

\begin{figure}[htbp]
\centering
% Original synthetic teaching figure; see corresponding module.
\begin{tikzpicture}
\begin{axis}[xmin=.5,xmax=6.5,ymin=5,ymax=18,xlabel={Number of rivals},ylabel={Sales (illustrative units)},legend style={at={(.5,1.08)},anchor=south,legend columns=2},xtick={1,...,6},ymajorgrids=true]
\addplot[IlliniBlue,mark=*,very thick] coordinates {(1,9)(2,8)(3,7)};
\addplot[IlliniOrange,mark=square*,very thick] coordinates {(4,16)(5,15)(6,14)};
\addplot[Storm,dashed,domain=1:6] {6+11*x/7};
\legend{Low demand,High demand,Pooled fit}
\node[anchor=west,font=\small] at (axis cs:1.1,6.3) {Within-group slope: $-1$};
\node[anchor=west,font=\small] at (axis cs:3.4,17) {Pooled slope: $+1.57$};
\end{axis}
\end{tikzpicture}

\caption{A positive pooled relationship can conceal a negative within-market-type relationship.}\label{fig:M03_omitted_demand}
\end{figure}

The within-group slope is \(-1\), while baseline demand differs between
groups. The positive pooled slope reflects that demand composition. In
an actual retail application, market demand may be unobserved, entry may
be forward-looking, and firms may respond strategically.

A manager considering entry should ask which outcomes would change if an
additional competitor entered the same market. Comparing different
markets does not automatically provide that counterfactual. Market fixed
effects, measured demand controls, an instrument, or a credible event
design may improve the comparison, but each uses different variation and
requires its own conditions. A highly predictive regression can remain
unsuitable for interpreting the coefficient as competition's causal
effect.

\hypertarget{M03-transformations-and-economically-meaningful-units}{%
\section{Transformations and Economically Meaningful
Units}\label{M03-transformations-and-economically-meaningful-units}}

\noindent A level-level specification gives an additive interpretation: one more
unit of \(X\) changes fitted \(Y\) by \(\beta_1\). A log-level
specification, \(\log Y=\beta_0+\beta_1X+\varepsilon\), describes
proportional differences in a positive outcome. A one-unit change in
\(X\) changes the fitted log outcome by \(\beta_1\); the corresponding
multiplicative change in the log-scale prediction is \(e^{\beta_1}\), or
\(100(e^{\beta_1}-1)\) percent. The approximation \(100\beta_1\) percent
is useful when the coefficient is small.

A log-log demand model, \(\log Q=\alpha+\beta\log P+\varepsilon\), gives
\(\beta\) an elasticity interpretation over the range where the
specification applies. If \(\beta=-1.5\), a 1\% higher price is
associated with approximately 1.5\% lower quantity. Sales revenue is
\(PQ\), so its local price elasticity is \(1+\beta=-0.5\). Profit also
depends on unit cost. A price that raises revenue can still lower
profit, and observed prices can be endogenous because firms set them in
response to demand.

Zero and negative outcomes need deliberate treatment. Taking a logarithm
of zero is undefined. Adding one changes the scale and interpretation,
especially when values are small or the unit of measurement changes. A
two-part model can separate purchase incidence from positive
expenditure. A suitable generalized linear model may model the
conditional mean directly. The choice should reflect the outcome and
intended estimand, with diagnostics for the selected specification.

Back-transforming a predicted log outcome also needs care.
Exponentiating the predicted conditional mean of \(\log Y\) generally
does not produce the conditional mean of \(Y\), because exponentiation
is nonlinear. Retransformation methods or a model for the mean on the
original scale may be needed when total sales is the target. Evaluate
predictions in the units used for the managerial decision.

\hypertarget{M03-interactions-describe-how-a-relationship-changes}{%
\section{Interactions Describe How a Relationship
Changes}\label{M03-interactions-describe-how-a-relationship-changes}}

\noindent Suppose an analyst expects a promotional message to work differently for
customers with long tenure. A linear interaction specification is

\[
Y_i=\beta_0+\beta_1D_i+\beta_2T_i+\beta_3D_iT_i+\varepsilon_i,
\]

where \(D_i\) is a message indicator and \(T_i\) is tenure in years. The
fitted message difference at tenure \(t\) is \(\beta_1+\beta_3t\). The
coefficient \(\beta_1\) is the difference at zero tenure, so centering
tenure at a meaningful value can make that coefficient easier to
interpret. The interaction coefficient tells how the fitted difference
changes per year of tenure.

With \(\beta_1=4\) dollars and \(\beta_3=-1\) dollar per year, the
fitted difference is \$4 at tenure zero, \$2 at two years, and zero at
four years. An average coefficient can conceal that variation.
Predictions beyond the observed tenure range should be treated as
extrapolation. In a nonlinear probability model, an interaction
coefficient on the model's index does not directly equal the interaction
in probability; compute and plot the predicted probabilities or marginal
contrasts.

\begin{figure}[htbp]
\centering
% Original synthetic teaching figure; see corresponding module.
\begin{tikzpicture}
\begin{axis}[xmin=0,xmax=5,ymin=-1.6,ymax=5,xlabel={Customer tenure (years)},ylabel={Fitted message difference (dollars)},xtick={0,...,5},ytick={-1,0,1,2,3,4},ymajorgrids=true]
\addplot[IlliniOrange,very thick,domain=0:5] {4-x};
\addplot[IlliniBlue,dashed] coordinates {(0,0)(5,0)};
\addplot[IlliniBlue,only marks,mark=*] coordinates {(0,4)(2,2)(4,0)};
\node[anchor=south west] at (axis cs:.1,4) {$\Delta(t)=4-t$};
\node[anchor=south west,align=left] at (axis cs:3.3,.6) {Zero contrast\\at four years};
\draw[edge] (axis cs:4,.6)--(axis cs:4,.1);
\end{axis}
\end{tikzpicture}

\caption{An interaction changes the relevant comparison across customer tenure.}\label{fig:M03_interaction}
\end{figure}

For tenure between zero and five years, the assumed contrast \(4-t\)
crosses zero at four years. In an estimated application, uncertainty
around the contrast and support at each tenure would determine whether
such a crossing is informative.

\hypertarget{M03-uncertainty-dependence-and-practical-importance}{%
\section{Uncertainty, Dependence, and Practical
Importance}\label{M03-uncertainty-dependence-and-practical-importance}}

\noindent A standard error describes sampling uncertainty under an estimator and
its assumptions. A confidence interval provides a range generated by a
procedure with a stated repeated-sampling coverage property. A p-value
assesses compatibility with a particular null under its reference model.
None of these quantities measures the managerial value of an effect by
itself. A precisely estimated increase of one cent may be unimportant; a
wide interval around a large potential benefit may justify another
experiment.

For an approximately normal coefficient estimate, a rough 95\% interval
is the estimate plus or minus 1.96 standard errors, with appropriate
adjustments when a finite-sample reference distribution is needed. If a
sales slope is 2.0 with a standard error of 0.5, that interval is
approximately {[}1.02, 2.98{]} in the slope's units. Report the units
and the comparison alongside the interval.

Heteroskedasticity means that residual variance differs across
observations. Heteroskedasticity-robust standard errors can address that
feature for inference under suitable sampling conditions; they do not
repair omitted-variable bias. Repeated store-week observations may share
persistent store shocks. Clustered standard errors allow dependence
within a specified cluster under their asymptotic conditions. The
relevant number of independent clusters can be much smaller than the
number of rows. With few clusters, conventional approximations may be
unreliable and design-specific methods are needed.

\hypertarget{M03-diagnostics-and-panel-data}{%
\section{Diagnostics and Panel
Data}\label{M03-diagnostics-and-panel-data}}

\noindent Inspect residuals against fitted values and important predictors for
nonlinear patterns and changing variance. Investigate influential
observations, errors in units, and unusual leverage. Multicollinearity
can make individual coefficients unstable even when predictions are
reasonably stable, because several predictors explain nearly the same
variation. Removing a variable solely to obtain a smaller p-value
changes the question and can introduce omitted-variable bias.

Store fixed effects absorb time-invariant differences between stores.
Calendar fixed effects absorb common period-specific shocks. In a
store-week panel, a model with both uses changes within stores relative
to common time movements. It still requires a valid account of
time-varying confounders. If a store raises advertising precisely when
an unobserved local event increases demand, the advertising coefficient
can remain biased.

A fitted model should be evaluated on the task for which it is used. For
prediction, compare held-out errors with a useful baseline. For a
descriptive coefficient, inspect support and sensitivity to
substantively justified specifications. For a causal effect, assess the
design and identifying conditions. A large \(R^2\) indicates in-sample
fit relative to an intercept-only benchmark; it does not establish a
credible intervention comparison.

For repeated food logs, a customer fixed effect can absorb persistent
differences in eating habits while relating changes in variety to
changes in recorded consumption. It leaves time-varying motivation,
activity, and reporting behavior as possible influences. The variety
analysis in \href{https://doi.org/10.1086/720453}{Nielsen and coauthors
(2022)} provides a concrete application. Percentage-change outcomes also
require attention to zero denominators: excluding days with zero prior
vegetable or snack consumption changes the sample relative to an
analysis of total calories. Interpret each coefficient for its own
observed population and outcome definition.

\hypertarget{M03-ai-assisted-implementation-and-practice}{%
\section{AI-Assisted Implementation and
Practice}\label{M03-ai-assisted-implementation-and-practice}}

\noindent Ask an assistant to fit the pooled and demand-group-adjusted regressions
in the six-observation example. Require a table with variable
definitions, units, coefficients, and the exact sample. Verify that the
adjusted rival coefficient is \(-1\) and the pooled coefficient is
\(11/7\approx1.57\). Ask it to explain why adding demand-group
indicators changes the comparison. A numerical match without that
explanation is an incomplete interpretation.

For a larger dataset, specify the formula and desired covariance
estimator before requesting code. Ask the assistant to list dropped rows
and explain how categorical variables were encoded. Have it calculate a
predicted contrast at two values of a moderator and compare the result
with the interaction formula. Treat diagnostic recommendations as
hypotheses to inspect, rather than automatic permission to remove
observations or alter the model.

For practice, consider an advertising coefficient measured per \$1,000
of spending. Re-express the predictor in dollars and explain how the
numerical coefficient changes while fitted sales remain the same. Next
explain why a robust standard error cannot eliminate the bias in the
market-demand example. Finally, distinguish a confidence interval for
the conditional mean from a prediction interval for a new store's
realized sales: the latter also reflects individual outcome variation.

\hypertarget{M03-readings}{%
\section{Readings}\label{M03-readings}}

\noindent \href{https://www.statlearning.com/}{An Introduction to Statistical
Learning} provides a foundation for regression, model assessment, and
nonlinear extensions. The
\href{https://www.statsmodels.org/stable/regression.html}{statsmodels
regression documentation} supports implementation.

\hypertarget{M03-mhealth-breadth-and-next-day-recorded-calories}{%
\section{mHealth: Breadth and Next-Day Recorded
Calories}\label{M03-mhealth-breadth-and-next-day-recorded-calories}}

\noindent The manager observes that users with broader daily food logs often
report more calories and considers using breadth to identify customers
who may want different suggestions. The relevant first comparison
concerns timing: does today's recorded breadth predict the next calendar
day's reported total? Taking the next available row would turn a one-day
comparison into an interval of variable length when customers skip days.

Let \(K_{it}\) be today's recorded category count and \(E_{i,t+1}\) the
next calendar day's reported calories. Among positive consecutive-day
pairs, estimate

\[
\log E_{i,t+1}=\alpha_i+\beta K_{it}+\gamma_{m(t)}+\delta_{w(t)}+u_{it}.
\]

The customer effect \(\alpha_i\) absorbs time-invariant differences in
typical logging and intake. Month and weekday effects describe calendar
variation. The coefficient compares days with different recorded breadth
for the same user, after these additive controls. It remains an
association: changing appetite, omitted meals, or recording effort can
affect both breadth and later totals. The
\href{https://epi.grants.cancer.gov/dietary-assessment-primer/learn/day.html}{NCI
discussion of weekday effects} motivates attention to calendar
composition.

The sample requires positive reported totals on both dates and at least
two eligible pairs per user. It contains 69,648 pairs from 370 users.
Standard errors are clustered by user to allow dependence among
overlapping pairs. The pooled model uses the same month and weekday
controls but omits customer effects. No prior-outcome control or calorie
winsorization is included in either baseline specification.

\begin{table}[H]
\centering\small
\begin{tabular}{@{}lrrr@{}}
\toprule
Specification & Pairs & Breadth Coefficient & Clustered 95\% Interval \\
\midrule

Pooled & 69,648 & 0.02771 & {[}0.01934, 0.03609{]} \\
Customer Fixed Effects & 69,648 & 0.01284 & {[}0.00997, 0.01572{]} \\
Fixed Effects, 100--5,000 on Both Days & 69,216 & 0.01136 & {[}0.00893,
0.01380{]} \\
\bottomrule
\end{tabular}
\end{table}

In the fixed-effects model, one additional category corresponds to a
1.29\% higher conditional geometric mean of next-day reported calories,
calculated as \(100(e^{\beta}-1)\). This percentage does not
automatically describe the arithmetic mean. The pooled estimate is
larger, consistent with persistent differences across users contributing
to the pooled association. Its reduction after adding fixed effects does
not identify a causal mechanism.

The third row applies the explicitly imposed review range from M1 to
both dates. It removes 432 pairs and produces a 1.14\% geometric-mean
association. This is a sensitivity comparison on a changed sample, not a
validated nutrition screen. The baseline outcome ranges from 0.64 to
9,190, so the restriction directly addresses observations influential on
the log scale. The adult and percentile screens in M2 answer different
sample questions and should be named separately when used.

\begin{figure}[htbp]
\centering
\begin{tikzpicture}
\begin{axis}[xlabel={Log-Outcome Coefficient per Additional Category},xmin=0,xmax=.04,scaled x ticks=false,xtick={0,.01,.02,.03,.04},xticklabel style={/pgf/number format/fixed,/pgf/number format/precision=2},ymin=.5,ymax=3.5,ytick={1,2,3},yticklabels={Restricted Fixed Effects,User Fixed Effects,Pooled},width=11cm,xmajorgrids=true]
\addplot[IlliniBlue,very thick,mark=|,mark size=6pt] coordinates {(.0193357,3)(.0360878,3)};
\addplot[IlliniOrange,very thick,mark=|,mark size=6pt] coordinates {(.0099734,2)(.0157151,2)};
\addplot[Patina,very thick,dashed,mark=|,mark size=6pt] coordinates {(.0089318,1)(.0137982,1)};
\addplot[only marks,IlliniBlue,mark=*,mark size=3pt] coordinates {(.0277117,3)};
\addplot[only marks,IlliniOrange,mark=square*,mark size=3pt] coordinates {(.0128443,2)};
\addplot[only marks,Patina,mark=triangle*,mark size=3pt] coordinates {(.011365,1)};
\end{axis}
\end{tikzpicture}

\caption{Customer controls and a calorie-range sensitivity restriction change the descriptive coefficient.}\label{fig:M03_mhealth_association}
\end{figure}

The intervals quantify sampling uncertainty under the regression and
clustering specification. They do not account for all reporting error or
make breadth an assigned treatment. A recommendation to induce greater
breadth needs evidence about how a particular intervention changes
behavior and outcomes.

The date shift below moves future outcomes back exactly one day before
the merge. The program then applies each sample rule explicitly and
reports pair and user counts. An AI-generated
\texttt{groupby(...).shift(-1)} without a date-gap check would answer a
different question.

\FloatBarrier
\Needspace{8\baselineskip}
Python source:
\href{../03_Data_and_Code/mHealth/recurring_examples/m03.py}{m03.py}.
From its directory, run \texttt{python\ m03.py}.

\begin{Verbatim}[breaklines,breakanywhere,fontsize=\small]
"""Estimate descriptive next-calendar-day associations with clustered uncertainty."""

import json
import numpy as np
import pandas as pd
import statsmodels.formula.api as smf
from common import load_inputs, save_result

food, daily, baskets = load_inputs()
future = daily[["Userid", "date", "food_calories"]].copy()
future["date"] -= pd.Timedelta(days=1)
future = future.rename(columns={"food_calories": "next_calories"})
panel = daily.merge(future, on=["Userid", "date"], validate="one_to_one")
panel = panel.loc[panel.food_calories.gt(0) & panel.next_calories.gt(0)].copy()
counts = panel.groupby("Userid").size()
panel = panel.loc[panel.Userid.isin(counts[counts >= 2].index)].copy()
panel["log_next"] = np.log(panel.next_calories)
panel["month"] = panel.date.dt.month
panel["weekday"] = panel.date.dt.dayofweek
assert (panel.next_calories > 0).all()
restricted = panel.loc[
    panel.food_calories.between(100, 5000) & panel.next_calories.between(100, 5000)
].copy()
counts = restricted.groupby("Userid").size()
restricted = restricted.loc[restricted.Userid.isin(counts[counts >= 2].index)]
fits = {}
base = "log_next ~ category_count + C(month) + C(weekday)"
specifications = [
    ("pooled", base, panel),
    ("user_fixed_effects", base + " + C(Userid)", panel),
    ("restricted_fixed_effects", base + " + C(Userid)", restricted),
]
for name, formula, sample in specifications:
    fit = smf.ols(formula, data=sample).fit(
        cov_type="cluster", cov_kwds={"groups": sample.Userid}, use_t=True
    )
    coefficient = float(fit.params["category_count"])
    low, high = map(float, fit.conf_int().loc["category_count"])
    fits[name] = {
        "pairs": len(sample),
        "users": int(sample.Userid.nunique()),
        "coefficient": coefficient,
        "se": float(fit.bse["category_count"]),
        "ci_low": low,
        "ci_high": high,
        "geometric_mean_percent": 100 * float(np.expm1(coefficient)),
    }
result = {
    "pairs": len(panel),
    "users": int(panel.Userid.nunique()),
    "min_next_calories": float(panel.next_calories.min()),
    "max_next_calories": float(panel.next_calories.max()),
    "models": fits,
}
save_result("M03", result)
print(json.dumps(result, indent=2))
\end{Verbatim}

\hypertarget{M03-python-calculations-for-the-synthetic-cases}{%
\section{Python Calculations for the Synthetic
Cases}\label{M03-python-calculations-for-the-synthetic-cases}}

\noindent The inputs below are imposed teaching values. Running the program
creates the labeled synthetic CSV files and prints the corresponding
calculations. They are separate from the mHealth estimates. Keep the
script beside \texttt{example\_helpers.py} in
\texttt{03\_Data\_and\_Code/Lecture\_Examples}.

\FloatBarrier
\Needspace{8\baselineskip}
Python source:
\href{../03_Data_and_Code/Lecture_Examples/m03_synthetic.py}{m03\_synthetic.py}.
From its directory, run \texttt{python\ m03\_synthetic.py}.

\begin{Verbatim}[breaklines,breakanywhere,fontsize=\small]
"""Synthetic numerical examples for M3; run from any working directory."""

import itertools
import json
import math
from example_helpers import write_csv, ols_slope, shares

r = {}
low = [(1, 9), (2, 8), (3, 7)]
high = [(4, 16), (5, 15), (6, 14)]
write_csv(
    "M03_rivals_synthetic.csv",
    ["demand_group", "rival_count", "sales_illustrative_units"],
    [("low", x, y) for x, y in low] + [("high", x, y) for x, y in high],
)
r["M03"] = {
    "pooled": ols_slope(low + high),
    "low": ols_slope(low),
    "high": ols_slope(high),
    "simple_sales_fit": ols_slope([(1, 3), (2, 5), (3, 7)]),
    "coefficient_ci": [2 - 1.96 * 0.5, 2 + 1.96 * 0.5],
}
assert math.isclose(r["M03"]["pooled"]["slope"], 11 / 7)
assert r["M03"]["low"]["slope"] == r["M03"]["high"]["slope"] == -1
print(json.dumps(r, indent=2))
\end{Verbatim}

\hypertarget{M03-price-elasticity-when-prices-respond-to-demand}{%
\section{Price Elasticity When Prices Respond to
Demand}\label{M03-price-elasticity-when-prices-respond-to-demand}}

\noindent A retailer considers raising a price and wants the resulting percentage
change in quantity. Historical prices can be high precisely when
managers expect strong demand. A regression of quantity on price then
combines the response to price with information used to choose it. Good
predictive fit does not separate these two sources of variation.

Consider a synthetic population of independent markets. Let \(p_i\) be
log price in dollars per unit and \(q_i\) log quantity sold. A
standardized cost shifter \(z_i\), an unobserved demand shock \(u_i\),
and independent shocks \(v_i,e_i\) are mutually independent standard
normal variables. The imposed price rule and demand equation are

\[
p_i=2+0.4z_i+0.6u_i+0.2v_i,
\qquad
q_i=5-1.5p_i+u_i+0.2e_i.
\]

The demand elasticity is \(-1.5\): a small 1\% increase in price changes
quantity by approximately \(-1.5\%\) with the demand shock held fixed.
The coefficient \(0.6\) makes price respond positively to unobserved
demand. The cost shifter changes price and has no direct effect on
quantity in this benchmark. These equations are imposed behavioral
relationships, without an estimated supply model or a solved pricing
equilibrium.

For the general linear equation \(q_i=\alpha+\beta p_i+\varepsilon_i\),
an instrumental-variable interpretation rests on the following
conditions.

\begin{assumption}[Price Instrument]
\label{ass:price-instrument}
Markets are independent draws with finite second moments. The demand equation has a common coefficient $\beta$. The cost shifter satisfies $\operatorname{Cov}(z_i,\varepsilon_i)=0$ and $\operatorname{Cov}(z_i,p_i)\ne0$.
\end{assumption}

Under Assumption \ref{ass:price-instrument}, the ratio
\(\operatorname{Cov}(z_i,q_i)/\operatorname{Cov}(z_i,p_i)\) equals
\(\beta\). The numerator measures quantity variation associated with the
cost shifter, and the denominator measures the induced price variation.
Ordinary least squares instead has probability limit
\(\beta+\operatorname{Cov}(p_i,\varepsilon_i)/\operatorname{Var}(p_i)\),
which is less negative in the benchmark because favorable demand shocks
raise price.

\begin{figure}[htbp]
\centering
\begin{tikzpicture}
\node[box,text width=2.2cm] (z) at (0,0) {Cost Shifter\\$z_i$};
\node[box,text width=2.2cm] (p) at (4.2,0) {Log Price\\$p_i$};
\node[box,text width=2.2cm] (q) at (8.4,0) {Log Quantity\\$q_i$};
\node[orangebox,text width=3cm] (u) at (5.7,2.4) {Unobserved Demand\\$u_i$};
\draw[edge] (z) -- node[above] {$0.4$} (p);
\draw[edge] (p) -- node[above] {$-1.5$} (q);
\draw[accent] (u) -- node[left] {$0.6$} (p);
\draw[accent] (u) -- node[right] {$1$} (q);
\draw[accent,dashed] (z.south) to[out=-45,in=-135] node[below,align=center,font=\footnotesize] {Direct Effect 0.5 Added Only in the Sensitivity Case} (q.south);
\end{tikzpicture}

\caption{A cost shifter changes price; unobserved demand changes both price and quantity. A direct cost-to-demand path violates exclusion.}\label{fig:M03_price_instrument}
\end{figure}

The sample contains 2,000 simulated markets with seed 534. OLS and
two-stage least squares use an intercept and the same log-price
regressor. The reported intervals use heteroskedasticity-robust HC1
covariance estimates and a normal approximation.

\begin{table}[H]
\centering\small
\begin{tabular}{@{}lrl@{}}
\toprule
Estimator and Data Construction & Elasticity & 95\% Interval \\
\midrule

OLS, Benchmark & -0.447 & {[}-0.483, -0.410{]} \\
IV, Benchmark & -1.483 & {[}-1.589, -1.378{]} \\
IV, Direct Cost Effect Added & -0.239 & {[}-0.283, -0.195{]} \\
\bottomrule
\end{tabular}
\end{table}

The third row adds \(0.5z_i\) directly to log quantity while preserving
the price rule and the structural price elasticity. Its first-stage F
statistic remains 815.85, identical to the benchmark. A strong first
stage therefore coexists with a severely misleading elasticity estimate
when exclusion fails. For example, a proposed cost instrument that also
changes product quality could affect demand through an additional path.

The
\href{https://pyblp.readthedocs.io/en/stable/_notebooks/tutorial/logit_nested.html}{PyBLP
logit tutorial} treats prices as endogenous in differentiated-product
demand. Its market-share models require additional structure beyond this
scalar simulation. The calculation below exposes the instrument logic
directly. With \(X\) containing the intercept and price and \(Z\) the
intercept and cost shifter, it forms \(\widehat X=Z(Z'Z)^{-1}Z'X\) and
computes \((\widehat X'X)^{-1}\widehat X'q\). Its covariance uses the
structural residual \(q-X\widehat\beta\), which differs from a
second-stage residual formed with fitted price.

\FloatBarrier
\Needspace{8\baselineskip}
Python source:
\href{../03_Data_and_Code/Lecture_Examples/price_iv.py}{price\_iv.py}.
From its directory, run \texttt{python\ price\_iv.py}.

\begin{Verbatim}[breaklines,breakanywhere,fontsize=\small]
"""Synthetic price endogeneity and an exclusion-restriction sensitivity."""

from pathlib import Path
import json
import numpy as np


def fit_iv(y, x, z):
    # Instruments include the intercept. Avoid constructing an n-by-n matrix.
    xhat = z @ np.linalg.lstsq(z, x, rcond=None)[0]
    bread = np.linalg.inv(xhat.T @ x)
    beta = bread @ xhat.T @ y
    residual = y - x @ beta
    scores = xhat * residual[:, None]
    covariance = bread @ (scores.T @ scores) @ bread.T
    covariance *= len(y) / (len(y) - x.shape[1])  # HC1 correction
    se = np.sqrt(np.diag(covariance))
    return {
        "elasticity": float(beta[1]),
        "robust_se": float(se[1]),
        "ci95": [float(beta[1] - 1.96 * se[1]), float(beta[1] + 1.96 * se[1])],
    }


def calculate():
    rng = np.random.default_rng(534)
    n = 2000
    cost, demand, pricing_noise, sales_noise = rng.normal(size=(4, n))
    log_price = 2 + 0.4 * cost + 0.6 * demand + 0.2 * pricing_noise
    log_quantity = 5 - 1.5 * log_price + demand + 0.2 * sales_noise
    x = np.column_stack([np.ones(n), log_price])
    z = np.column_stack([np.ones(n), cost])
    first = z @ np.linalg.lstsq(z, log_price, rcond=None)[0]
    residual_ss = np.square(log_price - first).sum()
    total_ss = np.square(log_price - log_price.mean()).sum()
    first_stage_f = (total_ss - residual_ss) / (residual_ss / (n - 2))
    result = {
        "data": "synthetic independent markets",
        "markets": n,
        "seed": 534,
        "true_elasticity": -1.5,
        "ols": fit_iv(log_quantity, x, x),
        "valid_iv": fit_iv(log_quantity, x, z),
        "invalid_iv_direct_cost_effect_0_5": fit_iv(log_quantity + 0.5 * cost, x, z),
        "first_stage_f": float(first_stage_f),
    }
    return result


if __name__ == "__main__":
    result = calculate()
    output = Path(__file__).resolve().parent / "price_iv_results.json"
    output.write_text(json.dumps(result, indent=2, allow_nan=False) + "\n")
    print(json.dumps(result, indent=2))
\end{Verbatim}

\hypertarget{M04-m4.-product-preferences-customer-segments-and-lifetime-value}{%
\chapter{Product Preferences, Customer Segments, and Lifetime
Value}\label{M04-m4.-product-preferences-customer-segments-and-lifetime-value}}

\noindent A subscription retailer must choose a product bundle, decide whom to
serve, and determine how much to spend acquiring or retaining customers.
These decisions are connected. A feature that attracts one segment may
be costly to deliver, and a segment with high initial spending may have
low retention. Conjoint analysis estimates tradeoffs among product
attributes, clustering summarizes patterns of similarity, and customer
lifetime value translates future behavior into expected contribution.

\hypertarget{M04-product-design-requires-tradeoffs}{%
\section{Product Design Requires
Tradeoffs}\label{M04-product-design-requires-tradeoffs}}

\noindent Asking consumers whether they like fast delivery often produces a
favorable response, but it does not reveal how much they will pay for it
or which other feature they will give up. Conjoint tasks present product
profiles whose attributes vary. Respondents may rate profiles, rank
them, or choose among alternatives. Choice-based tasks resemble an
assortment decision, while ratings can provide more observations per
respondent under a different response process.

Attributes should represent decisions the firm can implement and
descriptions respondents can understand. Levels should span a plausible
range without creating impossible combinations. If every premium-quality
product also has the highest price, the design provides little
independent information about the quality-price tradeoff. Randomization
and a suitable experimental design help separate attribute effects, but
realistic constraints may limit which combinations can be shown.

Include the decision's relevant alternatives. A forced choice between
two products estimates preference conditional on choosing one of them.
An outside option such as ``buy neither'' helps represent category
participation. A simulation that omits the outside option can overstate
category demand because every respondent must buy something. Survey
design, task burden, and hypothetical behavior also affect how well
stated choices transfer to market behavior.

\hypertarget{M04-utility-choice-probabilities-and-willingness-to-pay}{%
\section{Utility, Choice Probabilities, and Willingness to
Pay}\label{M04-utility-choice-probabilities-and-willingness-to-pay}}

\noindent A simple model assigns respondent \(i\) a systematic utility for
alternative \(j\):

\[
V_{ij}=\beta_q q_j+\beta_f f_j+\beta_p p_j,
\]

where \(q_j\) indicates premium quality, \(f_j\) indicates a delivery
feature, and \(p_j\) is price in dollars. Utility has an arbitrary
scale. Coefficients describe tradeoffs on that scale, and price is
expected to enter negatively in a conventional demand interpretation.
Alternative-specific constants can capture average preference for named
options relative to a reference alternative.

Under a multinomial-logit random-utility specification, the probability
of choosing \(j\) from the available set \(\mathcal J_i\) is

\[
P_{ij}=\frac{\exp(V_{ij})}{\sum_{k\in\mathcal J_i}\exp(V_{ik})}.
\]

Multinomial logit imposes proportional substitution across alternatives.
Whether this pattern fits the product category requires empirical
assessment. Repeated choices by one respondent are dependent; estimation
and validation should respect that respondent structure. Heterogeneous
preferences may be represented with interactions, latent classes, or
random coefficients when the data support the additional complexity.

For a binary feature whose coefficient is \(\beta_f\), a local
willingness-to-pay tradeoff under linear price utility is
\(-\beta_f/\beta_p\). If \(\beta_f=1.2\) and \(\beta_p=-0.1\) per
dollar, the model equates the feature with \$12 of price. This ratio is
conditional on the specification and utility scale cancellation. A price
coefficient close to zero makes the ratio unstable, and preferences
outside the experimental price range are extrapolations.

Attribute importance is often calculated from the range of an
attribute's estimated part-worths relative to the sum of ranges. The
result depends on the levels included in the study. Expanding the price
range can increase measured price importance even when underlying
preferences are unchanged. Report the design range when discussing
importance.

\hypertarget{M04-case-study-a-new-bundle-can-cannibalize-existing-sales}{%
\section{Case Study: A New Bundle Can Cannibalize Existing
Sales}\label{M04-case-study-a-new-bundle-can-cannibalize-existing-sales}}

\noindent Consider two subscription bundles and an outside option. Bundle A has
premium quality and costs \$10; bundle B has standard quality and costs
\$6. Let \(\beta_q=1.2\), \(\beta_p=-0.1\), and normalize outside-option
utility to zero. The systematic utilities are 0.2 for A and \(-0.6\) for
B. The logit probabilities are approximately 44.1\% for A, 19.8\% for B,
and 36.1\% for the outside option.

Now introduce bundle C with the same modeled attributes as A. The logit
model assigns A and C approximately 30.6\% each, B 13.7\%, and the
outside option 25.1\%. C's share is 30.6 percentage points, but the
total purchasing share rises only from 63.9\% to 74.9\%, an increase of
about 11.0 points. The remaining share comes from the existing products.
Forecasting a new product's sales without accounting for substitution
would overstate incremental demand.

\begin{figure}[htbp]
\centering
% Original synthetic teaching figure; see corresponding module.
\begin{tikzpicture}[x=11cm,y=1cm]
\node[anchor=east] at (-.02,1.5) {Before launch};
\node[anchor=east] at (-.02,0) {After launch};
\fill[IlliniBlue] (0,1.1) rectangle (.4409,1.9);
\fill[Patina] (.4409,1.1) rectangle (.639,1.9);
\fill[LightStorm] (.639,1.1) rectangle (1,1.9);
\node[text=white] at (.22045,1.5) {A: 44.1\%};
\node[text=white] at (.54,1.5) {B: 19.8\%};
\node at (.8195,1.5) {Outside: 36.1\%};
\fill[IlliniBlue] (0,-.4) rectangle (.306, .4);
\fill[Patina] (.306,-.4) rectangle (.443483,.4);
\fill[IlliniOrange] (.443483,-.4) rectangle (.749475,.4);
\fill[LightStorm] (.749475,-.4) rectangle (1,.4);
\node[text=white] at (.153,0) {A: 30.6\%};
\node[text=white,font=\scriptsize,align=center] at (.3747,0) {B:\\13.7\%};
\node[text=IlliniBlue] at (.5965,0) {C: 30.6\%};
\node[align=center] at (.87475,0) {Outside:\\25.1\%};
\draw[IlliniBlue,dashed] (.639,.95)--(.639,-.85);
\draw[<->,very thick,IlliniOrange] (.639,-.75)--(.749475,-.75);
\node[align=center] at (.53,-1.55) {C gains 30.6 percentage points of share.\\Category share grows by only about 11.0 points.};
\end{tikzpicture}

\caption{A new product’s sales include substitution from existing products.}\label{fig:M04_cannibalization}
\end{figure}

With utility coefficients and the potential-customer population fixed,
adding C changes only the available choice set. Because C has the same
modeled attributes as A, it is a close substitute. The logit model's
independence-of-irrelevant-alternatives property makes C draw share
proportionally from all existing alternatives. Whether that pattern is
plausible is an empirical question. Richer substitution models may be
needed when close substitutes compete especially strongly with one
another.

A profit calculation must attach contribution margins and fixed costs to
the simulated quantities. If A and C each contribute \$4 per purchase
and B contributes \$3, expected contribution per potential customer
rises from about \$2.36 to \$2.86 before product-development and
operating costs. The contribution comparison is conditional on the
specified utilities, margins, and fixed customer population. A pilot or
holdout choice tasks can help evaluate the proposed product and model
before a larger decision.

\hypertarget{M04-estimate-and-validate-preferences}{%
\section{Estimate and Validate
Preferences}\label{M04-estimate-and-validate-preferences}}

\noindent A conjoint workflow begins with the managerial choice and feasible
attributes, then constructs the task design and recruits a relevant
respondent sample. Inspect completion time, straight-line responses,
inconsistent attention, and whether participants understood the
attributes. Exclusion rules should be set and reported transparently.
Removing respondents because their preferences are inconvenient changes
the target population.

Estimate a transparent baseline before adopting a more elaborate model.
Check coefficient coding, reference levels, price units, and whether
probabilities sum to one for every choice set. Evaluate prediction on
held-out tasks or respondents according to the intended use. A model
that fits familiar respondents' repeated choices may generalize
differently to a new customer population.

A useful simulation changes one decision at a time and records what is
fixed: attributes, price, competitor options, market size, and margins.
Share of preference from a survey is not automatically market share.
Distribution, awareness, availability, purchase frequency, and sampling
can create differences. State the link between the conjoint output and
the demand quantity used in the business calculation.

\hypertarget{M04-segmentation-begins-with-a-definition-of-similarity}{%
\section{Segmentation Begins With a Definition of
Similarity}\label{M04-segmentation-begins-with-a-definition-of-similarity}}

\noindent A retailer may want segments for differentiated service, product
positioning, or campaign design. Those objectives should guide the
variables used to define similarity. Grouping customers by recent
spending answers a different question from grouping them by category
preferences or service needs. A segment should be interpretable,
sufficiently stable, and connected to an action the firm can take.

Recency, frequency, and monetary value provide a useful descriptive
starting point. Define recency as time since the last eligible purchase
at a common observation date, frequency as the number of eligible
purchases during a fixed window, and monetary value as sales or
contribution during that window. Customers with different observation
histories require adjustment or a deliberate eligibility rule. A new
customer observed for one month should not be compared mechanically with
one observed for a year.

K-means selects \(K\) centers to minimize within-cluster squared
Euclidean distance:

\[
\sum_{i=1}^n\left\|z_i-\mu_{c(i)}\right\|^2.
\]

Here \(z_i\) is the vector of transformed customer features, \(c(i)\) is
customer \(i\)'s assigned cluster, and \(\mu_{c(i)}\) is its cluster
center. The feature transformations determine the distance.
Standardization often helps prevent a large numerical scale from
dominating, but it does not establish that all features should have
equal managerial importance.

For three customers with spending-frequency pairs A=(100,1), B=(110,9),
and C=(200,1), raw Euclidean distance makes B much closer to A than C
is. If spending is expressed in hundreds of dollars while frequency
remains in purchases, C becomes closer to A. The customer records have
not changed; the implicit weights have. This is why units and scaling
belong in the substantive definition of a segment.

\begin{figure}[htbp]
\centering
% Distances within each panel use the stated feature units.
\begin{tikzpicture}
\begin{axis}[name=raw,width=6.3cm,height=6.1cm,xmin=.5,xmax=2.5,ymin=0,ymax=120,ybar,bar width=18pt,xtick={1,2},xticklabels={A to B,A to C},ylabel={Euclidean distance},title={Raw spending and visits},nodes near coords,nodes near coords style={font=\scriptsize}]
\addplot[fill=IlliniBlue,draw=IlliniBlue] coordinates {(1,12.806)(2,100)};
\end{axis}
\begin{axis}[at={(raw.east)},anchor=west,xshift=2.3cm,width=6.3cm,height=6.1cm,xmin=.5,xmax=2.5,ymin=0,ymax=10,ybar,bar width=18pt,xtick={1,2},xticklabels={A to B,A to C},ylabel={Euclidean distance},title={Spending divided by 100},nodes near coords,nodes near coords style={font=\scriptsize}]
\addplot[fill=IlliniOrange,draw=IlliniOrange] coordinates {(1,8.0006)(2,1)};
\end{axis}
\end{tikzpicture}

\caption{Changing feature scales changes which customers appear similar.}\label{fig:M04_scaling}
\end{figure}

Changing spending units reverses the nearest neighbor because it changes
the relative weight placed on spending and visits. A segmentation report
should explain that metric and examine whether the proposed actions
remain sensible under reasonable transformations.

\hypertarget{M04-choosing-and-evaluating-segments}{%
\section{Choosing and Evaluating
Segments}\label{M04-choosing-and-evaluating-segments}}

\noindent Use more than one criterion to choose \(K\). Within-cluster dispersion
falls mechanically as more clusters are added. A silhouette measure
compares cohesion with separation under the chosen distance, but a high
silhouette does not establish business value. Examine segment sizes,
profiles, stability across resamples or time, and whether the segments
imply different feasible actions.

Hierarchical clustering provides a nested set of groupings whose
interpretation depends on the linkage and distance. Density-based
methods can identify irregular groups and mark some observations as
noise, with sensitivity to neighborhood parameters. Mixture models
provide probabilistic membership under distributional assumptions. These
alternatives address different patterns; no algorithm discovers a
uniquely correct segmentation independent of representation and purpose.

For customer activation, avoid using future purchase outcomes to define
a segment that supposedly exists before the campaign. For evaluating
segmentation, examine behavior not used mechanically to construct it. A
high-spending segment will have high spending by design if spending is a
clustering input. A more useful question is whether the grouping helps
predict future needs or improves the result of a tested service policy.

Clustering shopping trips on basket size, category composition, and
spending can help separate trip types before comparing channels. Fit the
representation without allowing the desired channel comparison to
determine the grouping. Examine profiles and sensitivity, because a
trip-classification algorithm can also create misleading comparisons if
its features or thresholds differ across channels.

\hypertarget{M04-customer-lifetime-value-depends-on-timing-and-margin}{%
\section{Customer Lifetime Value Depends on Timing and
Margin}\label{M04-customer-lifetime-value-depends-on-timing-and-margin}}

\noindent Customer lifetime value is the discounted expected contribution from a
relationship, under a stated time horizon and behavioral model. For a
newly acquired customer, a simple finite-horizon expression is

\[
\operatorname{CLV}=\sum_{t=1}^{T}\frac{\Pr(\text{active at }t)\,m_t}{(1+d)^t}-a.
\]

Here \(t\) indexes periods after acquisition, \(m_t\) is expected
contribution conditional on being active in period \(t\), \(d\) is the
discount rate per period, and \(a\) is acquisition cost paid at time
zero. Contribution should deduct relevant variable service and
fulfillment costs. The expression assumes the conditional margin and
activity probability have been defined consistently; more general models
can use the expected product of activity and customer-specific margin.

Assume a stationary subscription process: a currently active customer
survives each period with probability \(r\), earns a constant
contribution \(m\) at the end of a period only if still active, and has
an infinite horizon. With \(r/(1+d)<1\), the geometric series gives

\[
\operatorname{CLV}=\frac{mr}{1+d-r}-a.
\]

This convention excludes an immediate current-period contribution.
Including one would add that contribution and change the formula.
Monthly retention requires monthly margins and a monthly discount rate.
Combining monthly retention with an annual discount rate produces an
inconsistent calculation.

With monthly contribution \(m=60\), retention \(r=0.80\), monthly
discount rate \(d=0.01\), and acquisition cost \(a=50\), CLV is
approximately \$178.57. Increasing retention to 0.90 raises it to
approximately \$440.91. The gain is nonlinear because retention affects
many future payments. These values are conditional on constant retention
and margin; changes in retention or margin over time would change the
valuation.

\begin{figure}[htbp]
\centering
% Original synthetic teaching figure; see corresponding module.
\begin{tikzpicture}
\begin{axis}[xmin=.5,xmax=.95,ymin=0,ymax=1000,xlabel={Per-period retention probability},ylabel={Lifetime value after acquisition cost (dollars)},xtick={.5,.6,.7,.8,.9,.95},ytick={0,200,400,600,800,1000},ymajorgrids=true]
\addplot[IlliniBlue,domain=.5:.95,samples=150,very thick] {60*x/(1.01-x)-50};
\addplot[IlliniOrange,only marks,mark=*,mark size=3pt] coordinates {(.8,178.5714)(.9,440.9091)};
\node[anchor=north west] at (axis cs:.8,170) {$r=.80$: \$178.57};
\node[anchor=south east] at (axis cs:.89,480) {$r=.90$: \$440.91};
\node[align=left,anchor=north west] at (axis cs:.52,920) {$m=60,\ d=.01,\ a=50$\\Payments begin after first retention};
\end{axis}
\end{tikzpicture}

\caption{The value of retention is nonlinear under a stationary subscription model.}\label{fig:M04_clv}
\end{figure}

With monthly contribution, discounting, and acquisition cost fixed, a
one-point retention increase has greater value at a higher baseline
retention rate. Uncertainty about long-run retention can therefore
dominate a CLV forecast.

\hypertarget{M04-subscription-and-noncontractual-relationships}{%
\section{Subscription and Noncontractual
Relationships}\label{M04-subscription-and-noncontractual-relationships}}

\noindent In a subscription setting, cancellation is often observed. Survival
analysis can model time to cancellation and account for right censoring
when observation ends before cancellation. A hazard describes the
probability of exit during a period conditional on survival to that
period. Cohort effects, tenure dependence, and changes in service
quality can make constant retention implausible.

In a noncontractual retail setting, the last purchase does not reveal
whether a customer has permanently left. Models such as BG/NBD represent
purchase frequency and latent dropout under specific stochastic
assumptions. A separate spending model may estimate value per
transaction. Those assumptions require checking, including the treatment
of returns and any assumed relationship between purchase frequency and
monetary value. A customer-level count alone does not identify permanent
churn without a model.

Evaluate value forecasts on a holdout period with the same horizon used
in the decision. Compare aggregate and customer-level calibration, rank
ordering, and errors for new or infrequent customers. A model can rank
customers well while overpredicting total value. For an acquisition
budget, the aggregate monetary error may be more consequential than a
small improvement in ranking.

\hypertarget{M04-integrated-case-which-segment-should-receive-retention-support}{%
\section{Integrated Case: Which Segment Should Receive Retention
Support?}\label{M04-integrated-case-which-segment-should-receive-retention-support}}

\noindent Suppose segment H has high predicted CLV but little response to service
outreach, while segment M has moderate predicted CLV and a larger
response. Targeting only H can be inferior even if its value forecast is
accurate. The relevant quantity is the change in expected future
contribution caused by outreach, net of the outreach cost. CLV describes
a baseline relationship; a retention intervention changes that
relationship under a separate causal model.

A useful analysis would estimate or specify segment-specific
preferences, margins, and retention, then test the proposed offer.
Product design might attract a segment with expensive service needs; a
seemingly attractive acquisition price might therefore have low
contribution. Segment labels should help organize the decision without
replacing these calculations. A final recommendation should show which
assumption changes the preferred segment or product.

For practice, recompute the subscription value with \(r=0.85\) and
verify the result before acquisition cost is subtracted. Add an
immediate \$60 contribution and explain the change in timing. In the
conjoint example, remove the outside option and observe that simulated
purchase incidence becomes one by construction. In the clustering
example, compare similarity using category shares instead of spending
levels and explain which managerial question the new representation
serves.

\hypertarget{M04-ai-assisted-implementation-and-readings}{%
\section{AI-Assisted Implementation and
Readings}\label{M04-ai-assisted-implementation-and-readings}}

\noindent Ask a coding assistant to calculate choice probabilities from a supplied
profile table, check that they sum to one, and separate a new product's
share from incremental category demand. For clustering, require
documented transformations, segment profiles, and a sensitivity
comparison. For CLV, request an explicit payment timeline and verify the
closed-form calculation against a long finite sum. These checks catch
common mistakes that can survive apparently reasonable charts.

\href{https://sawtoothsoftware.com/academics/teaching-resources}{Sawtooth's
academic materials} provide conjoint examples.
\href{https://link.springer.com/book/10.1007/978-981-10-8818-6}{Dolnicar,
Grün, and Leisch's segmentation book} develops the connection between
segmentation design and use.
\href{https://www.pymc-marketing.io/en/latest/notebooks/clv/clv_quickstart.html}{PyMC-Marketing's
CLV documentation} offers implementation examples whose model
assumptions and software versions should be checked before adaptation.

\hypertarget{M04-mhealth-behavioral-segments-and-product-economics}{%
\section{mHealth: Behavioral Segments and Product
Economics}\label{M04-mhealth-behavioral-segments-and-product-economics}}

\noindent A product manager may want different interfaces for occasional loggers,
frequent loggers with narrow recorded diets, and frequent loggers with
broad recorded diets. Segmenting on January through June behavior
creates groups whose later logging can be described without using later
observations to define the groups. A segment is useful when it
corresponds to a decision the manager can implement, such as simplifying
entry or offering a broader set of suggestions.

Eligibility requires at least 30 observed days before July 1 and
nonmissing average variety. Each user contributes three features:
observed days divided by the 182 calendar days in the first half of
2016, mean recorded category count, and mean calorie-share variety.
Standardization gives each feature unit variance before three-cluster
k-means. Three clusters are imposed for a manageable descriptive
comparison. Cluster labels are ordered by their mean prior logging
fraction.

\begin{table}[H]
\centering\small
\begin{tabular}{@{}lrrrr@{}}
\toprule
Segment & Users & Prior Logging Fraction & Mean Breadth & Later Logging
Fraction \\
\midrule

1 & 126 & 0.426 & 10.84 & 0.338 \\
2 & 90 & 0.759 & 6.99 & 0.594 \\
3 & 115 & 0.870 & 12.80 & 0.654 \\
\bottomrule
\end{tabular}
\end{table}

Later logging uses the 184 calendar days from July through December,
with zero observed days for eligible users absent from that period. The
differences describe persistence in recorded use. They are not treatment
responses, contract renewal rates, or purchase probabilities. In
particular, Segment 2 combines frequent logging with relatively narrow
breadth. A single ranking by activity would conceal that distinction.

Repeating k-means with a second seed produces an adjusted Rand index of
1.0 for these inputs. The partition is stable to this initialization
check. That result does not establish three natural customer types or
prove stability to other features, calendar windows, or measurement
screens. The calorie-based variety and breadth features are related, so
including both places additional emphasis on recorded dietary diversity.
Nutrition screening can also alter the segmentation inputs; compare its
consequences before naming a segment around a health interpretation.

\begin{figure}[htbp]
\centering
\begin{tikzpicture}
\begin{axis}[xlabel={Prior Fraction of Calendar Days Logged},ylabel={Mean Recorded Category Count},xmin=.25,xmax=1,ymin=5,ymax=15,grid=major]
\addplot[only marks,IlliniBlue,mark=*,mark size=5pt] coordinates {(.425999,10.8354)};
\addplot[only marks,IlliniOrange,mark=square*,mark size=5pt] coordinates {(.758608,6.989)};
\addplot[only marks,Patina,mark=triangle*,mark size=5pt] coordinates {(.869804,12.7997)};
\node[anchor=south west] at (axis cs:.44,11.1) {Segment 1: 126 Users};
\node[anchor=north east] at (axis cs:.75,6.7) {Segment 2: 90 Users};
\node[anchor=south east] at (axis cs:.88,13.2) {Segment 3: 115 Users};
\end{axis}
\end{tikzpicture}

\caption{Prior logging frequency and recorded breadth distinguish groups that a single activity ranking would combine. Both axes use January through June; later behavior does not define the partition.}\label{fig:M04_mhealth_segments}
\end{figure}

The food logs supply no observed subscription price, product-choice
experiment, margin, acquisition cost, or payment-based retention. A
hypothetical paid-plan calculation can nevertheless clarify the extra
inputs needed. Set the outside-option utility to zero, paid-plan benefit
to 0.8, price coefficient to -0.1 per dollar, and price to \$5. The
imposed logit model yields a paid-plan share of 0.5744. These inputs are
assumptions, not conjoint estimates from the logs.

For a separate hypothetical subscription, suppose monthly contribution
conditional on retention is \$6, the monthly discount rate is 1\%,
acquisition cost is \$10, and the first contribution occurs one month
after acquisition. With constant monthly retention \(r\), value is
\(6r/(1.01-r)-10\). Assumed retention of 0.80, 0.85, and 0.90 gives
values of \$12.86, \$21.87, and \$39.09. These scenarios are not
assigned to the three behavioral segments. Doing so would equate logging
persistence with paid retention without evidence.

The code fits the scaler and clusters using only the earlier period,
then joins later logging for description. An AI assistant should explain
each feature's units, how missing later records enter the denominator,
and which economic inputs remain hypothetical before proposing a
targeting rule.

\FloatBarrier
\Needspace{8\baselineskip}
Python source:
\href{../03_Data_and_Code/mHealth/recurring_examples/m04.py}{m04.py}.
From its directory, run \texttt{python\ m04.py}.

\begin{Verbatim}[breaklines,breakanywhere,fontsize=\small]
"""Segment January-June behavior and describe later logging."""

import json
import numpy as np
import pandas as pd
from sklearn.cluster import KMeans
from sklearn.preprocessing import StandardScaler
from sklearn.metrics import adjusted_rand_score
from common import OUTPUT, SEED, load_inputs, save_result

food, daily, baskets = load_inputs()
prior = daily.loc[daily.date.lt("2016-07-01")]
features = prior.groupby("Userid").agg(
    logged_days=("date", "size"),
    breadth=("category_count", "mean"),
    variety=("variety", "mean"),
)
features = features.loc[features.logged_days.ge(30)].dropna()
features["logging_fraction"] = features.logged_days / 182
columns = ["logging_fraction", "breadth", "variety"]
scaled = StandardScaler().fit_transform(features[columns])
model = KMeans(n_clusters=3, n_init=20, random_state=SEED).fit(scaled)
alternative = KMeans(n_clusters=3, n_init=20, random_state=SEED + 1).fit(scaled)
# Number clusters by their mean prior logging fraction for readable output.
order = (
    features.assign(raw=model.labels_)
    .groupby("raw")
    .logging_fraction.mean()
    .sort_values()
    .index
)
labels = {raw: rank + 1 for rank, raw in enumerate(order)}
features["segment"] = [labels[value] for value in model.labels_]
later = daily.loc[daily.date.ge("2016-07-01")].groupby("Userid").size()
features["later_logging_fraction"] = later.reindex(features.index, fill_value=0) / 184
table = features.groupby("segment").agg(
    users=("segment", "size"),
    prior_logging=("logging_fraction", "mean"),
    breadth=("breadth", "mean"),
    variety=("variety", "mean"),
    later_logging=("later_logging_fraction", "mean"),
)
OUTPUT.mkdir(exist_ok=True)
table.to_csv(OUTPUT / "M04_segments.csv")
# Imposed monthly financial assumptions, separate from the log estimates.
retention = np.array([0.80, 0.85, 0.90])
clv = 6 * retention / (1.01 - retention) - 10
utility = np.array([0.0, 0.8 - 0.1 * 5])  # Outside option, hypothetical paid plan.
choice = np.exp(utility - utility.max())
choice /= choice.sum()
result = {
    "users": len(features),
    "features": columns,
    "seed_agreement_ari": float(
        adjusted_rand_score(model.labels_, alternative.labels_)
    ),
    "segments": table.reset_index().to_dict(orient="records"),
    "hypothetical_retention": retention.tolist(),
    "hypothetical_clv": clv.tolist(),
    "hypothetical_paid_plan_share": float(choice[1]),
}
save_result("M04", result)
print(json.dumps(result, indent=2))
\end{Verbatim}

\hypertarget{M04-python-calculations-for-the-synthetic-cases}{%
\section{Python Calculations for the Synthetic
Cases}\label{M04-python-calculations-for-the-synthetic-cases}}

\noindent The inputs below are imposed teaching values. Running the program
creates the labeled synthetic CSV files and prints the corresponding
calculations. They are separate from the mHealth estimates. Keep the
script beside \texttt{example\_helpers.py} in
\texttt{03\_Data\_and\_Code/Lecture\_Examples}.

\FloatBarrier
\Needspace{8\baselineskip}
Python source:
\href{../03_Data_and_Code/Lecture_Examples/m04_synthetic.py}{m04\_synthetic.py}.
From its directory, run \texttt{python\ m04\_synthetic.py}.

\begin{Verbatim}[breaklines,breakanywhere,fontsize=\small]
"""Synthetic numerical examples for M4; run from any working directory."""

import itertools
import json
import math
from example_helpers import write_csv, ols_slope, shares

r = {}
before = shares({"A": 0.2, "B": -0.6, "outside": 0})
after = shares({"A": 0.2, "B": -0.6, "C": 0.2, "outside": 0})
r["M04"] = {
    "shares_before": before,
    "shares_after": after,
    "category_share_increase": before["outside"] - after["outside"],
    "contribution_before": 4 * before["A"] + 3 * before["B"],
    "contribution_after": 4 * (after["A"] + after["C"]) + 3 * after["B"],
    "wtp_dollars": 1.2 / 0.1,
    "distance_raw_AB": math.hypot(10, 8),
    "distance_raw_AC": 100,
    "distance_rescaled_AB": math.hypot(0.1, 8),
    "distance_rescaled_AC": 1,
    "clv_after_acquisition": {
        str(retention): 60 * retention / (1.01 - retention) - 50
        for retention in [0.8, 0.85, 0.9]
    },
}
write_csv(
    "M04_segmentation_synthetic.csv",
    ["customer", "spending_dollars", "visits"],
    [("A", 100, 1), ("B", 110, 9), ("C", 200, 1)],
)
assert math.isclose(sum(before.values()), 1) and math.isclose(sum(after.values()), 1)
print(json.dumps(r, indent=2))
\end{Verbatim}

\hypertarget{M04-mhealth-time-until-a-logging-gap}{%
\section{mHealth: Time Until a Logging
Gap}\label{M04-mhealth-time-until-a-logging-gap}}

\noindent A retention manager may want to intervene before a long interruption in
activity. The analytical object must specify when the interruption has
actually occurred. A user who is absent for two days and returns on the
third day has a different history from one who remains absent for
several weeks. The end of the data window can arrive before either
pattern is fully observed.

Define eligibility using at least 30 recorded days before July 1. Among
eligible users who return during July, set time zero to that user's
first July logging day. The event occurs on the 14th consecutive day
without a subsequent log, and observation ends on August 31. The event
date is the day the gap is confirmed. A later return does not erase a
previously confirmed gap. Customers without a confirmed gap by August 31
are administratively censored at their elapsed follow-up time.

This definition concerns an interruption in observed logging. Permanent
customer departure, failure to eat, and cancellation of a paid
subscription are separate events. July return is also a selection
condition: 331 users satisfy prior eligibility, and 297 have a July time
origin. The curve describes these returners, including a mixture of
early and late July start dates.

To allow different amounts of follow-up, define \(n_t\) as users still
under observation and gap-free immediately before elapsed day \(t\), and
\(d_t\) as confirmed first gaps on that day. The Kaplan--Meier estimate
is

\[
\widehat S(t)=\prod_{s\le t}\left(1-\frac{d_s}{n_s}\right).
\]

A censored user contributes to earlier risk sets and leaves subsequent
ones. Censoring itself causes no downward step. When an event and
administrative censoring share a day, the event is counted before
censoring. A population interpretation of the pooled curve requires
censoring to be independent of gap timing, possibly conditional on
relevant start-date characteristics. Calendar trends can undermine that
condition because later July starters have shorter follow-up.

\begin{table}[H]
\centering\small
\begin{tabular}{@{}
  >{\raggedright\arraybackslash}p{(\columnwidth - 10\tabcolsep) * \real{0.13}}
  >{\raggedleft\arraybackslash}p{(\columnwidth - 10\tabcolsep) * \real{0.17}}
  >{\raggedleft\arraybackslash}p{(\columnwidth - 10\tabcolsep) * \real{0.17}}
  >{\raggedleft\arraybackslash}p{(\columnwidth - 10\tabcolsep) * \real{0.17}}
  >{\raggedleft\arraybackslash}p{(\columnwidth - 10\tabcolsep) * \real{0.17}}
  >{\raggedleft\arraybackslash}p{(\columnwidth - 10\tabcolsep) * \real{0.19}}@{}}
\toprule
Confirmed Gap & Users & Events by August 31 & Censored & Gap-Free at Day
30 & Gap-Free at Day 60 \\
\midrule

14 Days & 297 & 93 & 204 & 0.8182 & 0.6867 \\
28 Days & 297 & 33 & 264 & 0.9630 & 0.8877 \\
\bottomrule
\end{tabular}
\end{table}

\begin{figure}[htbp]
\centering
\begin{tikzpicture}
\begin{axis}[xlabel={Days Since First July Log},ylabel={Estimated Gap-Free Fraction},xmin=0,xmax=63,ymin=.6,ymax=1.01,xtick={0,14,28,45,60},ytick={.6,.7,.8,.9,1},yticklabel style={/pgf/number format/fixed,/pgf/number format/precision=1},grid=major,legend style={at={(.02,.03)},anchor=south west}]
\addplot[IlliniOrange,solid,const plot,very thick] coordinates {(0,1) (14,0.95622896) (15,0.92592593) (17,0.91919192) (18,0.91582492) (19,0.91245791) (20,0.90909091) (21,0.90572391) (22,0.89562290) (23,0.88215488) (24,0.85858586) (25,0.85521886) (27,0.84511785) (28,0.83164983) (29,0.82491582) (30,0.81818182) (31,0.81481481) (32,0.81144781) (33,0.80808081) (34,0.80471380) (35,0.80134680) (36,0.79797980) (37,0.78787879) (38,0.77773443) (39,0.77435297) (40,0.76420861) (41,0.76420861) (42,0.76079696) (44,0.75397367) (45,0.73675966) (46,0.73330070) (47,0.72638277) (48,0.71946483) (49,0.71946483) (50,0.71251348) (51,0.70556213) (52,0.70556213) (53,0.70196233) (54,0.69832522) (55,0.69101291) (56,0.69101291) (57,0.69101291) (58,0.69101291) (59,0.69101291) (60,0.68669408) (61,0.68211612)};
\addlegendentry{14-Day Gap}
\addplot[IlliniOrange,only marks,mark=+,mark size=2.4pt,forget plot] coordinates {(37,0.78787879) (40,0.76420861) (41,0.76420861) (44,0.75397367) (45,0.73675966) (49,0.71946483) (51,0.70556213) (52,0.70556213) (53,0.70196233) (54,0.69832522) (55,0.69101291) (56,0.69101291) (57,0.69101291) (58,0.69101291) (59,0.69101291) (60,0.68669408) (61,0.68211612)};
\addplot[IlliniBlue,dashed,const plot,very thick] coordinates {(0,1) (28,0.97643098) (29,0.96296296) (31,0.95959596) (32,0.95622896) (33,0.95286195) (35,0.95286195) (36,0.94948301) (37,0.94610407) (38,0.93593090) (40,0.93593090) (41,0.93252752) (42,0.92911167) (44,0.92569582) (45,0.92569582) (46,0.92222879) (47,0.91876177) (48,0.91529474) (49,0.91529474) (50,0.91529474) (51,0.91176078) (52,0.90440787) (53,0.90065514) (54,0.89687088) (55,0.89687088) (56,0.89687088) (57,0.89687088) (58,0.89687088) (59,0.88771913) (60,0.88771913) (61,0.88264645)};
\addlegendentry{28-Day Gap}
\addplot[IlliniBlue,only marks,mark=+,mark size=2.4pt,forget plot] coordinates {(35,0.95286195) (37,0.94610407) (40,0.93593090) (41,0.93252752) (44,0.92569582) (45,0.92569582) (48,0.91529474) (49,0.91529474) (50,0.91529474) (51,0.91176078) (52,0.90440787) (53,0.90065514) (54,0.89687088) (55,0.89687088) (56,0.89687088) (57,0.89687088) (58,0.89687088) (59,0.88771913) (60,0.88771913) (61,0.88264645)};
\node[anchor=north east,font=\footnotesize] at (axis cs:62,1.005) {+ Administrative Censoring};
\end{axis}\end{tikzpicture}

\caption{Gap-free estimates depend on the interruption threshold; users censored at the data boundary leave the risk set without becoming events.}\label{fig:M04_logging_survival}
\end{figure}

The 28-day definition produces fewer confirmed events because an
interruption must persist longer. It also confirms an event later even
when the same interruption satisfies both thresholds. At elapsed day 60,
the risk sets contain 160 and 187 users for the 14- and 28-day
definitions, respectively. The table therefore requires an event
definition and a time origin to support a retention comparison.
Converting either curve into customer lifetime value would additionally
require a commercial relationship, margins, discounting, and a model
linking logging interruptions to that relationship.

The
\href{https://github.com/CamDavidsonPilon/lifelines/blob/master/examples/Customer\%20Churn.ipynb}{lifelines
customer-churn example} uses explicit event and duration fields. The
\href{https://lifelines.readthedocs.io/en/latest/Survival\%20analysis\%20with\%20lifelines.html}{Kaplan--Meier
documentation} supplies the corresponding risk-set estimator. The code
below computes each risk set explicitly so a small hand-constructed
event/censoring table can verify the calculation.

\FloatBarrier
\Needspace{8\baselineskip}
Python source:
\href{../03_Data_and_Code/mHealth/recurring_examples/logging_gaps.py}{logging\_gaps.py}.
From its directory, run \texttt{python\ logging\_gaps.py}.

\begin{Verbatim}[breaklines,breakanywhere,fontsize=\small]
"""Kaplan-Meier estimates for a defined logging-gap event, with censoring."""

import json
import pandas as pd
from common import OUTPUT, load_food, make_daily, save_result


def first_gap(dates, start, end, gap_days):
    dates = set(dates)
    run = 0
    for day in pd.date_range(start + pd.Timedelta(days=1), end):
        run = 0 if day in dates else run + 1
        if run == gap_days:
            return (day - start).days, True
    return (end - start).days, False


def kaplan_meier(records):
    # At a tied time, events precede administrative censoring.
    survival = 1.0
    rows = []
    for time in sorted(records.duration.unique()):
        risk = int(records.duration.ge(time).sum())
        events = int((records.duration.eq(time) & records.event).sum())
        censored = int((records.duration.eq(time) & ~records.event).sum())
        survival *= 1 - events / risk
        rows.append(
            {
                "day": int(time),
                "at_risk": risk,
                "events": events,
                "censored": censored,
                "gap_free": survival,
            }
        )
    return pd.DataFrame(rows)


if __name__ == "__main__":
    daily = make_daily(load_food())
    prior = daily.loc[daily.date < "2016-07-01"].groupby("Userid").size()
    eligible = prior.index[prior >= 30]
    july = daily.loc[
        daily.Userid.isin(eligible) & daily.date.between("2016-07-01", "2016-07-31")
    ]
    starts = july.groupby("Userid").date.min()
    observed = daily.groupby("Userid").date.agg(set)
    end = pd.Timestamp("2016-08-31")
    result = {
        "eligible_before_july": len(eligible),
        "july_returners": len(starts),
        "administrative_end": str(end.date()),
        "thresholds": [],
    }
    for gap in [14, 28]:
        records = pd.DataFrame(
            [
                first_gap(observed[user], start, end, gap)
                for user, start in starts.items()
            ],
            columns=["duration", "event"],
        )
        curve = kaplan_meier(records)
        summary = {
            "gap_days": gap,
            "users": len(records),
            "events": int(records.event.sum()),
            "censored": int((~records.event).sum()),
        }
        for horizon in [30, 45, 60]:
            before = curve.loc[curve.day <= horizon]
            summary[f"gap_free_day_{horizon}"] = (
                float(before.gap_free.iloc[-1]) if len(before) else 1.0
            )
            summary[f"at_risk_day_{horizon}"] = int(records.duration.ge(horizon).sum())
        result["thresholds"].append(summary)
        OUTPUT.mkdir(exist_ok=True)
        curve.to_csv(OUTPUT / f"logging_gap_{gap}.csv", index=False)
    save_result("logging_gaps", result)
    print(json.dumps(result, indent=2))
\end{Verbatim}

\hypertarget{M05-m5.-text-social-media-and-customer-networks}{%
\chapter{Text, Social Media, and Customer
Networks}\label{M05-m5.-text-social-media-and-customer-networks}}

\noindent A retailer receives thousands of reviews and support messages each week.
The service manager needs to distinguish delivery delays, product
defects, and billing problems, then decide which issues deserve
attention. A count of negative words is an incomplete representation:
``the refund was fast, but the replacement failed'' contains different
evaluations of two aspects of the experience. The analytical task is to
turn language into a defensible measurement of a concept relevant to
action.

Text analysis also raises a sampling problem. Customers who post reviews
can differ from those who remain silent, and the platform can affect
what gets posted or displayed. Social-network analysis adds another
layer: observed connections and similar behavior can reflect influence,
common circumstances, or selection into relationships.

\hypertarget{M05-define-the-construct-and-the-document}{%
\section{Define the Construct and the
Document}\label{M05-define-the-construct-and-the-document}}

\noindent A construct is the concept the analysis intends to measure, such as
delivery dissatisfaction, perceived authenticity, or purchase intent. A
label is the operational category assigned to an observation. A document
might be a complete review, one sentence, a support ticket, or a
conversation. Those units should follow the construct. A review-level
sentiment label can obscure a positive service experience and a negative
product experience in the same review.

Write an annotation codebook before producing labels at scale. It should
define each category, inclusion and exclusion rules, treatment of
ambiguity, and examples near category boundaries. If a message can
contain several problems, a multi-label task may be appropriate. If the
decision requires routing each ticket to one team, define a primary
issue or a routing priority separately from the presence of every issue.

Metadata can be informative and potentially misleading. Ratings may help
validate a sentiment measure but can also leak the intended label into a
prediction exercise. Product category, language, channel, and date can
reveal domain differences. Remove or protect personal information before
sending text to an external service, and use approved data access. A
scraped review is not automatically free of personal or contractual
restrictions.

\hypertarget{M05-representing-language-with-counts-and-weights}{%
\section{Representing Language With Counts and
Weights}\label{M05-representing-language-with-counts-and-weights}}

\noindent A bag-of-words representation counts tokens while largely ignoring their
order. An \(n\)-gram representation retains local sequences, so ``not
good'' can differ from ``good.'' Tokenization, case normalization,
punctuation, and stop-word rules change the resulting features. Removing
negation words can destroy information needed for sentiment. Removing
product names may help focus on attributes or may erase the very object
whose complaints the manager needs to identify.

Term frequency-inverse document frequency gives greater weight to terms
that are frequent in a document but less common across the corpus. One
simple convention is

\[
\operatorname{tfidf}_{wd}=\operatorname{tf}_{wd}\log\left(\frac{N}{\operatorname{df}_w}\right),
\]

where \(\operatorname{tf}_{wd}\) is the frequency of word \(w\) in
document \(d\), \(N\) is the number of documents, and
\(\operatorname{df}_w\) is the number of documents containing the word.
Implementations vary in smoothing, normalization, and term-frequency
definitions. Record the convention rather than assuming a library's
output matches this exact expression.

If ``delivery'' appears in two of three documents, its
inverse-document-frequency term under this convention is
\(\log(3/2)\approx0.405\). A word appearing in every document receives
zero weight. That weighting can help discriminate documents, but a
common term may still be substantively important. Statistical
distinctiveness and managerial importance are different criteria.

Fit a learned vocabulary and IDF weights on training text when
evaluating prediction. Learning them from the test corpus exposes
information about the evaluation distribution. The leakage may appear
minor compared with using labels, yet it still changes the task. A
pipeline makes the training-only transformation explicit and
reproducible.

\hypertarget{M05-embeddings-and-semantic-similarity}{%
\section{Embeddings and Semantic
Similarity}\label{M05-embeddings-and-semantic-similarity}}

\noindent An embedding maps a document or token into a vector learned from
language data. Similar vectors can capture related meanings even when
documents use different words. Cosine similarity compares the angle
between vectors:

\[
\operatorname{cos}(u,v)=\frac{u^\top v}{\|u\|\,\|v\|}.
\]

The normalization removes vector length from the comparison. This is
useful when direction represents meaning in the selected model, but it
does not guarantee that the similarity corresponds to the marketing
construct. Two messages about refunds can be semantically close while
expressing opposite evaluations. Retrieval quality should be checked
with examples that matter for the downstream task.

A simple TF-IDF classifier can be a strong baseline for a well-defined
domain. Embeddings may help with paraphrases, sparse labels, and
semantic retrieval, but they introduce dependence on a pretrained model
and its training domain. Compare alternatives on the same held-out data
and operational criteria. Record model versions because an updated
representation can change nearest neighbors and labels even when the
input text is unchanged.

For interpretation, inspect representative successes and failures. A
model may associate a brand name with complaints because one product had
a recall during training. It can then misclassify neutral mentions after
the issue is resolved. Temporal validation and product-level holdouts
help reveal whether performance depends on such shortcuts.

\hypertarget{M05-sentiment-aspect-labels-and-topic-models}{%
\section{Sentiment, Aspect Labels, and Topic
Models}\label{M05-sentiment-aspect-labels-and-topic-models}}

\noindent A sentiment lexicon assigns scores to words and combines them using
rules. It offers a transparent baseline, but sarcasm, domain-specific
meanings, and context can cause errors. A supervised classifier learns a
mapping from labeled examples to outcomes. Its labels and training
sample define what it learns; a generic ``positive/negative'' dataset
may not represent customer-service severity.

Aspect-based analysis separates the object being evaluated from the
evaluation itself. For the sentence ``delivery was late, but the support
agent solved the problem,'' a useful record might contain
delivery-negative and service-positive labels. A single overall score
discards that distinction. If the action is a logistics intervention,
the aspect-specific measure is closer to the decision.

Topic models describe recurring patterns of word co-occurrence or
document similarity. A topic is an estimated pattern, not automatically
a stable consumer need. In probabilistic models, documents can combine
topics; in clustering approaches, assignments may be discrete or
probabilistic. Analysts supply topic names after inspecting high-weight
words and representative documents. Validate those names with actual
text and examine stability across preprocessing and model choices.

A manager should avoid interpreting an increase in a topic's share
without checking its denominator. More delivery complaints can mean more
delivery problems, more deliveries overall, a changed review prompt, or
fewer reviews about other topics. Combining text measures with
transaction or service-volume data helps distinguish these explanations.

\hypertarget{M05-case-study-a-complaint-detector-under-changing-prevalence}{%
\section{Case Study: A Complaint Detector Under Changing
Prevalence}\label{M05-case-study-a-complaint-detector-under-changing-prevalence}}

\noindent Suppose a classifier identifies a rare but costly product defect. In a
synthetic evaluation population of 1,000 tickets, 200 truly mention the
defect. The classifier detects 160 of them and incorrectly flags 80 of
the 800 other tickets. Sensitivity, or recall, is 80\%; specificity is
90\%; precision is 160/240, or 66.7\%.

\begin{table}[H]
\centering\small
\begin{tabular}{@{}lrrrr@{}}
\toprule
Population & True Defect Tickets & True Positives & False Positives &
Precision \\
\midrule

Evaluation sample & 200 & 160 & 80 & 66.7\% \\
Deployment sample & 10 & 8 & 99 & 7.5\% \\
\bottomrule
\end{tabular}
\end{table}

The second row holds sensitivity at 80\% and specificity at 90\%, but
reduces true prevalence to 1\%. Among 990 nondefect tickets, 99 are
flagged. Only eight of the 107 flagged tickets are true positives. A
manager using the original precision to plan review staffing would be
badly misled.

Bayes' rule makes the dependence explicit. If prevalence is \(\rho\),
sensitivity is \(s\), and specificity is \(q\), then

\[
\Pr(\text{defect}\mid\text{flag})=\frac{s\rho}{s\rho+(1-q)(1-\rho)}.
\]

The denominator counts both correct and incorrect flags. For a rare
issue, even a modest false-positive rate can dominate that total. This
is why validation should represent the deployment population or
explicitly adjust for a changed sampling design.

\begin{figure}[htbp]
\centering
% Original synthetic teaching figure; see corresponding module.
\begin{tikzpicture}
\begin{axis}[xmin=0,xmax=.3,ymin=0,ymax=1,xlabel={Prevalence of the target complaint},ylabel={Precision among flagged reviews},xtick={0,.1,.2,.3},ytick={0,.2,.4,.6,.8,1},ymajorgrids=true]
\addplot[IlliniBlue,very thick,domain=0:.3,samples=150] {.8*x/(.8*x+.1*(1-x))};
\addplot[IlliniOrange,only marks,mark=*,mark size=3pt] coordinates {(.01,.0747664)(.2,.6666667)};
\node[align=left,anchor=north west] at (axis cs:.12,.96) {Sensitivity fixed at 80\%\\Specificity fixed at 90\%};
\node[anchor=west] at (axis cs:.018,.13) {1\% prevalence: 7.5\% precision};
\node[anchor=north west] at (axis cs:.145,.58) {20\% prevalence: 66.7\%};
\end{axis}
\end{tikzpicture}

\caption{Precision falls as the target issue becomes rarer, even with unchanged sensitivity and specificity.}\label{fig:M05_prevalence}
\end{figure}

Sensitivity and specificity remain at 0.80 and 0.90 in both populations.
The decline in precision comes entirely from the lower prevalence.
Operational evaluation should also consider the cost of missed defects,
the burden of false alarms, and whether uncertain cases can be routed to
human review.

\hypertarget{M05-human-and-llm-annotation}{%
\section{Human and LLM Annotation}\label{M05-human-and-llm-annotation}}

\noindent LLMs can help draft a codebook, suggest candidate labels, or classify
text using explicit instructions. For a defensible measurement exercise,
preserve the original text, the codebook version, model identifier,
prompt, output schema, and date. Require a bounded output format with an
allowed ``uncertain'' category where appropriate. Validate the output
against the schema and inspect cases that violate the instructions.

A prompt should identify the construct, unit, categories, boundary
rules, and examples. For delivery dissatisfaction, it should distinguish
a complaint about an actual delay from a question about the promised
delivery date. Ask for a short supporting span from the supplied text
when that improves auditability. A fluent explanation is not evidence
that the label is correct; compare it with independent human judgments
on a relevant sample.

Human agreement is informative because the construct may itself be
ambiguous. Double-code a sample, examine disagreements, and refine
definitions before scaling. Report category-specific errors, prevalence,
and performance across languages or product groups. An average accuracy
can conceal systematic failure on the minority category that motivates
the analysis.

Treat text being classified as data. A review may contain language
resembling an instruction to the model. It should not change the
annotation task, request external actions, or supply new system rules.
Separate the fixed task instructions from the document content,
constrain outputs, and inspect unexpected behavior. Use synthetic or
approved text for classroom exercises involving external services.

\hypertarget{M05-networks-connections-reach-and-brokerage}{%
\section{Networks: Connections, Reach, and
Brokerage}\label{M05-networks-connections-reach-and-brokerage}}

\noindent A network represents entities as nodes and relationships as edges. Nodes
might be customers, creators, brands, or products. An edge might
indicate following, communication, co-purchase, or a similarity
criterion. Directed and weighted networks preserve different information
from undirected binary networks. The edge definition should follow the
substantive relationship the analyst intends to study.

Degree counts a node's direct connections. Betweenness summarizes how
often a node lies on shortest paths between other nodes. A person with
few direct contacts can have high betweenness if those contacts connect
otherwise separate groups. Eigenvector-based measures give more weight
to connections with well-connected nodes. These metrics describe network
structure under their definitions; they do not by themselves establish
persuasion, credibility, or the effect of paying a creator.

\begin{figure}[htbp]
\centering
% Original synthetic teaching figure; see corresponding module.
\begin{tikzpicture}[v/.style={circle,draw=IlliniBlue,fill=white,minimum size=9mm},b/.style={circle,draw=IlliniOrange,very thick,fill=white,minimum size=10mm}]
\node[v] (a) at (0,1.2) {A};\node[v] (b) at (0,-1.2) {B};\node[v] (c) at (1.5,0) {C};\node[v] (d) at (-1.5,0) {D};
\node[b] (r) at (4,0) {R};
\node[v] (e) at (6.5,0) {E};\node[v] (f) at (8,1.2) {F};\node[v] (g) at (8,-1.2) {G};\node[v] (h) at (9.5,0) {H};
\foreach \x/\y in {a/b,a/c,a/d,b/c,b/d,c/d,e/f,e/g,e/h,f/g,f/h,g/h} {\draw[IlliniBlue] (\x)--(\y);}
\draw[IlliniOrange,very thick] (c)--(r)--(e);
\node[align=center] at (0,-2.2) {Four connected customers};
\node[align=center] at (8,-2.2) {Four connected customers};
\node[align=center,text width=4cm] at (4,2.3) {R has only two links,\\but connects the groups.};
\draw[accent] (4,1.6)--(r.north);
\end{tikzpicture}

\caption{A broker with few direct links can connect two otherwise separate customer groups.}\label{fig:M05_network_bridge}
\end{figure}

Broker R has two direct connections, while each within-group connector
has four. Removing R disconnects the two groups. A customer with fewer
direct contacts can therefore be essential for reaching a different
audience. A marketing campaign may value either role, depending on
message transmission, audience overlap, and the intended outcome.

Observed similarity among connected customers can arise from influence,
homophily, or common exposure. Friends may adopt a product because they
affect one another, because they already share tastes, or because they
see the same promotion. Regressing adoption on friends' adoption does
not distinguish those mechanisms without additional design. Randomized
encouragement, timing restrictions, or credible natural variation can
help, with attention to spillovers and the unit of assignment.

\hypertarget{M05-social-media-case-engagement-and-commercial-value}{%
\section{Social-Media Case: Engagement and Commercial
Value}\label{M05-social-media-case-engagement-and-commercial-value}}

\noindent A brand compares two creators. One receives many comments, but those
comments largely come from an audience that already follows the brand.
The other reaches a smaller group with less overlap. A high engagement
count does not establish more incremental sales. The comparison requires
a target population, exposure definition, outcome horizon, and a way to
measure what would have happened without the collaboration.

Text can help describe comment content, and network overlap can describe
audience duplication. Those measurements support campaign design, but
they should be combined with an evaluation strategy. A test might
randomize eligible regions or audience groups when feasible, or compare
outcomes using a credible observational design. If communication crosses
groups, spillovers can complicate the interpretation of an untreated
group.

For classroom practice, construct two small audience sets and calculate
their overlap. Then compare comments per exposed user with total
comments. Explain why neither is a direct estimate of incremental
profit.

\hypertarget{M05-ai-assisted-workflow-practice-and-readings}{%
\section{AI-Assisted Workflow, Practice, and
Readings}\label{M05-ai-assisted-workflow-practice-and-readings}}

\noindent Start a text exercise with a small manually inspected sample and a
codebook. Build a transparent baseline, reserve an evaluation sample,
and compare a more complex method on the same outcomes. Ask an assistant
to produce a confusion matrix and representative errors, then
independently recompute precision and recall from the counts. Keep model
selection and prompt refinement away from the final test sample.

For practice, calculate accuracy in the first complaint-detector sample:
88\%. Explain why a classifier that always predicts ``no defect'' would
achieve 99\% accuracy in the rare-prevalence deployment sample while
detecting no defects. Next identify how removing ``not'' changes the
meaning of ``not reliable.'' Finally, explain why two highly connected
adopters may share preferences even if neither caused the other's
adoption.

The
\href{https://scikit-learn.org/stable/modules/feature_extraction.html\#text-feature-extraction}{scikit-learn
text-feature documentation} supports count and TF-IDF representations.
\href{https://github.com/cjhutto/vaderSentiment}{VADER's author
repository} provides a rule-based sentiment baseline.
\href{https://doi.org/10.1007/s11002-020-09535-7}{Hagen and colleagues
(2020)} discuss ways machine learning can support behavioral marketing
research.

\hypertarget{M05-mhealth-classifying-new-food-descriptions}{%
\section{mHealth: Classifying New Food
Descriptions}\label{M05-mhealth-classifying-new-food-descriptions}}

\noindent The manager wants to reduce the effort required to categorize a food
entry. A lookup table may work well for familiar descriptions, while a
text classifier may help with a new description. These are distinct
deployment conditions. Repeated descriptions in both training and test
records can make a random entry split look successful even when the
system will struggle with unseen language.

Normalize descriptions by lowercasing, collapsing whitespace, and
trimming the ends. Retain nonempty descriptions and nonmissing source
categories, then sample 200,000 eligible entries with seed 534. The
sample contains 37,570 distinct normalized descriptions; 1,725 appear
with more than one source category. Those inconsistencies create label
ambiguity. The source category is the prediction target, not a validated
nutritional interpretation of the food.

The first split assigns approximately 20\% of entries to testing. The
second uses a fixed SHA-256 rule to hold out entire normalized
descriptions. Its test entries therefore share no exact normalized
description with training. Both evaluations weight food entries equally,
so frequent descriptions receive more weight. A second evaluation
weighting unique descriptions equally would answer a different product
question.

A training-only lookup chooses the most frequent category for each
description, breaking count ties alphabetically and falling back to the
most frequent training category. The text alternative fits a TF-IDF
vocabulary of unigrams and bigrams on distinct training descriptions,
caps the vocabulary at 20,000 features, and requires a term to occur in
at least two training descriptions. Multinomial naive Bayes uses
smoothing of 0.1 and training description-category counts as weights.

\begin{table}[H]
\centering\small
\begin{tabular}{@{}
  >{\raggedright\arraybackslash}p{(\columnwidth - 10\tabcolsep) * \real{0.13}}
  >{\raggedleft\arraybackslash}p{(\columnwidth - 10\tabcolsep) * \real{0.17}}
  >{\raggedleft\arraybackslash}p{(\columnwidth - 10\tabcolsep) * \real{0.17}}
  >{\raggedleft\arraybackslash}p{(\columnwidth - 10\tabcolsep) * \real{0.17}}
  >{\raggedleft\arraybackslash}p{(\columnwidth - 10\tabcolsep) * \real{0.17}}
  >{\raggedleft\arraybackslash}p{(\columnwidth - 10\tabcolsep) * \real{0.19}}@{}}
\toprule
Evaluation & Test Entries & Familiar Descriptions & Lookup Accuracy &
Text Accuracy & Text Macro F1 \\
\midrule

Entry Split & 40,273 & 88.65\% & 84.90\% & 84.94\% & 0.5531 \\
Description Split & 36,987 & 0\% & 3.82\% & 67.42\% & 0.3465 \\
\bottomrule
\end{tabular}
\end{table}

Macro F1 averages across the 312 labels present in the combined sampled
data, assigning zero to undefined class scores. Rare labels therefore
matter more than in accuracy. The gap between the two evaluations
reflects both familiarity and the changed test composition; it is not a
causal estimate of leakage's effect. For a service processing new
descriptions, the description-level evaluation is the closer test of
that requirement.

\begin{figure}[htbp]
\centering
\begin{tikzpicture}
\begin{axis}[ybar,bar width=21pt,ymin=0,ymax=1,ylabel={Entry-Weighted Accuracy},symbolic x coords={Entry Split,Description Split},xtick=data,enlarge x limits=.55,legend style={at={(.5,-.19)},anchor=north,legend columns=2},ymajorgrids=true]
\addplot[fill=IlliniBlue,draw=IlliniBlue] coordinates {(Entry Split,.848956)(Description Split,.038176)};
\addplot[fill=IlliniOrange,draw=IlliniOrange,postaction={pattern=north east lines,pattern color=white}] coordinates {(Entry Split,.849378)(Description Split,.674183)};
\legend{Description Lookup,TF-IDF + Naive Bayes}
\end{axis}
\end{tikzpicture}

\caption{Exact-description familiarity changes the task measured by a text-classification test.}\label{fig:M05_mhealth_text}
\end{figure}

The lookup nearly matches the text classifier on the entry split because
most test descriptions already appear in training. Its collapse on
unseen descriptions reveals how much of that performance depends on
memory. A product could use the lookup for verified familiar terms and
route uncertain new terms for review. Nutrition-based food composition
matching would additionally require portion units, preparation details,
and validated food codes. A category prediction alone supplies none of
those quantities.

An AI assistant can help inspect ambiguous examples and propose a
labeling rubric. Use only permitted local processing, preserve the
original category, and evaluate proposed relabeling against
independently reviewed examples. The code below makes no external model
request and learns every lookup and text transformation from training
rows only.

\FloatBarrier
\Needspace{8\baselineskip}
Python source:
\href{../03_Data_and_Code/mHealth/recurring_examples/m05.py}{m05.py}.
From its directory, run \texttt{python\ m05.py}.

\begin{Verbatim}[breaklines,breakanywhere,fontsize=\small]
"""Contrast familiar-description and unseen-description classification."""

import json
import numpy as np
from sklearn.feature_extraction.text import TfidfVectorizer
from sklearn.naive_bayes import MultinomialNB
from sklearn.metrics import accuracy_score, f1_score
from common import SEED, description_holdout, load_food, save_result

food = load_food()
usable = food.loc[food.description.ne("") & food.category.ne("Unclassified")]
sample = usable.sample(n=min(200000, len(usable)), random_state=SEED).reset_index(
    drop=True
)
rng = np.random.default_rng(SEED)
masks = {
    "entry_split": rng.random(len(sample)) < 0.2,
    "description_split": sample.description.map(description_holdout).to_numpy(),
}
results = {}
for name, mask in masks.items():
    train, test = sample.loc[~mask], sample.loc[mask]
    counts = train.groupby(["description", "category"]).size().reset_index(name="count")
    # Break tied category counts alphabetically; learn the lookup only on training rows.
    lookup = counts.sort_values(["count", "category"], ascending=[False, True])
    lookup = lookup.drop_duplicates("description").set_index("description").category
    majority = train.category.value_counts().sort_index().idxmax()
    baseline = test.description.map(lookup).fillna(majority)
    vectorizer = TfidfVectorizer(ngram_range=(1, 2), max_features=20000, min_df=2)
    vectorizer.fit(train.description.drop_duplicates())
    classifier = MultinomialNB(alpha=0.1)
    classifier.fit(
        vectorizer.transform(counts.description),
        counts.category,
        sample_weight=counts["count"],
    )
    prediction = classifier.predict(vectorizer.transform(test.description))
    labels = sorted(set(train.category) | set(test.category))
    results[name] = {
        "train_rows": len(train),
        "test_rows": len(test),
        "familiar_description_fraction": float(
            test.description.isin(lookup.index).mean()
        ),
        "lookup_accuracy": float(accuracy_score(test.category, baseline)),
        "nb_accuracy": float(accuracy_score(test.category, prediction)),
        "nb_macro_f1": float(
            f1_score(
                test.category,
                prediction,
                labels=labels,
                average="macro",
                zero_division=0,
            )
        ),
        "test_label_unseen_in_train_fraction": float(
            (~test.category.isin(train.category.unique())).mean()
        ),
        "metric_classes": len(labels),
    }
ambiguity = sample.groupby("description").category.nunique()
result = {
    "sample_rows": len(sample),
    "descriptions": len(ambiguity),
    "descriptions_with_multiple_labels": int(ambiguity.gt(1).sum()),
    "splits": results,
}
save_result("M05", result)
print(json.dumps(result, indent=2))
\end{Verbatim}

\hypertarget{M05-python-calculations-for-the-synthetic-cases}{%
\section{Python Calculations for the Synthetic
Cases}\label{M05-python-calculations-for-the-synthetic-cases}}

\noindent The inputs below are imposed teaching values. Running the program
creates the labeled synthetic CSV files and prints the corresponding
calculations. They are separate from the mHealth estimates. Keep the
script beside \texttt{example\_helpers.py} in
\texttt{03\_Data\_and\_Code/Lecture\_Examples}.

\FloatBarrier
\Needspace{8\baselineskip}
Python source:
\href{../03_Data_and_Code/Lecture_Examples/m05_synthetic.py}{m05\_synthetic.py}.
From its directory, run \texttt{python\ m05\_synthetic.py}.

\begin{Verbatim}[breaklines,breakanywhere,fontsize=\small]
"""Synthetic numerical examples for M5; run from any working directory."""

import itertools
import json
import math
from example_helpers import write_csv, ols_slope, shares

r = {}
r["M05"] = {}
for positive in [200, 10]:
    tp = 0.8 * positive
    fn = positive - tp
    fp = 0.1 * (1000 - positive)
    tn = 1000 - positive - fp
    r["M05"][str(positive / 1000)] = {
        "tp": tp,
        "fp": fp,
        "fn": fn,
        "tn": tn,
        "precision": tp / (tp + fp),
        "accuracy": (tp + tn) / 1000,
    }
write_csv(
    "M05_annotation_synthetic.csv",
    [
        "prevalence",
        "true_positive",
        "false_positive",
        "false_negative",
        "true_negative",
    ],
    [(p, v["tp"], v["fp"], v["fn"], v["tn"]) for p, v in r["M05"].items()],
)
print(json.dumps(r, indent=2))
\end{Verbatim}

\hypertarget{M05-mhealth-text-topics-and-their-denominators}{%
\section{mHealth: Text Topics and Their
Denominators}\label{M05-mhealth-text-topics-and-their-denominators}}

\noindent A product manager wants to organize a large food-description vocabulary
into interpretable groups. Repeated entries make a frequently logged
description prominent, while a catalog of distinct descriptions
emphasizes the range of language available. These populations answer
different questions. A topic model fitted to distinct descriptions can
summarize vocabulary, but its topic shares need a separate choice of
observation unit.

Use all 50,860 distinct nonempty normalized descriptions recorded before
July 1. Fit TF-IDF on these descriptions, retaining at most 3,000
unigram and bigram features, removing English stop words, and requiring
a feature in at least three documents. The maximum document-frequency
threshold is 0.95. Fix the number of nonnegative matrix factorization
components at four as an illustrative specification. No July description
is used to learn the vocabulary, inverse-document frequencies, or
component vectors.

Let \(V\) contain nonnegative TF-IDF features, with one row per training
description. NMF approximates \(V\) by \(WH\), where the rows of \(H\)
contain term patterns and each row of \(W\) gives the associated
nonnegative loadings. The chosen specification minimizes squared
reconstruction error subject to nonnegative factors. Multiplying a row
of \(H\) by a constant and dividing its column of \(W\) by the same
constant preserves the approximation. Before comparing loadings, the
code therefore scales them to correspond to component term vectors with
unit Euclidean norm.

July contains 17,042 distinct nonempty descriptions and 79,957 entries,
including 5,208 descriptions absent from training. Apply the frozen
vectorizer and factorization to July descriptions. Assign each
description to its largest scaled loading when positive. A zero-vector
description remains unassigned. The resulting groups have the following
leading terms and shares.

\begin{table}[H]
\centering\small
\begin{tabular}{@{}lrr@{}}
\toprule
Leading Terms & Share of Distinct July Descriptions & Share of July
Entries \\
\midrule

Chicken, Breast, Soup & 14.83\% & 10.12\% \\
Cheese, Cheddar, Pizza & 32.52\% & 29.92\% \\
Chocolate, Bar, Dark & 32.22\% & 38.45\% \\
Salad, Dressing, Greek & 17.01\% & 19.15\% \\
Unassigned & 3.42\% & 2.36\% \\
\bottomrule
\end{tabular}
\end{table}

\begin{figure}[htbp]
\centering
\begin{tikzpicture}
\begin{axis}[xbar,bar width=7pt,width=11.5cm,height=7cm,xlabel={Share of July Observations (Percent)},xmin=0,xmax=43,ymin=-.6,ymax=4.6,ytick={0,1,2,3,4},yticklabels={Unassigned,Salad Terms,Chocolate Terms,Cheese Terms,Chicken Terms},xtick={0,10,20,30,40},xmajorgrids=true,legend style={at={(.5,1.13)},anchor=south,legend columns=2}]
\addplot[IlliniBlue,fill=IlliniBlue,bar shift=5pt] coordinates {(14.833940,4) (32.519657,3) (32.220397,2) (17.005046,1) (3.420960,0)};
\addlegendentry{Distinct Descriptions}
\addplot[IlliniOrange,fill=IlliniOrange,bar shift=-5pt] coordinates {(10.115437,4) (29.924835,3) (38.454419,2) (19.150293,1) (2.355016,0)};
\addlegendentry{Food Entries}
\end{axis}\end{tikzpicture}

\caption{Weighting descriptions by their observed entry frequency changes the apparent prominence of the same fixed topics.}\label{fig:M05_topic_denominators}
\end{figure}

The chocolate-related component receives a larger share when repeated
entries determine the weights. This result concerns the language and
repetition of recorded entries. The component includes other terms and
does not define a validated food group, calorie share, or measure of
nutritional quality. Assigning a dominant component also discards
information about mixtures. The 583 unassigned descriptions expose a
vocabulary-coverage problem that would be hidden by forcing every text
into a topic.

\href{https://github.com/blueprints-for-text-analytics-python/blueprints-text/blob/master/ch08/Topic_Modeling_Clustering.ipynb}{Blueprints
for Text Analytics, Chapter 8} provides an NMF topic-model example.
\href{https://maartengr.github.io/BERTopic/getting_started/representation/representation.html}{BERTopic's
representation documentation} offers separate methods for describing
topics. An LLM-generated topic name should remain attached to the
component's terms and inspected examples. A fluent name supplies no
independent evidence that a component measures preference,
healthfulness, or a commercial segment.

A useful AI-assisted extension changes one choice at a time: the number
of topics, frequency weighting during fitting, or the vocabulary
threshold. Compare representative descriptions and unmatched texts as
well as reconstruction error. Improving the numerical objective alone
does not establish that a proposed topic label answers the manager's
question.

\FloatBarrier
\Needspace{8\baselineskip}
Python source:
\href{../03_Data_and_Code/mHealth/recurring_examples/topics.py}{topics.py}.
From its directory, run \texttt{python\ topics.py}.

\begin{Verbatim}[breaklines,breakanywhere,fontsize=\small]
"""Local topic decomposition of food descriptions with a future-text check."""

import json
import numpy as np
from sklearn.decomposition import NMF
from sklearn.feature_extraction.text import TfidfVectorizer
from common import SEED, load_food, save_result


if __name__ == "__main__":
    food = load_food()
    prior = food.loc[(food.date < "2016-07-01") & food.description.ne("")]
    training_text = sorted(prior.description.unique())
    july = food.loc[
        food.date.between("2016-07-01", "2016-07-31") & food.description.ne("")
    ]
    counts = july.groupby("description").size().sort_index()
    vectorizer = TfidfVectorizer(
        min_df=3,
        max_df=0.95,
        max_features=3000,
        ngram_range=(1, 2),
        stop_words="english",
    )
    train_matrix = vectorizer.fit_transform(training_text)
    model = NMF(
        n_components=4, init="nndsvda", random_state=SEED, max_iter=500, tol=1e-4
    )
    model.fit(train_matrix)
    future_matrix = vectorizer.transform(counts.index)
    # Put topic term vectors on a common L2 scale before comparing loadings.
    weights = model.transform(future_matrix) * np.linalg.norm(model.components_, axis=1)
    assigned = weights.sum(axis=1) > 1e-12
    labels = np.where(assigned, weights.argmax(axis=1), -1)
    terms = vectorizer.get_feature_names_out()
    rows = []
    for topic in range(4):
        selected = labels == topic
        top = np.argsort(-model.components_[topic], kind="stable")[:6]
        rows.append(
            {
                "topic": topic + 1,
                "terms": terms[top].tolist(),
                "distinct_description_share": float(selected.mean()),
                "food_entry_share": float(
                    counts.to_numpy()[selected].sum() / counts.sum()
                ),
            }
        )
    result = {
        "training_distinct_descriptions": len(training_text),
        "training_entries": len(prior),
        "vocabulary_terms": len(terms),
        "july_distinct_descriptions": len(counts),
        "july_entries": int(counts.sum()),
        "july_new_descriptions": int((~counts.index.isin(training_text)).sum()),
        "july_zero_vocabulary_descriptions": int(
            (future_matrix.getnnz(axis=1) == 0).sum()
        ),
        "july_unassigned_descriptions": int((~assigned).sum()),
        "july_unassigned_entries": int(counts.to_numpy()[~assigned].sum()),
        "iterations": int(model.n_iter_),
        "topics": rows,
    }
    assert model.n_iter_ < model.max_iter
    assert np.isclose(
        sum(row["food_entry_share"] for row in rows)
        + result["july_unassigned_entries"] / counts.sum(),
        1,
    )
    save_result("topics", result)
    print(json.dumps(result, indent=2))
\end{Verbatim}

\hypertarget{M06-m6.-causal-effects-and-marketing-experiments}{%
\chapter{Causal Effects and Marketing
Experiments}\label{M06-m6.-causal-effects-and-marketing-experiments}}

\noindent A retailer wants to know whether an email campaign increases
contribution. Customers who open marketing emails often purchase more,
but they may already be more engaged. Comparing openers with nonopeners
can therefore reward a campaign for behavior it did not cause. The
decision requires a comparison between outcomes under alternative
actions for the same target population.

\hypertarget{M06-define-the-treatment-outcome-and-estimand}{%
\section{Define the Treatment, Outcome, and
Estimand}\label{M06-define-the-treatment-outcome-and-estimand}}

\noindent Let \(Y_i(1)\) be customer \(i\)'s 30-day contribution if assigned an
email and \(Y_i(0)\) the contribution if not assigned. The individual
effect is \(Y_i(1)-Y_i(0)\). Only one of these outcomes is observed for
a given customer under a single assignment. A causal design supplies a
defensible way to learn an average comparison despite that missing
counterfactual.

For a target population, the average treatment effect is

\[
\operatorname{ATE}=E[Y_i(1)-Y_i(0)].
\]

The expectation refers to the stated population and outcome window. An
average treatment effect on the treated uses the population that
actually received treatment. An intention-to-treat effect uses
assignment as the treatment, even if some assigned customers never open
the email. These are different estimands. The managerial decision to
send an email often corresponds closely to the assignment effect,
because opening is outside the firm's direct control.

A well-defined treatment includes the offer, timing, channel, and
implementation. ``Digital marketing'' is too broad to identify a single
comparison. The outcome should include economically relevant costs and
have enough follow-up to capture returns or delayed purchases. If
customers can affect one another, an individual's outcome may depend on
other customers' assignments. The unit and design then need to account
for interference rather than silently assuming it away.

\hypertarget{M06-why-the-treated-untreated-difference-can-be-biased}{%
\section{Why the Treated-Untreated Difference Can Be
Biased}\label{M06-why-the-treated-untreated-difference-can-be-biased}}

\noindent The observed difference between treated and untreated customers can be
decomposed as

\[
E[Y\mid D=1]-E[Y\mid D=0]
=E[Y(1)-Y(0)\mid D=1]
+\{E[Y(0)\mid D=1]-E[Y(0)\mid D=0]\}.
\]

Here \(D\) indicates treatment. The first term is the average treatment
effect among treated customers. The bracketed term compares their
untreated outcomes with untreated customers' outcomes. If the firm
targets people who would purchase anyway, that selection term can be
positive. If it targets customers in serious decline, the term can be
negative. The direction follows the assignment process, not a universal
rule about observational data.

An unobserved engagement variable may affect both email opening and
purchasing. Adjusting for pre-treatment engagement measures can be
useful if they capture the relevant confounding, but controlling for a
variable does not guarantee this condition. A post-treatment variable
such as clicking the email can mediate part of its effect. Conditioning
on customers who click changes the population and can introduce
selection.

\begin{figure}[htbp]
\centering
% Original synthetic teaching figure; see corresponding module.
\begin{tikzpicture}
\node[box,text width=2.8cm] (z) at (0,0) {Randomized\\assignment $Z$};
\node[orangebox,text width=2.8cm] (o) at (5,0) {Email opening $O$};
\node[box,text width=2.8cm] (y) at (10,0) {Purchase $Y$};
\node[box,text width=3cm] (u) at (7.5,2.5) {Prior\\engagement $U$};
\draw[edge] (z)--(o);\draw[edge] (o)--(y);\draw[edge] (u)--(o);\draw[edge] (u)--(y);
\draw[edge] (z.south)--(0,-1.4)--node[below,font=\scriptsize] {Other pathways from assignment} (10,-1.4)--(y.south);
\node[align=center,text width=11cm] at (5,-2.9) {Opening depends on engagement and assignment.\\Keep all assigned customers in the analysis.};
\end{tikzpicture}

\caption{Assignment, engagement, and opening define different causal comparisons.}\label{fig:M06_dag}
\end{figure}

The firm randomizes email assignment, while customers choose whether to
open. Assignment can affect opening, and engagement can affect both
opening and purchases. Restricting attention to openers creates a
selected group. An assignment-based analysis preserves the randomized
comparison when its implementation and outcome measurement are valid.

\hypertarget{M06-designing-an-experiment-around-the-decision}{%
\section{Designing an Experiment Around the
Decision}\label{M06-designing-an-experiment-around-the-decision}}

\noindent An experiment needs an eligible population, assignment unit, treatment
implementation, primary outcome, horizon, and analysis plan. Randomize
at a unit that corresponds to exposure and likely spillovers. A
customer-level email experiment may be appropriate when accounts act
independently. A store-level display intervention requires store
assignment or a design that accounts for shared exposure. Household,
geographic, or network spillovers can require broader clusters.

Balance treatment groups through randomization, with blocking or
stratification when important pre-treatment characteristics are known.
For example, randomize within customer-tenure groups to improve balance
and precision. Record assignment before contact and retain every
assigned unit in the primary analysis when estimating the
intention-to-treat effect. Excluding nonopeners, customers who later
unsubscribe, or stores with disappointing execution can undermine the
randomized comparison.

Specify a primary outcome and guardrails before viewing results.
Incremental contribution might be primary, with unsubscribe rates,
complaints, and returns as guardrails. Multiple outcomes can be
informative, but selecting the most favorable result after examining
many tests exaggerates evidence. Decide how multiple testing, sequential
monitoring, and stopping will be handled. Ordinary fixed-horizon
p-values do not generally retain their advertised interpretation under
unrestricted repeated peeking and outcome selection.

\hypertarget{M06-case-study-estimate-an-email-assignment-effect}{%
\section{Case Study: Estimate an Email Assignment
Effect}\label{M06-case-study-estimate-an-email-assignment-effect}}

\noindent In a synthetic experiment, 1,000 customers are assigned an email and
1,000 are assigned no email. Within 30 days, 130 assigned customers and
100 controls make at least one eligible purchase. The estimated
assignment effect on purchase incidence is

\[
\hat\tau=\frac{130}{1000}-\frac{100}{1000}=0.03,
\]

or three percentage points. Relative to the control rate of 10\%, this
is a 30\% relative increase. The two descriptions use different
denominators; reporting only ``30\% lift'' can obscure the absolute
change.

Under independent customer assignments and the usual large-sample
approximation, an unpooled standard error is

\[
\widehat{\operatorname{SE}}(\hat\tau)
=\sqrt{\frac{0.13(0.87)}{1000}+\frac{0.10(0.90)}{1000}}
\approx0.0143.
\]

An approximate 95\% interval is {[}0.0021, 0.0579{]}, or about 0.21 to
5.79 percentage points. This interval describes sampling uncertainty
under the stated design and approximation. It does not account for an
unmeasured outcome-definition error or interference between assigned
customers.

Suppose an additional purchase contributes \$40 and email delivery costs
\$0.50 per assigned customer. If the offer has no additional discount
cost and purchase incidence is the relevant incremental quantity,
estimated incremental contribution is \$40(0.03)-\$0.50=\$0.70 per
assigned customer. Applying the same transformation to the effect
interval gives approximately \$-0.42 to \$1.82. A positive point
estimate therefore coexists with economically consequential uncertainty.
A decision can combine that uncertainty with rollout costs, the value of
another test, and guardrail outcomes.

The margin assumption matters. If treatment also changes basket size,
returns, or discount expense, purchase incidence alone is insufficient.
An experiment measuring contribution directly may better match the
objective.

\hypertarget{M06-power-precision-and-the-minimum-useful-effect}{%
\section{Power, Precision, and the Minimum Useful
Effect}\label{M06-power-precision-and-the-minimum-useful-effect}}

\noindent Power planning begins with the smallest effect that could change the
decision, not with a convenient number of rows. For two equally sized
independent groups and a binary outcome with baseline rate near \(p\), a
rough per-group sample-size approximation is

\[
n\approx\frac{2p(1-p)(z_{1-\alpha/2}+z_{1-\beta})^2}{\delta^2},
\]

where \(\delta\) is the absolute effect in probability units, \(\alpha\)
is the two-sided significance level, and \(1-\beta\) is target power.
This approximation uses a common baseline variance and normal critical
values. More precise planning can use the expected rates under the
alternative and the actual assignment design.

With \(p=0.10\), \(\delta=0.02\), \(\alpha=0.05\), and 80\% power, using
1.96 and 0.84 gives about 3,528 customers per group. Halving the target
effect approximately quadruples the required sample under this formula.
A two-percentage-point effect is 0.02, not 2; a unit error can produce a
meaningless sample-size recommendation.

Clustering reduces the amount of independent information. A common
approximation with equal cluster size \(m\) and intracluster correlation
\(\rho\) multiplies individual-level variance by \(1+(m-1)\rho\). For
\(m=20\) and \(\rho=0.05\), that factor is 1.95. The number of clusters,
unequal cluster sizes, assignment design, and small-sample inference
still require explicit consideration. Thousands of customers across four
stores do not provide thousands of independent store-level treatment
assignments.

\hypertarget{M06-implementation-checks-that-protect-the-comparison}{%
\section{Implementation Checks That Protect the
Comparison}\label{M06-implementation-checks-that-protect-the-comparison}}

\noindent Verify that the assignment counts match the intended allocation and
investigate material sample-ratio discrepancies. Check whether tracking
differs between treatment and control, whether treatment changes the
probability an outcome is observed, and whether exclusions were applied
consistently. Pre-treatment balance is a diagnostic, not proof that
every causal assumption holds. In a correctly randomized experiment,
chance imbalances can occur.

Noncompliance changes the estimand. If only 40\% of assigned customers
open the email, the intention-to-treat effect remains the effect of
assignment under the observed implementation. Dividing it by an
opening-rate difference identifies a complier effect only under
additional instrumental-variable conditions, including an appropriate
exclusion restriction and monotonicity. The firm may care about the
assignment effect even when a different scientific question concerns
actual exposure.

Use pre-treatment outcomes to improve precision when justified. A
regression adjustment can compare outcomes while accounting for baseline
spending or stratification variables. Avoid selecting an adjustment
because it makes the treatment coefficient significant. Report the
planned analysis and explain any departures. Analyze heterogeneity
cautiously, especially when many subgroups are examined after seeing the
main result.

\hypertarget{M06-observational-adjustment-and-its-limits}{%
\section{Observational Adjustment and Its
Limits}\label{M06-observational-adjustment-and-its-limits}}

\noindent When randomization is unavailable, a common approach assumes that
treatment is conditionally independent of potential outcomes given
observed pre-treatment covariates, with overlap in treatment
probabilities. Regression adjustment, matching, and inverse-probability
weighting implement different ways to construct a comparison under that
condition. The central question is whether the observed covariates make
the treated and untreated outcomes comparable.

A propensity score is the treatment probability given those covariates.
Extreme estimated scores indicate limited overlap and can create
unstable weights. Trimming or restricting the sample changes the
population for which the effect is estimated. Report that population
explicitly. A flexible prediction algorithm for treatment assignment
does not establish that all confounders were measured.

A doubly robust estimator combines an outcome model and a treatment
model and can be consistent when one nuisance model is correctly
specified, under the underlying identification conditions and regularity
assumptions. This robustness does not cover unmeasured confounding. For
advanced implementations, sample splitting and cross-fitting help
control overfitting of nuisance functions; the design remains the
foundation of the causal interpretation.

\hypertarget{M06-evaluating-a-recommendation-policy-from-logged-exposure}{%
\subsection{Evaluating a Recommendation Policy From Logged
Exposure}\label{M06-evaluating-a-recommendation-policy-from-logged-exposure}}

\noindent A retailer has records from an existing recommendation policy and wants
to evaluate a different one. The existing policy shows products
selectively: frequent buyers may usually see A, while occasional buyers
usually see B. Comparing their observed click rates mixes product
response with the customer populations receiving each recommendation.
\href{https://proceedings.mlr.press/v48/schnabel16.html}{Schnabel and
coauthors (2016)} formulate recommendation learning and evaluation as
problems of selected observation and use propensity weighting to address
that selection under explicit conditions.

Consider a single recommendation opportunity per customer. Let \(X_i\)
be information available before display, \(A_i\) the displayed item, and
\(Y_i\) a click within the specified response window. The logging policy
assigns item \(a\) with probability \(p_0(a\mid X_i)\). A target policy
assigns it with probability \(\pi(a\mid X_i)\). The target is the
expected clicks per opportunity if eligible customers receive
recommendations from \(\pi\). An inverse-propensity estimator is

\[
\widehat V(\pi)=\frac{1}{n}\sum_{i=1}^{n}
\frac{\pi(A_i\mid X_i)}{p_0(A_i\mid X_i)}Y_i.
\]

The ratio changes how much each logged outcome represents the target
policy. For a policy that always recommends A, outcomes following B
receive zero weight. An A outcome observed under a 10\% display
probability receives weight ten. The denominator \(n\) counts all
eligible logged opportunities, including those assigned zero weight.
Dividing only by the number of matching displays estimates a different
quantity.

Interpretation requires a well-defined action and outcome, conditional
exchangeability of display given \(X_i\), and positive logging
probability wherever the target policy assigns positive probability. A
known randomized logging policy can supply the assignment condition.
Estimating probabilities from observational logs requires the additional
substantive claim that the recorded context captures the relevant
selection. The target population and response process must also match
the logged environment. For a slate of products, position and other
displayed items can affect the response; a single-item propensity
generally does not describe assignment of the complete slate.

The weighting identity explains the adjustment. Conditional on \(X=x\),
averaging over logged actions gives
\(\sum_a p_0(a\mid x)\{\pi(a\mid x)/p_0(a\mid x)\}E[Y(a)\mid X=x]\). The
logging probability cancels, leaving the response weighted by the target
policy. Cancellation requires positive support. If an item was never
eligible for a customer group, multiplying the observed records by
larger weights cannot reveal that group's response to it.

\hypertarget{M06-case-study-exposure-adjustment-reverses-the-preferred-item}{%
\subsection{Case Study: Exposure Adjustment Reverses the Preferred
Item}\label{M06-case-study-exposure-adjustment-reverses-the-preferred-item}}

\noindent Suppose 1,000 synthetic customers comprise 500 frequent buyers and 500
occasional buyers. The logging policy randomly assigns A within each
group with probabilities 0.90 and 0.10, respectively. B receives the
remaining probability. The constructed counts are:

\begin{table}[H]
\centering\small
\begin{tabular}{@{}llrrr@{}}
\toprule
Customer Group & Item & Display Probability & Displays & Clicks \\
\midrule

Frequent & A & 0.90 & 450 & 270 \\
Frequent & B & 0.10 & 50 & 25 \\
Occasional & A & 0.10 & 50 & 10 \\
Occasional & B & 0.90 & 450 & 180 \\
\bottomrule
\end{tabular}
\end{table}

A receives 280 clicks from 500 displays, or 56\%. B receives 205 from
500, or 41\%. Yet A's within-group rates are 60\% and 20\%, while B's
are 50\% and 40\%. Applied to the same equally weighted customer
population, A gives 40\% and B gives 45\%. A's displayed customers
disproportionately come from the group with high response to either
item.

The inverse-propensity calculations recover the common-population
comparison:

\[
\widehat V(A)=\frac{270/0.90+10/0.10}{1000}=0.40,
\qquad
\widehat V(B)=\frac{25/0.10+180/0.90}{1000}=0.45.
\]

\begin{figure}[htbp]
\centering
\begin{tikzpicture}
\begin{axis}[ybar,bar width=21pt,width=13cm,height=7cm,
  symbolic x coords={Displayed customers,Common population},
  xtick=data,ymin=0,ymax=65,ytick={0,10,20,30,40,50,60},
  ylabel={Clicks per 100 opportunities},
  enlarge x limits=.45,ymajorgrids=true,
  nodes near coords,point meta=y,
  every node near coord/.append style={font=\small},
  legend style={at={(.5,-.15)},anchor=north,legend columns=2}]
\addplot[fill=IlliniBlue,draw=IlliniBlue] coordinates
  {(Displayed customers,56) (Common population,40)};
\addplot[fill=IlliniOrange,draw=IlliniOrange] coordinates
  {(Displayed customers,41) (Common population,45)};
\legend{Item A,Item B}
\end{axis}
\end{tikzpicture}

\caption{Selected display populations favor A, while comparison on the same population favors B.}\label{fig:M06_exposure_reversal}
\end{figure}

The raw bars describe response among each item's displayed customers.
The adjusted bars describe the same 50:50 customer mix for both
policies. The five-percentage-point advantage for B is imposed by the
synthetic cell counts and randomized logging design. It is not an
estimate from the cited study.

Large weights reduce precision. For A, the 50 occasional buyers each
represent ten eligible customers, so a small number of outcomes carries
substantial influence. The diagnostic \((\sum_i w_i)^2/\sum_i w_i^2\)
equals 180 for either policy in this example, despite 1,000 logged
opportunities. This effective sample size describes weight
concentration; it does not replace a variance calculation that respects
assignment and dependence. Weight clipping trades variance for bias, and
self-normalization changes the finite-sample estimator. Report these
choices and retain an untouched evaluation sample when choosing between
policies.

\hypertarget{M06-difference-in-differences-and-a-counterfactual-trend}{%
\section{Difference-in-Differences and a Counterfactual
Trend}\label{M06-difference-in-differences-and-a-counterfactual-trend}}

\noindent A retailer introduces an app in treated markets while comparable markets
remain untreated. A before-after difference in treated markets mixes the
app's effect with common demand changes. Difference-in-differences
subtracts the comparison group's change:

\[
\hat\tau_{\mathrm{DiD}}=(\bar Y_{T,post}-\bar Y_{T,pre})-(\bar Y_{C,post}-\bar Y_{C,pre}).
\]

In a synthetic example, treated-market sales rise from 120 to 150 and
comparison-market sales rise from 100 to 110. The DiD estimate is 20
sales units. Under the identifying parallel-trends condition, the
treated markets' untreated counterfactual would have risen to 130.
Different starting levels are compatible with this design; the condition
concerns untreated changes.

\begin{figure}[htbp]
\centering
% Original synthetic teaching figure; see corresponding module.
\begin{tikzpicture}
\begin{axis}[xmin=-.12,xmax=1.38,ymin=90,ymax=162,xtick={0,1},xticklabels={Before,After},ylabel={Outcome (illustrative units)},legend style={at={(.5,1.08)},anchor=south,legend columns=3}]
\addplot[IlliniOrange,mark=*,very thick] coordinates {(0,120)(1,150)};
\addplot[IlliniBlue,mark=square*,very thick] coordinates {(0,100)(1,110)};
\addplot[Storm,dashed,mark=o,very thick] coordinates {(0,120)(1,130)};
\legend{Treated,Comparison,Treated counterfactual}
\draw[<->,IlliniOrange,very thick] (axis cs:1.08,130)--(axis cs:1.08,150);
\node[anchor=west,align=left] at (axis cs:1.1,140) {Effect\\$=20$};
\end{axis}
\end{tikzpicture}

\caption{Difference-in-differences subtracts the comparison group’s change to construct the counterfactual.}\label{fig:M06_did}
\end{figure}

The dashed line is a counterfactual implied by parallel untreated
trends, not an observed sales series. The vertical gap between observed
treated sales and that line is the estimated effect. A
treatment-specific local shock at the same time would challenge the
interpretation even if the pre-treatment levels looked similar.

Inspect multiple pre-treatment periods, institutional timing, possible
anticipation, spillovers, and changes in sample composition. Similar
pre-trends can support the design's plausibility, but a nonsignificant
pre-trend test does not establish parallel counterfactual trends. With
staggered adoption and heterogeneous effects, conventional two-way
fixed-effects coefficients can combine comparisons in problematic ways.
Use estimators aligned with cohort and time-specific effects, with a
clear comparison group, when extending beyond the simple two-period
design.

\hypertarget{M06-other-designs-discontinuities-and-instruments}{%
\section{Other Designs: Discontinuities and
Instruments}\label{M06-other-designs-discontinuities-and-instruments}}

\noindent A regression discontinuity design can use a treatment rule at a
threshold, such as eligibility based on a pre-existing score. Customers
just above and below the cutoff may be comparable if other determinants
of the outcome vary smoothly there and the score is not precisely
manipulated. The resulting effect is local to the cutoff and the
relevant compliance behavior. Bandwidth, functional form, density
patterns, and discontinuities in other predetermined variables deserve
inspection.

An instrumental variable changes treatment while affecting the outcome
only through the treatment channel, under the design's conditions.
Relevance is an empirical first-stage requirement; exclusion is a
substantive condition. A strong first stage cannot establish exclusion.
In a binary-treatment setting with independence, exclusion, and
monotonicity, the familiar local average treatment effect concerns
compliers whose treatment changes with the instrument. It need not equal
the population ATE.

Marketing instruments require particularly careful reasoning because
pricing, distribution, and exposure decisions often respond to demand. A
variable correlated with treatment is not sufficient. Explain who makes
the relevant decision, when the instrument is determined, and which
other channels could connect it to outcomes. If those channels remain
credible, the causal interpretation needs revision.

\hypertarget{M06-spillovers-and-experiment-units}{%
\section{Spillovers and Experiment
Units}\label{M06-spillovers-and-experiment-units}}

\noindent An advertising campaign can shift purchases between stores or products
rather than increase total demand. A customer-level outcome may show a
gain while chain-wide contribution changes little. Define the outcome at
the level relevant to the managerial objective and examine displacement
where it is plausible. The same issue arises when customers share offers
or when a platform changes recommendations for multiple users.

\begin{figure}[htbp]
\centering
% Original synthetic teaching figure; see corresponding module.
\begin{tikzpicture}[v/.style={circle,draw=IlliniBlue,minimum size=11mm},t/.style={v,fill=IlliniOrange}]
\node[t] (t) at (0,0) {T};\node[v] (a) at (2,1.5) {U};\node[v] (b) at (2,-1.5) {U};\node[v] (c) at (4,0) {U};
\draw[accent] (t)--(a);\draw[accent] (t)--(b);\draw[edge,dashed] (a)--(c);\draw[edge,dashed] (b)--(c);
\draw[IlliniBlue,dashed,rounded corners] (-1,-2.2) rectangle (5,2.2);
\node[align=center] at (2,-3) {One connected household or market\\can form a randomization cluster.};
\node[box,text width=4cm] at (8,1) {T: assigned treatment\\U: assigned control};
\node[align=left,text width=4cm] at (8,-1.1) {Control status does not guarantee absence of exposure through other people.};
\end{tikzpicture}

\caption{A treatment can affect both assigned units and nearby untreated units.}\label{fig:M06_spillovers}
\end{figure}

An assigned store can gain sales while diverting purchases from another
store. If the comparison store loses customers, subtracting its decline
can overstate the increase in total chain sales. Cluster randomization,
geographic separation, or designs that vary treatment saturation may
help identify the relevant effects, subject to their own practical
constraints.

\hypertarget{M06-case-study-recommendations-supply-information-and-exposure}{%
\subsection{Case Study: Recommendations Supply Information and
Exposure}\label{M06-case-study-recommendations-supply-information-and-exposure}}

\noindent A recommendation can increase consumption by bringing a product to
attention or by changing beliefs about its quality.
\href{https://www.econstor.eu/handle/10419/271773}{Aridor, Gonçalves,
Kluver, Kong, and Konstan's MovieLens experiment}, in its November 2022
working-paper version, separates recommendation and exposure-only
conditions and elicits users' beliefs. The results support an
informational role: recommendations influence beliefs and consumption
beyond simply exposing a movie. MovieLens provides recommendations
without selling or hosting the movies, which helps define the incentives
in this setting.

A retailer studying the same mechanism could distinguish ordinary
display from an explicit personalized endorsement while keeping the
eligible product set and available information controlled. The primary
outcome would be purchases or contribution per assigned customer over a
fixed horizon. Belief measures would address how the endorsement
operates. The experiment must account for the possibility that asking
about a product itself increases attention to it. Comparisons based only
on customers who subsequently click would condition on a response that
the endorsement can change.

The displayed information can also affect the labels collected
afterward.
\href{https://experts.umn.edu/en/publications/is-seeing-believing-how-recommender-interfaces-affect-users-opini/}{Cosley
and coauthors (2003)} experimentally vary predicted ratings shown to
MovieLens users and find that subsequent ratings tend to move toward the
displayed prediction. A rating collected after a recommendation
therefore reflects the evaluation context as well as the user's
experience. A model trained to reproduce those ratings can partly
reproduce its own previous influence. Collecting some evaluations before
showing predicted ratings provides a comparison for that mechanism.

An experiment on an adaptive recommender also needs a rule for
retraining. If both treatment and control data update a shared model,
one customer's assignment may change another customer's recommendations.
Freezing the model, maintaining separate models, or assigning
sufficiently separated populations evaluates different policies. State
which policy and learning process the experiment implements. A short
experiment with a frozen model identifies a different intervention from
deploying a model that learns from its own recommendations over several
months.

\hypertarget{M06-ai-assisted-workflow-and-practice}{%
\section{AI-Assisted Workflow and
Practice}\label{M06-ai-assisted-workflow-and-practice}}

\noindent Ask an assistant to state the estimand, assignment unit, outcome
horizon, and comparison group before writing estimation code. For the
email example, require the conversion effect in percentage points, the
unpooled standard error, and the contribution calculation. Check the
arithmetic independently. Ask it to explain why restricting the sample
to email openers changes the analysis.

For observational data, draw or describe the assignment process before
selecting controls. Request an overlap diagnostic and a report of which
observations receive large weights. For DiD, ask the assistant to show
observed and counterfactual series separately and identify the
assumptions that connect them. A convincing plot should not substitute
for those conditions.

If comparison-market sales rise by 25 rather than 10, the DiD estimate
becomes five. If email cost doubles to \$1, the point estimate of
incremental contribution falls to \$0.20. If treatment raises returns,
explain why a conversion effect can overstate the contribution effect.
If two stores share most customers, identify why store assignment may
create interference.

For the exposure example, ask an assistant to compute both the raw click
rates and inverse-propensity policy values from the four aggregate rows,
using the display counts when calculating denominators and weight
concentration. Verify the weighted results independently by averaging
the two group-specific response rates. Then set A's display probability
among occasional buyers to zero in a proposed logging design and explain
why the always-A policy becomes unidentified from those logs. The
implementation should flag the unsupported action rather than silently
drop that group.

\hypertarget{M06-readings}{%
\section{Readings}\label{M06-readings}}

\noindent \href{https://matheusfacure.github.io/python-causality-handbook/landing-page.html}{Facure's
causal-inference handbook} supplies accessible implementations of
several designs.
\href{https://exp-platform.com/Documents/GuideControlledExperiments.pdf}{Kohavi
and colleagues' practical guide} discusses controlled online
experiments.
\href{https://doi.org/10.1016/j.jeconom.2020.12.001}{Callaway and
Sant'Anna (2021)} develop difference-in-differences with multiple
periods.

\hypertarget{M06-mhealth-planning-a-randomized-suggestion-test}{%
\section{mHealth: Planning a Randomized Suggestion
Test}\label{M06-mhealth-planning-a-randomized-suggestion-test}}

\noindent The manager has a candidate suggestion policy and wants to know whether
offering it changes customer behavior. The historical logs contain no
documented randomized assignment or recommendation exposure. They can
supply a distribution for planning precision and reveal how a proposed
outcome is measured. An effect estimate requires the assignment and
outcome records from a suitable new design.

Consider a user-randomized trial comparing the current interface with an
offer of the new suggestions. Define the primary logging outcome as the
fraction of the next 30 calendar days with at least one recorded food
entry. Analyze all assigned eligible users under their original
assignment, including users who never return. A suggestion may increase
logging without changing consumption. A nutrition outcome therefore
needs its own validated measurement and a plan for missingness.
Post-assignment completeness screens could remove users differently
across arms and alter the intention-to-treat comparison.

For historical planning, require at least 30 recorded days during
January through May. Among the 327 eligible users, June's logging
fraction averages 0.6408 with a standard deviation of 0.3542.
Twenty-three users have no June entries and receive a zero logging
fraction. That zero is meaningful for observed logging. It says nothing
about zero consumption, nutritional adequacy, or contract cancellation.

With independent users, equal arm sizes, a two-sided 5\% test, 80\%
power, and the historical standard deviation treated as the prospective
outcome standard deviation, the normal approximation is

\[
n_{\text{per arm}}\approx
\frac{2\sigma^2(z_{0.975}+z_{0.80})^2}{\Delta^2}.
\]

An imposed minimum detectable effect of 0.05 in the logging fraction,
equivalent to 1.5 days out of 30, gives 788 users per arm after rounding
upward. This is a planning scenario using historical variance. Future
eligibility, outcome distributions, unequal assignment, clustering, and
the intended analysis can change the required size. The 327-user
historical subset does not supply that many independent trial
participants.

\begin{figure}[htbp]
\centering
\begin{tikzpicture}
\begin{axis}[xlabel={Independent Users per Arm},ylabel={Detectable Change in Logging Fraction},xmin=100,xmax=1000,ymin=0,ymax=.16,ytick={0,.05,.10,.15},yticklabel style={/pgf/number format/fixed,/pgf/number format/precision=2},grid=major]
\addplot[IlliniBlue,domain=100:1000,samples=80] {sqrt(2*0.354206889^2*(1.959963985+0.841621234)^2/x)};
\addplot[IlliniOrange,dashed] coordinates {(100,.05)(1000,.05)};
\addplot[only marks,IlliniOrange,mark=square*,mark size=3pt] coordinates {(788,.05)};
\node[anchor=south west,font=\footnotesize] at (axis cs:600,.063) {788 per arm for 0.05};
\node[anchor=north east,font=\footnotesize,align=right] at (axis cs:980,.151) {Historical SD: 0.3542\\Two-sided 5\% test; 80\% power};
\end{axis}
\end{tikzpicture}

\caption{The curve holds the historical outcome variance and testing assumptions fixed. More independent users narrow the detectable effect; adding days within the same user does not supply the same information.}\label{fig:M06_mhealth_precision}
\end{figure}

A randomization simulation assigns the same observed June outcomes to
nearly equal arms 2,000 times while imposing a sharp zero effect. The
standard deviation of the resulting mean difference is 0.0395, and the
central 95\% interval is approximately {[}-0.0798, 0.0740{]}. These
differences arise from assignment variation alone. Repeatedly treating
each user-day as an independent trial participant would conceal
dependence within users and overstate the effective sample size.

The protocol should specify randomization before any policy exposure,
record assignment and actual exposure separately, and keep the outcome
window fixed. For nutrition assessments, collect reliable intake or
appropriate validation measures independently of how often the app is
used. Domain-based restrictions should be specified in advance, with
their effect on the target population and missing-outcome analysis made
explicit.

The program uses a local random generator with seed 534. It reports the
assumed effect separately from the simulated differences and never
labels an artificial assignment as an observed intervention.

\FloatBarrier
\Needspace{8\baselineskip}
Python source:
\href{../03_Data_and_Code/mHealth/recurring_examples/m06.py}{m06.py}.
From its directory, run \texttt{python\ m06.py}.

\begin{Verbatim}[breaklines,breakanywhere,fontsize=\small]
"""Use historical logging variation to plan, not estimate, a randomized trial."""

import json
import numpy as np
from scipy.stats import norm
from common import SEED, load_inputs, save_result

food, daily, baskets = load_inputs()
prior = daily.loc[daily.date.between("2016-01-01", "2016-05-31")]
counts = prior.groupby("Userid").size()
eligible = counts[counts >= 30].index.sort_values()
june = daily.loc[daily.date.between("2016-06-01", "2016-06-30")]
# A zero means no recorded day in June, not zero consumption or poor health.
outcome = june.groupby("Userid").size().reindex(eligible, fill_value=0) / 30
sigma = float(outcome.std(ddof=1))
effect = 0.05  # Imposed minimum detectable increase in the fraction of logged days.
critical = norm.ppf(0.975) + norm.ppf(0.80)
per_arm = int(np.ceil(2 * sigma**2 * critical**2 / effect**2))
# Illustrative balanced assignment under a sharp zero-effect simulation.
rng = np.random.default_rng(SEED)
values = outcome.to_numpy()
differences = []
for repetition in range(2000):
    order = rng.permutation(len(values))
    treated, control = np.array_split(order, 2)
    differences.append(values[treated].mean() - values[control].mean())
result = {
    "eligible_users": len(eligible),
    "june_mean_logging_fraction": float(outcome.mean()),
    "june_sd": sigma,
    "zero_logged_days": int(outcome.eq(0).sum()),
    "assumed_effect": effect,
    "normal_approximation_users_per_arm": per_arm,
    "null_simulation_mean_difference": float(np.mean(differences)),
    "null_simulation_sd": float(np.std(differences, ddof=1)),
    "null_simulation_95_percent_interval": np.quantile(
        differences, [0.025, 0.975]
    ).tolist(),
}
save_result("M06", result)
print(json.dumps(result, indent=2))
\end{Verbatim}

\hypertarget{M06-python-calculations-for-the-synthetic-cases}{%
\section{Python Calculations for the Synthetic
Cases}\label{M06-python-calculations-for-the-synthetic-cases}}

\noindent The inputs below are imposed teaching values. Running the program
creates the labeled synthetic CSV files and prints the corresponding
calculations. They are separate from the mHealth estimates. Keep the
script beside \texttt{example\_helpers.py} in
\texttt{03\_Data\_and\_Code/Lecture\_Examples}.

\FloatBarrier
\Needspace{8\baselineskip}
Python source:
\href{../03_Data_and_Code/Lecture_Examples/m06_synthetic.py}{m06\_synthetic.py}.
From its directory, run \texttt{python\ m06\_synthetic.py}.

\begin{Verbatim}[breaklines,breakanywhere,fontsize=\small]
"""Synthetic numerical examples for M6; run from any working directory."""

import itertools
import json
import math
from example_helpers import write_csv, ols_slope, shares

r = {}
se = math.sqrt(0.13 * 0.87 / 1000 + 0.1 * 0.9 / 1000)
ci = [0.03 - 1.96 * se, 0.03 + 1.96 * se]
r["M06"] = {
    "assignment_effect": 0.03,
    "standard_error": se,
    "effect_ci": ci,
    "contribution_per_assigned_customer": 40 * 0.03 - 0.5,
    "contribution_ci": [40 * x - 0.5 for x in ci],
    "sample_per_arm": 2 * 0.1 * 0.9 * (1.96 + 0.84) ** 2 / 0.02**2,
    "design_effect": 1 + (20 - 1) * 0.05,
    "difference_in_differences": (150 - 120) - (110 - 100),
}
write_csv(
    "M06_experiment_synthetic.csv",
    ["assignment", "assigned_customers", "purchasers"],
    [("email", 1000, 130), ("control", 1000, 100)],
)
write_csv(
    "M06_did_synthetic.csv",
    ["group", "before", "after"],
    [("treated", 120, 150), ("comparison", 100, 110)],
)

exposures = [
    ("frequent", "A", 0.9, 450, 270),
    ("frequent", "B", 0.1, 50, 25),
    ("occasional", "A", 0.1, 50, 10),
    ("occasional", "B", 0.9, 450, 180),
]
write_csv(
    "M06_exposure_synthetic.csv",
    ["customer_group", "item", "logging_probability", "displays", "clicks"],
    exposures,
)


def constant_policy_value(rows, item):
    groups = {row[0] for row in rows}
    selected = [row for row in rows if row[1] == item]
    if {row[0] for row in selected if row[2] > 0} != groups:
        raise ValueError(
            "Target action lacks positive logging support in a customer group"
        )
    n = sum(row[3] for row in rows)
    weighted_clicks = sum(clicks / prob for _, _, prob, count, clicks in selected)
    sum_weights = sum(count / prob for _, _, prob, count, clicks in selected)
    sum_squared_weights = sum(count / prob**2 for _, _, prob, count, clicks in selected)
    raw = sum(row[4] for row in selected) / sum(row[3] for row in selected)
    return {
        "raw_rate": raw,
        "ips_value": weighted_clicks / n,
        "effective_sample_size": sum_weights**2 / sum_squared_weights,
    }


r["M06"]["recommendation_policies"] = {
    item: constant_policy_value(exposures, item) for item in ["A", "B"]
}
for item, target in [("A", 0.4), ("B", 0.45)]:
    result = r["M06"]["recommendation_policies"][item]
    # Independent standardization uses the known 50:50 population mix.
    standardized = sum(
        0.5 * clicks / count for group, j, prob, count, clicks in exposures if j == item
    )
    assert math.isclose(result["ips_value"], target) and math.isclose(
        result["ips_value"], standardized
    )
    assert math.isclose(result["effective_sample_size"], 180)
unsupported = [
    (g, j, 0 if (g, j) == ("occasional", "A") else p, n, y)
    for g, j, p, n, y in exposures
]
try:
    constant_policy_value(unsupported, "A")
except ValueError:
    pass
else:
    raise AssertionError("An unsupported target policy must be rejected")
print(json.dumps(r, indent=2))
\end{Verbatim}

\hypertarget{M06-mhealth-prior-activity-and-experiment-precision}{%
\section{mHealth: Prior Activity and Experiment
Precision}\label{M06-mhealth-prior-activity-and-experiment-precision}}

\noindent A manager wants an experiment to detect a useful change in logging with
fewer participants. Users differ persistently in how often they record
food. Random assignment balances those differences in expectation, while
a particular assignment can still place more frequent loggers in one
arm. A pre-treatment activity measure can explain part of the outcome
variation and improve precision.

For a planned June experiment, retain users with at least 30 recorded
days before June 1. Define \(X_i\) as the fraction of May days with a
log and \(Y_i\) as the fraction of June days with a log. All 327
eligible users remain in both monthly denominators, including users with
zero observed logging days. Fit a coefficient using March logging
fractions to predict April fractions, and freeze it before June. The
CUPED-adjusted outcome is

\[
Y_i^{\mathrm{adj}}=Y_i-\theta(X_i-\overline X),
\qquad
\widehat\tau_{\mathrm{adj}}=overline Y_1^{\mathrm{adj}}-overline Y_0^{\mathrm{adj}}.
\]

The subscripts 1 and 0 denote randomized assignment arms. With a fixed
coefficient and a pre-treatment covariate unaffected by assignment, the
adjustment removes a mean-zero difference between arms in expectation.
It preserves the pooled outcome mean. Randomization supplies the
treatment-effect comparison; covariate adjustment changes its precision.
The arithmetic is the control-variate approach underlying
\href{https://exp-platform.com/cuped/}{CUPED}, introduced by Deng, Xu,
Kohavi, and Walker (2013).

For a fixed coefficient, the adjusted variance equals

\[
\operatorname{Var}(Y^{\mathrm{adj}})=\operatorname{Var}(Y)+\theta^2\operatorname{Var}(X)-2\theta\operatorname{Cov}(X,Y).
\]

At the population-optimal coefficient
\(\theta^*=\operatorname{Cov}(X,Y)/\operatorname{Var}(X)\), this becomes
\(\operatorname{Var}(Y)(1-\rho^2)\). The historical coefficient used
here is learned from different months, so its gain must be measured in
the later comparison. A poorly transferred coefficient can increase
variance. A covariate affected by treatment can also change the estimand
and bias the treatment-effect calculation.

\begin{table}[H]
\centering\small
\begin{tabular}{@{}lrr@{}}
\toprule
Historical Planning Quantity & Unadjusted & CUPED \\
\midrule

Variance of June Outcome & 0.12546 & 0.06315 \\
Exact Sharp-Null Randomization SD & 0.03918 & 0.02779 \\
SD Across 2,000 Random Assignments & 0.04032 & 0.02840 \\
Approximate Users per Arm for 0.05 & 788 & 397 \\
\bottomrule
\end{tabular}
\end{table}

The March--April coefficient is 0.55386. The table compares the same
June users and outcomes, with an adjusted-to-unadjusted variance ratio
of 0.50334. The prospective sample-size calculation assumes independent
users, equal arm sizes, a two-sided 5\% test, 80\% power, and a 0.05
absolute increase in logging fraction. It uses historical variance as a
planning input and a normal approximation. Actual effects, variance
shifts, interference, and attrition can change the needed sample.

\begin{figure}[htbp]
\centering
\begin{tikzpicture}
\begin{axis}[xlabel={Independent Users per Arm},ylabel={Detectable Change in Logging Fraction},xmin=100,xmax=1000,ymin=0,ymax=.16,ytick={0,.05,.10,.15},yticklabel style={/pgf/number format/fixed,/pgf/number format/precision=2},grid=major,legend style={at={(.98,.98)},anchor=north east}]
\addplot[IlliniBlue,solid,domain=100:1000,samples=80] {sqrt(2*0.1254625200*(1.959963985+0.841621234)^2/x)};
\addlegendentry{Unadjusted}
\addplot[only marks,IlliniBlue,mark=square*,mark size=3pt,forget plot] coordinates {(788,.05)};
\node[anchor=north,font=\footnotesize] at (axis cs:788,.045) {788};
\addplot[IlliniOrange,dashed,domain=100:1000,samples=80] {sqrt(2*0.0631502639*(1.959963985+0.841621234)^2/x)};
\addlegendentry{Prior-Period Adjustment}
\addplot[only marks,IlliniOrange,mark=square*,mark size=3pt,forget plot] coordinates {(397,.05)};
\node[anchor=north,font=\footnotesize] at (axis cs:397,.045) {397};
\addplot[Storm,densely dotted,forget plot] coordinates {(100,.05)(1000,.05)};
\end{axis}\end{tikzpicture}

\caption{Historical covariate adjustment lowers the detectable change at a given sample size; both curves use the same outcome, significance level, and power.}\label{fig:M06_cuped_precision}
\end{figure}

No experimental assignment is observed in mHealth. The randomization
exercise imposes a sharp zero-effect null, holding every recorded June
outcome fixed while repeatedly assigning 163 users to treatment and 164
to control. Its dispersion illustrates the design's sampling
variability. The analytical finite-population variance is
\(S^2(1/n_1+1/n_0)\) for each fixed outcome, and provides an independent
comparison with the simulation.

\href{https://github.com/google-marketing-solutions/feedx}{FeedX}
combines experimental design with prior-period adjustment for
shopping-feed tests. Its application also raises a choice of
randomization unit. A test of generated product descriptions may
randomize products, whereas an mHealth reminder experiment could
randomize users. Shared exposure and spillovers must be addressed at the
unit actually assigned. Repeated fixed-horizon significance tests during
monitoring require a sequentially valid procedure or a prespecified
analysis time; software output does not resolve that design choice.

\FloatBarrier
\Needspace{8\baselineskip}
Python source:
\href{../03_Data_and_Code/mHealth/recurring_examples/cuped.py}{cuped.py}.
From its directory, run \texttt{python\ cuped.py}.

\begin{Verbatim}[breaklines,breakanywhere,fontsize=\small]
"""Historical precision planning with a coefficient frozen before June."""

import json
import math
import numpy as np
import pandas as pd
from scipy.stats import norm
from common import SEED, load_food, make_daily, save_result


def adjust(outcome, baseline, theta):
    return outcome - theta * (baseline - baseline.mean())


if __name__ == "__main__":
    daily = make_daily(load_food())
    prior = daily.loc[daily.date < "2016-06-01"].groupby("Userid").size()
    users = prior.index[prior >= 30]
    monthly = {}
    for month in ["2016-03", "2016-04", "2016-05", "2016-06"]:
        period = pd.Period(month, freq="M")
        part = daily.loc[daily.date.dt.to_period("M").eq(period)]
        monthly[month] = (
            part.groupby("Userid").size().reindex(users, fill_value=0).to_numpy()
            / period.days_in_month
        )
    march, april = monthly["2016-03"], monthly["2016-04"]
    theta = np.cov(march, april, ddof=1)[0, 1] / np.var(march, ddof=1)
    baseline, outcome = monthly["2016-05"], monthly["2016-06"]
    adjusted = adjust(outcome, baseline, theta)
    assert np.isclose(adjusted.mean(), outcome.mean())
    n, n_treated = len(users), len(users) // 2
    rng = np.random.default_rng(SEED)
    differences = []
    for _ in range(2000):
        treated = np.zeros(n, dtype=bool)
        treated[rng.choice(n, n_treated, replace=False)] = True
        # Imposed sharp null: these recorded outcomes do not change with assignment.
        differences.append(
            [
                values[treated].mean() - values[~treated].mean()
                for values in [outcome, adjusted]
            ]
        )
    differences = np.asarray(differences)
    power_multiplier = 2 * (norm.ppf(0.975) + norm.ppf(0.8)) ** 2 / 0.05**2
    result = {
        "users": n,
        "coefficient_from_march_april": float(theta),
        "may_june_correlation": float(np.corrcoef(baseline, outcome)[0, 1]),
        "june_mean": float(outcome.mean()),
        "sharp_null_draws": len(differences),
        "methods": [],
    }
    for column, (name, values) in enumerate(
        [("unadjusted", outcome), ("cuped", adjusted)]
    ):
        variance = np.var(values, ddof=1)
        result["methods"].append(
            {
                "method": name,
                "outcome_variance": float(variance),
                "simulated_effect_mean": float(differences[:, column].mean()),
                "simulated_effect_sd": float(differences[:, column].std(ddof=1)),
                "exact_randomization_sd": float(
                    np.sqrt(variance * (1 / n_treated + 1 / (n - n_treated)))
                ),
                "approximate_users_per_arm_for_0_05": math.ceil(
                    power_multiplier * variance
                ),
            }
        )
    result["variance_ratio"] = float(np.var(adjusted, ddof=1) / np.var(outcome, ddof=1))
    save_result("cuped", result)
    print(json.dumps(result, indent=2))
\end{Verbatim}

\hypertarget{M07-m7.-prediction-machine-learning-and-marketing-actions}{%
\chapter{Prediction, Machine Learning, and Marketing
Actions}\label{M07-m7.-prediction-machine-learning-and-marketing-actions}}

\noindent A retailer wants to anticipate next month's customer purchases and
allocate limited service capacity. A model that predicts historical
purchases accurately may still fail when deployed: its features may
contain future information, its test set may reuse the same customers as
training, or its objective may reward errors that are costly for the
business. Prediction should therefore be designed around an explicit
deployment situation.

\hypertarget{M07-specify-the-forecast-origin-and-target}{%
\section{Specify the Forecast Origin and
Target}\label{M07-specify-the-forecast-origin-and-target}}

\noindent A prediction task needs a target variable, a forecast origin, a horizon,
and an eligible population. For customer churn, define what counts as
churn and when that status becomes observable. In a subscription
business, cancellation may be explicit. In noncontractual retail, ``no
purchase for 90 days'' is an operational label, not proof that a
relationship has ended permanently.

Construct predictors using information available at the forecast origin.
If predicting purchases during November, features might use transactions
through October 31. November support contacts, returns, or campaign
interactions are unavailable at that origin. A database may contain them
when the analyst extracts the file, making leakage easy to overlook.
Store feature timestamps or document their measurement windows.

The deployment population determines the validation design. Predicting
future behavior for existing customers differs from predicting behavior
for newly acquired customers. A random row split can place the same
customer in training and test data, giving a favorable assessment for a
task that actually requires generalization to new customers. A time
split addresses future periods, while a customer-group split addresses
new customers. Some tasks require both.

\begin{figure}[htbp]
\centering
% Original synthetic teaching figure; see corresponding module.
\begin{tikzpicture}[x=1.55cm,y=.85cm]
\foreach \x/\lab in {0/Period 1,1/Period 2,2/Period 3} {\node at (\x,1) {\lab};\node at (\x+4.5,1) {\lab};}
\node[font=\small\bfseries] at (1,2) {Random rows};\node[font=\small\bfseries] at (5.5,2) {Future-period holdout};
\foreach \y/\lab in {0/A,-1/B,-2/C} {\node[anchor=east] at (-.55,\y) {\lab};\node[anchor=east] at (3.95,\y) {\lab};}
\foreach \x/\y in {0/0,2/0,1/-1,2/-1,0/-2,1/-2} {\fill[IlliniBlue] (\x-.43,\y-.36) rectangle (\x+.43,\y+.36);}
\foreach \x/\y in {1/0,0/-1,2/-2} {\fill[IlliniOrange] (\x-.43,\y-.36) rectangle (\x+.43,\y+.36);}
\foreach \x in {4.5,5.5} {\foreach \y in {0,-1,-2} {\fill[IlliniBlue] (\x-.43,\y-.36) rectangle (\x+.43,\y+.36);}}
\foreach \y in {0,-1,-2} {\fill[IlliniOrange] (6.07,\y-.36) rectangle (6.93,\y+.36);}
\node[align=center,text width=12cm] at (3.1,-3.4) {Blue: training. Orange: test. Rows identify customers.\\For new-customer deployment, hold out complete customer rows instead.};
\end{tikzpicture}

\caption{The validation split should reproduce the information available at deployment.}\label{fig:M07_splits}
\end{figure}

A random row split can reuse a customer's history in training and test
data. Holding out complete customers evaluates generalization to new
customers; holding out later periods evaluates future performance. The
intended use determines which comparison is relevant.

\hypertarget{M07-baselines-loss-functions-and-regularization}{%
\section{Baselines, Loss Functions, and
Regularization}\label{M07-baselines-loss-functions-and-regularization}}

\noindent A baseline should be operationally meaningful. For sales forecasting,
compare with the previous period, the same period last year, or a simple
customer-history average when those quantities are available. For
classification, compare with a constant probability and a transparent
logistic regression. A complex model's incremental value is easier to
assess when these comparisons are reported on the same test
observations.

A supervised learner minimizes a loss that quantifies prediction error.
Squared error penalizes large deviations strongly; absolute error gives
a different compromise under skewness. For a binary outcome, log loss is

\[
L=-\frac{1}{n}\sum_{i=1}^n\{y_i\log\hat p_i+(1-y_i)\log(1-\hat p_i)\}.
\]

Here \(y_i\) is the observed binary outcome and \(\hat p_i\) is its
predicted probability. Confident mistakes receive large penalties. The
loss evaluates probability forecasts; it does not by itself encode the
cost of a marketing action. Select a loss appropriate to the prediction
task, then evaluate how the forecasts support the downstream decision.

Regularization constrains a flexible model to reduce overfitting. Ridge
regression adds a squared-coefficient penalty, while lasso adds an
absolute-coefficient penalty and can set some coefficients to zero.
Standardize features appropriately when the penalty depends on
coefficient scale, and fit that transformation using training data. The
regularization strength is a tuning choice selected with validation
data, not with the final test sample.

Bias and variance describe an important tradeoff. A rigid model may
systematically miss a pattern; a very flexible model may fit noise
specific to the training sample. More data can help, but only if it
represents the deployment problem and adds relevant information. A
million duplicated or highly dependent rows can provide less useful
information than their count suggests.

\hypertarget{M07-logistic-regression-trees-and-ensembles}{%
\section{Logistic Regression, Trees, and
Ensembles}\label{M07-logistic-regression-trees-and-ensembles}}

\noindent Logistic regression maps a linear index to a probability:

\[
\hat p_i=\frac{1}{1+\exp[-(\beta_0+x_i^\top\beta)]}.
\]

The coefficients describe changes in log odds under the model. A
coefficient's effect on probability depends on the starting probability
and other predictors. Predicted probabilities are often easier to
communicate than odds ratios, especially when the business decision uses
expected counts or contribution. Nonlinear features and interactions can
extend the model while preserving a relatively transparent structure.

A decision tree partitions the feature space into regions using
sequential splits. It can represent interactions without specifying them
manually: a spending threshold might matter only among recent customers.
Deep trees can fit small idiosyncratic groups, making pruning, minimum
leaf sizes, and validation important. A split point near an observed
value should not automatically be interpreted as a behavioral threshold.

Random forests average trees fitted with randomness in observations and
candidate features, reducing instability relative to a single tree in
many settings. Gradient boosting builds an ensemble sequentially to
improve a chosen loss, with tuning choices such as tree depth, learning
rate, and number of iterations. These methods can be effective on
tabular data, but performance remains empirical. More complexity also
increases the burden of checking leakage, calibration, computation, and
reproducibility.

Neural networks learn layered representations and can be especially
useful for text, images, or large complex datasets. Their flexibility
does not eliminate the need for baselines or careful evaluation. In a
modest customer table, a well-tuned simpler model may be competitive and
easier to maintain. Choose the method based on task performance and
practical constraints rather than the apparent modernity of the
algorithm.

\hypertarget{M07-validation-cross-validation-and-leakage}{%
\section{Validation, Cross-Validation, and
Leakage}\label{M07-validation-cross-validation-and-leakage}}

\noindent Use training data to fit parameters, validation data or cross-validation
to select models and tuning choices, and a final test sample to estimate
performance after those choices are fixed. Reusing the test sample
repeatedly turns it into a validation sample, even if the file is still
named ``test.'' Report the procedure so the claimed holdout remains
meaningful.

All learned preprocessing belongs inside the training process. This
includes imputation, scaling, feature selection, vocabulary
construction, dimensionality reduction, and target-based encodings. In
cross-validation, fit each transformation on the training fold and apply
it to that fold's validation observations. A pipeline makes the sequence
explicit. The
\href{https://scikit-learn.org/stable/common_pitfalls.html}{scikit-learn
pitfalls guide} documents implementation patterns for this separation.

Time-series and grouped data need adapted splits. Rolling-origin
evaluation trains on earlier periods and predicts later ones,
potentially repeating this across several origins. Grouped validation
holds out complete customers, households, stores, or other relevant
units. If a target window overlaps a training outcome window, a temporal
gap may be necessary. Document the dates and units assigned to each
split.

A useful audit asks when each feature became observable. In a
marketing-call dataset, call duration can be useful after a call but
unavailable when deciding whom to call. A model using that feature
answers a different prediction problem. The feature's legitimacy follows
the decision time, not whether a package accepts it or whether it
improves accuracy.

\hypertarget{M07-case-study-accuracy-can-conceal-a-useless-model}{%
\section{Case Study: Accuracy Can Conceal a Useless
Model}\label{M07-case-study-accuracy-can-conceal-a-useless-model}}

\noindent Suppose 5\% of 1,000 customers make a purchase in the target window. A
model predicting no purchase for everyone achieves 95\% accuracy while
identifying no purchasers. A different model flags 100 customers,
including 30 purchasers and 70 nonpurchasers. It has 30 true positives,
20 false negatives, 70 false positives, and 880 true negatives.

The second model's precision is 30\%, recall is 60\%, and accuracy is
91\%. Its lower accuracy coexists with a potentially useful ranking for
a limited-capacity task. Whether it improves profit depends on the
action, costs, and whether the purchase probabilities under that action
are the relevant quantities. Accuracy alone cannot make that decision.

ROC curves compare true-positive and false-positive rates across
thresholds. Precision-recall curves focus on the tradeoff between
finding positives and the composition of the flagged group, which can be
informative when positives are rare. Both depend on the evaluation
setting; precision is especially sensitive to prevalence. Lift at a
specified contact fraction compares the selected group's outcome rate
with the population rate, but observed lift is not automatically an
incremental treatment effect.

Report metrics with their operational units. For a service team that can
review only 100 cases per week, precision among the top 100 and the
number of important cases missed may be more relevant than a global
ranking metric. Use a holdout that reflects the weekly population and
inspect uncertainty in the metric when counts are small.

\hypertarget{M07-ranking-and-calibration-answer-different-questions}{%
\section{Ranking and Calibration Answer Different
Questions}\label{M07-ranking-and-calibration-answer-different-questions}}

\noindent A ranking model orders cases from more likely to less likely. A
calibrated probability model assigns probabilities that correspond to
observed frequencies in relevant groups. If customers assigned
probability 0.20 purchase about 20\% of the time in a suitable
evaluation population, that group is calibrated. Calibration can vary
across segments and over time even when overall performance looks
satisfactory.

A strictly increasing transformation of scores preserves their ranking.
Squaring probabilities between zero and one therefore leaves a
ranking-based AUC unchanged, apart from practical handling of ties or
finite precision, while generally changing calibration. A predicted
probability of 0.50 becomes 0.25. If the original 0.50 group purchases
half the time, the transformed probabilities understate its frequency.

\begin{figure}[htbp]
\centering
% Original synthetic teaching figure; see corresponding module.
\begin{tikzpicture}
\begin{axis}[xmin=0,xmax=1,ymin=0,ymax=1,xlabel={Reported purchase probability},ylabel={Implied observed purchase frequency},xtick={0,.25,.5,.75,1},ytick={0,.25,.5,.75,1},legend style={at={(.5,1.08)},anchor=south,legend columns=2},ymajorgrids=true]
\addplot[IlliniBlue,dashed,domain=0:1] {x};
\addplot[IlliniOrange,very thick,domain=0:1,samples=151] {sqrt(x)};
\legend{Calibrated score,Squared score}
\addplot[IlliniBlue,only marks,mark=*] coordinates {(.25,.5)};
\node[align=left,anchor=north west,font=\scriptsize] at (axis cs:.48,.25) {The original 0.50 becomes 0.25.\\Ranking is unchanged.};
\end{axis}
\end{tikzpicture}

\caption{A monotonic score transformation preserves ranking while changing probability calibration.}\label{fig:M07_calibration}
\end{figure}

For calibrated forecasts, predicted probabilities equal observed
frequencies. After squaring an initially calibrated score \(p\), the
reported probability is \(q=p^2\), while the observed frequency remains
\(p=\sqrt{q}\). Expected counts and profit calculations depend on the
probability scale, so unchanged ranking performance does not guarantee
unchanged decision quality.

Inspect calibration plots and proper scoring rules, such as log loss or
the Brier score, on held-out data. Calibration methods themselves
require fitting and should use data separate from the final evaluation.
If the deployment population changes, recalibration may be necessary. A
single global correction may be insufficient when the relationship
changes differently across groups.

\hypertarget{M07-thresholds-should-reflect-the-decision}{%
\section{Thresholds Should Reflect the
Decision}\label{M07-thresholds-should-reflect-the-decision}}

\noindent Consider a classification action in which a false positive costs
\(C_{FP}\) and a false negative costs \(C_{FN}\), with zero cost for
correct classifications. For a calibrated probability \(p\), predicting
positive has expected cost \(C_{FP}(1-p)\) and predicting negative has
expected cost \(C_{FN}p\). The cost-minimizing rule predicts positive
when

\[
p>\frac{C_{FP}}{C_{FP}+C_{FN}}.
\]

If an unnecessary review costs one unit and missing a true issue costs
nine units, the threshold is 0.10 under this simplified loss. A
threshold of 0.50 would encode a different cost balance. Capacity
constraints, review effectiveness, and heterogeneous costs require a
more detailed rule. The formula is useful because it makes the decision
assumptions explicit.

A marketing intervention usually needs a causal comparison beyond this
classification loss. For a coupon, let \(p_i(1)\) and \(p_i(0)\) be
purchase probabilities with and without the coupon, \(m\) the
contribution before the discount, \(d\) the discount paid on a treated
purchase, and \(c\) the contact cost. A simple incremental expected
contribution is

\[
\Delta\pi_i=m\{p_i(1)-p_i(0)\}-d\,p_i(1)-c.
\]

The discount is paid on all treated purchases in this setup, including
purchases that would have happened anyway. The equation assumes one
relevant purchase opportunity and a fixed discount and contribution. It
can be extended to spending, multiple purchases, returns, and
heterogeneous margins.

\hypertarget{M07-case-study-targeting-the-most-responsive-customers-can-lose-money}{%
\section{Case Study: Targeting the Most Responsive Customers Can Lose
Money}\label{M07-case-study-targeting-the-most-responsive-customers-can-lose-money}}

\noindent Take \(m=40\), \(d=5\), and \(c=2\). Customer A has purchase
probabilities 0.80 with the coupon and 0.78 without it. Customer B has
probabilities 0.30 and 0.15. A has the higher treated response
probability, but incremental contribution is \$-5.20 for A and \$2.50
for B.

\begin{table}[H]
\centering\small
\begin{tabular}{@{}
  >{\raggedright\arraybackslash}p{(\columnwidth - 6\tabcolsep) * \real{0.20}}
  >{\raggedleft\arraybackslash}p{(\columnwidth - 6\tabcolsep) * \real{0.26}}
  >{\raggedleft\arraybackslash}p{(\columnwidth - 6\tabcolsep) * \real{0.26}}
  >{\raggedleft\arraybackslash}p{(\columnwidth - 6\tabcolsep) * \real{0.28}}@{}}
\toprule
Customer & Purchase With Coupon & Purchase Without Coupon & Incremental
Contribution \\
\midrule

A & 0.80 & 0.78 & -\$5.20 \\
B & 0.30 & 0.15 & \$2.50 \\
\bottomrule
\end{tabular}
\end{table}

\begin{figure}[htbp]
\centering
% Original synthetic teaching figure; see corresponding module.
\begin{tikzpicture}
\begin{axis}[ybar,bar width=15pt,xmin=.4,xmax=2.6,ymin=-7.5,ymax=8.5,xtick={1,2},xticklabels={Customer A,Customer B},ylabel={Contribution per customer (dollars)},ymajorgrids=true,legend style={at={(.5,1.13)},anchor=south,legend columns=3},nodes near coords,nodes near coords style={font=\scriptsize}]
\addplot[fill=IlliniBlue,draw=IlliniBlue] coordinates {(1,.8)(2,6)};
\addplot[fill=LightStorm,draw=Storm] coordinates {(1,-6)(2,-3.5)};
\addplot[fill=IlliniOrange,draw=IlliniOrange] coordinates {(1,-5.2)(2,2.5)};
\legend{Incremental purchases,Discount plus contact,Net contribution}
\end{axis}
\end{tikzpicture}

\caption{Baseline purchasing and discount expense can reverse a response-based targeting rule.}\label{fig:M07_targeting}
\end{figure}

Incremental purchasing contributes \$0.80 for A and \$6.00 for B, while
discount and contact costs are \$6.00 and \$3.50, respectively. These
values follow from the assumed purchase probabilities. Estimating their
difference from customer data requires treatment variation and an
identification strategy. Uplift or heterogeneous-treatment-effect models
address that causal target under their conditions; they are not ordinary
response models with a different name.

Evaluate a targeting policy on data that support its value comparison. A
randomized holdout or a carefully designed policy evaluation can
estimate incremental contribution for the selected group. Estimating a
model and reporting its value on the same observations used to choose
the policy can overstate performance. Budget constraints also matter: a
customer with positive expected value may still be excluded when another
customer offers greater value per unit of scarce capacity.

\hypertarget{M07-interpretation-fairness-and-deployment}{%
\section{Interpretation, Fairness, and
Deployment}\label{M07-interpretation-fairness-and-deployment}}

\noindent Feature importance describes a relationship between a fitted model and
its inputs under a chosen method. Permutation importance measures the
change in predictive performance after disrupting a feature, with
complications when predictors are correlated. Partial-dependence plots
average predictions over a reference distribution and can involve
unrealistic combinations of correlated features. Local explanation
methods describe the fitted prediction under their conventions. These
tools do not automatically identify causal effects of changing a
feature.

Evaluate performance across groups relevant to the application and
examine how errors translate into service or opportunity. A model can be
accurate on average while performing poorly for new customers, small
stores, or a language group with sparse training data. Protected or
sensitive attributes and proxies require appropriate institutional and
legal review for the intended use. The consequences of errors depend on
the decision and the affected customers.

Deployment requires monitoring the input distribution, missingness,
outcome frequency, calibration, and decision value. A change in model
accuracy can reflect changing customer behavior, instrumentation, or
policy. Record versions and retain a reproducible baseline. If an
intervention changes the behavior that generates future training data,
the model's environment becomes endogenous to its use, because its
recommendations change subsequent exposure and behavior.

\hypertarget{M07-recommendation-models-and-the-candidate-set}{%
\section{Recommendation Models and the Candidate
Set}\label{M07-recommendation-models-and-the-candidate-set}}

\noindent A recommendation decision chooses a small set of products for a
particular customer at a particular time. Candidate generation first
identifies eligible items, scoring predicts their relevance, and ranking
selects their order. A final adjustment may impose stock availability,
repetition limits, category variety, or a customer-stated goal. An
accurate score cannot recover a useful product that candidate generation
has already excluded. Evaluation should therefore record both the
eligible set and the final list.

Content-based models use item attributes and the customer's history. A
grocery system might represent foods by ingredients, category, price, or
an embedding of their descriptions. Collaborative models instead use
patterns across customer-item interactions. Matrix factorization
represents a customer and an item with short vectors and scores their
compatibility through an inner product:

\[
\hat r_{uj}=\mu+b_u+b_j+p_u^\top q_j.
\]

Here \(u\) indexes a customer, \(j\) an item, \(\mu\) is an overall
level, \(b_u\) and \(b_j\) are customer and item offsets, and the
vectors \(p_u\) and \(q_j\) summarize interaction patterns in a common
latent space. The score's units depend on the fitted target. With
squared-error training on star ratings, it approximates rating units. A
ranking score fitted to implicit interactions has no automatic
probability interpretation. Popularity offsets capture broad appeal,
while the inner product permits customers with different histories to
prefer different products.

Ratings are explicit feedback. Clicks, views, and purchases are implicit
feedback whose meaning depends on exposure and opportunity. A missing
interaction is not a verified rejection. Treating every missing
customer-item pair as a negative can encode the previous recommender's
display policy. Negative sampling changes the training distribution, so
the sampling rule and any probability correction belong in the analysis.
Logging impressions and eligible alternatives makes these distinctions
more tractable.

New customers and new products create a cold-start problem. Content
features, initial preference questions, and popularity baselines can
provide information before interaction histories become informative. A
hybrid model combines these sources. Sequential models additionally
condition on the order of past events, allowing a recent purchase to
change what is useful next. The validation sample must preserve that
order and exclude later purchases, ratings, and catalog information from
the forecast origin.

\hypertarget{M07-ranking-metrics-depend-on-position-and-available-labels}{%
\section{Ranking Metrics Depend on Position and Available
Labels}\label{M07-ranking-metrics-depend-on-position-and-available-labels}}

\noindent Precision at \(K\) is the fraction of the first \(K\) recommendations
labeled relevant. Recall at \(K\) is the fraction of all labeled
relevant items retrieved in that list. Hit rate records whether at least
one relevant item appears. These metrics use different denominators and
can disagree. A hit rate of one can coexist with low precision when a
long list contains one relevant item.

Normalized discounted cumulative gain gives greater weight to relevant
items near the top. For nonnegative relevance labels \(r_1,\ldots,r_K\),
one convention is

\[
\mathrm{DCG}@K=\sum_{k=1}^{K}\frac{2^{r_k}-1}{\log_2(k+1)},
\qquad
\mathrm{NDCG}@K=\frac{\mathrm{DCG}@K}{\mathrm{IDCG}@K}.
\]

The position discount reflects the greater value assigned to placing
useful items earlier. The ideal DCG uses the best ordering of the
labeled eligible items at the same cutoff. Report the gain convention
and the treatment of customers with no relevant labels, for whom this
ratio has a zero denominator. Averaging only over customers with
positive labels changes the evaluation population.

Consider a synthetic three-item candidate set in which A is irrelevant
and B and C are relevant. The lists A, B, C and B, C, A both have
precision \(2/3\) and recall one. The first list has DCG
\(1/\log_2(3)+1/\log_2(4)\approx1.131\). The second has DCG
\(1+1/\log_2(3)\approx1.631\), which is also the ideal value. Their NDCG
values are approximately 0.693 and 1.000. The same number of relevant
items can produce different ranking quality because their positions
differ.

Historical labels describe items observed under earlier exposure and
consumption decisions. High NDCG on those labels can coexist with poor
discovery of items rarely displayed. Evaluate a common candidate set,
retain customer and time boundaries, compare with a popularity baseline,
and report performance for customers and products with little history.
The policy-value calculation in M6, Section 6.7, addresses selected
exposure under different assumptions from an ordinary held-out ranking
calculation. Neither calculation alone establishes customer satisfaction
or the value of changing the deployed system.

\hypertarget{M07-case-study-food-recommendations-variety-and-customer-goals}{%
\section{Case Study: Food Recommendations, Variety, and Customer
Goals}\label{M07-case-study-food-recommendations-variety-and-customer-goals}}

\noindent A food app can predict what a customer usually eats while learning
little about what would help that customer pursue a stated goal. The
distinction changes the training sample, the model inputs, and the
relevant evaluation outcome.
\href{https://doi.org/10.1145/3485447.3514193}{Ling, Nie, Nielsen,
Knäuper, Yang, and Dubé (2022)} develop a recommendation system using
daily food histories and information about calories and body-mass-index
groups. Their 2016 app sample restricts users to ages 18--65 with
substantial logging histories and removes sparse daily records. Table 2
reports 7,457 users in the resulting model sample.

The model represents both the foods within a day and the sequence of
days. Separate models learn from general histories, particular
body-mass-index groups, and periods associated with the relevant weight
outcome. Their predictions can be combined. The last basket supplies the
test observation, the preceding basket supplies validation, and earlier
baskets supply training. This design evaluates whether the model
predicts foods eaten during periods associated with the intended
outcome. It does not randomize customers to receive the recommendations,
so the reported ranking improvements do not identify a causal effect on
weight.

Selecting successful periods can help define a target pattern while also
selecting motivation, physical activity, reporting behavior, or other
determinants of the outcome. A customer who follows a similar food
pattern may differ on those other dimensions. For deployment, the goal
and all required inputs must be available when recommendations are made.
A later realized weight outcome can define a retrospective evaluation
stratum; it cannot serve as a known input for a new customer unless
supplied as an explicit scenario rather than an observed future fact.

\href{https://doi.org/10.1086/720453}{Nielsen, Yang, Dubé, Knäuper,
Ling, and Nie (2022)} examine variety in food logs and develop FOODVAR.
Their empirical analysis relates daily variety to subsequent changes in
calories and the composition of food consumption. Variety measures
include item counts and a reverse-coded concentration index based on
calorie shares. Two days with the same item count can therefore differ
in variety if one is dominated by a single food. The reported
relationships are conditional associations, and the denominators and
sample restrictions differ across the total-calorie, vegetable, and
snack analyses.

FOODVAR incorporates a variety-change indicator into sequential food
recommendation. Its evaluation uses 30 days of history, seven subsequent
validation days, and a final test day, with the indicator representing
the intended variety direction. Table 5 compares hit-rate and NDCG
measures against the baseline across weight-outcome groups. The results
establish predictive differences within that evaluation design. A
prospective test would elicit the variety goal before displaying
recommendations and measure subsequent behavior. The synthetic
category-bonus example in M8, Section 8.2, is a simpler list-selection
exercise with a different algorithm.

\hypertarget{M07-a-goal-weight-can-change-the-preferred-item-abruptly}{%
\subsection{A Goal Weight Can Change the Preferred Item
Abruptly}\label{M07-a-goal-weight-can-change-the-preferred-item-abruptly}}

\noindent Suppose three products have preference scores and goal-compatibility
scores on a common illustrative scale. A scores \((0.90,0.10)\), B
scores \((0.75,0.60)\), and C scores \((0.60,0.90)\). The scores are
imposed inputs, not estimated purchase probabilities or health effects.
A weighted ranking rule is

\[
s_j(\alpha)=(1-\alpha)p_j+\alpha g_j,\qquad 0\leq\alpha\leq1,
\]

where \(p_j\) measures preference compatibility, \(g_j\) measures
compatibility with the stated goal, and \(\alpha\) determines their
relative importance. At \(\alpha=0\), A ranks first. At \(\alpha=0.25\),
the scores are 0.700, 0.7125, and 0.675, so B ranks first. At
\(\alpha=0.50\), C ranks first with 0.750. Solving the pairwise score
equalities gives the optimal switches at \(\alpha=3/13\) and
\(\alpha=1/3\).

\begin{figure}[htbp]
\centering
\begin{tikzpicture}
\begin{axis}[width=13cm,height=7.4cm,xmin=0,xmax=1,ymin=0,ymax=1.12,
  xtick={0,.25,.5,.75,1},ytick={0,.25,.5,.75,1},
  xlabel={Goal weight $\alpha$},ylabel={Combined score},
  legend style={at={(.5,-.21)},anchor=north,legend columns=3}]
\addplot[IlliniBlue,very thick,domain=0:1,samples=2]{.9-.8*x};
\addplot[Patina,very thick,dashed,domain=0:1,samples=2]{.75-.15*x};
\addplot[IlliniOrange,very thick,dashdotted,domain=0:1,samples=2]{.6+.3*x};
\legend{Item A,Item B,Item C}
\draw[Storm,densely dotted] (axis cs:.230769,0)--(axis cs:.230769,1.03);
\draw[Storm,densely dotted] (axis cs:.333333,0)--(axis cs:.333333,1.03);
\node[font=\footnotesize,anchor=south] at (axis cs:.1154,1.02) {A first};
\node[font=\footnotesize,anchor=south] at (axis cs:.2821,1.02) {B};
\node[font=\footnotesize,anchor=south] at (axis cs:.6667,1.02) {C first};
\fill[IlliniBlue] (axis cs:.230769,.715385) circle (2pt);
\fill[IlliniOrange] (axis cs:.333333,.7) circle (2pt);
\end{axis}
\end{tikzpicture}

\caption{Smooth changes in the objective weight produce discrete changes in the top-ranked item.}\label{fig:M07_goal_blend}
\end{figure}

Every score changes continuously, yet the chosen item changes
discretely. A small adjustment near a crossing can therefore replace the
displayed recommendation. B is the preferred compromise only over a
narrow weight interval. A useful audit varies the weight and inspects
which customers experience large changes in their lists. The chosen
scales and weights define the tradeoff; combining scores does not itself
establish that a goal will be achieved.

\hypertarget{M07-llm-explanations-need-factual-and-behavioral-evaluation}{%
\section{LLM Explanations Need Factual and Behavioral
Evaluation}\label{M07-llm-explanations-need-factual-and-behavioral-evaluation}}

\noindent A shopping assistant may generate a persuasive reason for recommending a
product after the product has already been selected. That text can be
useful, but its fluency does not establish that it describes the ranking
model's actual reason. Three questions remain distinct: whether the
explanation contains correct product facts, whether it corresponds to
the system's decision process, and whether it helps the customer make a
satisfactory choice.

\href{https://arxiv.org/abs/2502.16759v3}{Wang, Li, and Chen's RecPIE
study}, in its May 2026 version, jointly trains recommendation
prediction and natural-language explanation. Prediction performance
informs explanation generation, and explanation representations feed
into subsequent predictions. The empirical setting includes Google Maps
point-of-interest recommendations and a separate human assessment of
explanations. Reported improvements in ranking and explanation
preference concern those evaluations. They do not by themselves estimate
the effect of deploying the system on purchases, retention, or customer
welfare.

\href{https://aclanthology.org/2024.findings-emnlp.247/}{Qin, Huang,
Deng, Lei, and Chua (2024)} study credibility in conversational
recommendation. Their PC-CRS approach uses credibility-aware persuasive
strategies and subsequent revision of explanations. The comparison
motivates separate evaluation of persuasive success and support for
factual claims. An assistant can increase acceptance by asserting an
appealing attribute that the product does not possess. Acceptance alone
would reward that error.

Consider a synthetic catalog with prices, delivery times, and verified
product attributes. A customer asks for delivery within two days at a
price below \$30. First enforce those eligibility constraints using the
catalog. Then rank eligible products and generate explanations that
refer to specific supported attributes. ``Arrives tomorrow'' is
unsupported when the recorded delivery window is three to five days,
even if the rest of the explanation is relevant. ``You value quiet
operation'' is unsupported when the customer has expressed no such
preference and no permitted history supplies it.

An audit should inspect factual support, constraint violations, ranking
performance, and customer outcomes separately. Remove or alter a cited
attribute and examine whether the explanation and ranking change as
claimed; this provides a behavioral diagnostic of correspondence, with
scope limited to the perturbations tested. Hold out customers, products,
or time periods according to deployment. Keep the catalog version and
the evidence used for each explanation so the claim can be checked after
product information changes.

\hypertarget{M07-ai-assisted-workflow-practice-and-readings}{%
\section{AI-Assisted Workflow, Practice, and
Readings}\label{M07-ai-assisted-workflow-practice-and-readings}}

\noindent Ask an assistant to create an evaluation plan before choosing a model.
Specify forecast origin, horizon, unit, grouping, baseline, loss, and
final holdout. Require a feature-availability table and a pipeline that
fits transformations only on training data. Have it produce both a
metric table and an explanation of the operational decision each metric
supports.

For practice, verify the 91\% accuracy and 60\% recall in the
rare-purchase example. Square a small set of predicted probabilities and
confirm that their order is unchanged while expected purchase counts
change. Recalculate the coupon case with a \$1 discount and explain
which customer becomes more attractive. Finally, identify why a model
trained on all customers' latest outcomes cannot honestly claim to
predict those same outcomes out of sample.

For a recommender exercise, have an assistant build a popularity
baseline, a simple content score, and the stated goal-weight rule on
synthetic items. Require a chronological split, a documented candidate
set, and a report of excluded or unavailable products. Verify the two
NDCG calculations by hand, then enumerate the goal-weight switches
independently. Ask why an unseen item was labeled irrelevant and require
the code to distinguish an observed rejection from absent exposure. A
working implementation of the simple score is not a replication of
FOODVAR, HHFR, or RecPIE.

For an explanation exercise, give the assistant a small catalog and a
customer's explicit constraints. Request a recommendation with an
attribute-level source for each factual claim. Include one unavailable
attribute and one product that violates a constraint. Check whether the
system acknowledges missing information and excludes the ineligible
product. Score factual support separately from how persuasive the
explanation sounds.

\href{https://www.statlearning.com/}{An Introduction to Statistical
Learning} develops the central learning methods. The
\href{https://scikit-learn.org/stable/modules/model_evaluation.html}{scikit-learn
evaluation guide} provides metric definitions and implementation
references. \href{https://doi.org/10.1007/s11002-020-09535-7}{Hagen and
colleagues (2020)} connect machine-learning tools to behavioral
marketing questions.

\hypertarget{M07-mhealth-predicting-a-returning-users-next-basket}{%
\section{mHealth: Predicting a Returning User's Next
Basket}\label{M07-mhealth-predicting-a-returning-users-next-basket}}

\noindent The manager needs a recommendation baseline that uses only information
available when suggestions are chosen. Global category popularity is
easy to implement. A personal frequency rule can exploit repeated food
choices without estimating a complex model. The comparison should
preserve complete baskets and evaluate both rules on the same users,
candidate set, and future outcomes.

A basket is the set of distinct source categories recorded by a user on
a day; missing category labels are excluded. Freeze the candidate set at
the 312 categories observed before July 1. For each forecast origin,
require at least 30 prior observed basket days and select the first
observed basket during the following month as the target. Validation
starts July 1 and testing starts September 1. Every score uses dates
strictly before its forecast origin. The test concerns returning users
who subsequently log a basket, conditional on that observation; it does
not evaluate people who never return.

Personal frequency is the fraction of a user's observed days in a
specified history window containing a category. It does not divide by
all calendar days. A tiny popularity term resolves most frequency ties;
exact remaining ties use alphabetical order. If a user has no observed
day in a short window, personal frequencies are zero and the popularity
term supplies the order. Global popularity uses the corresponding
fraction across all prior observed user-days. Scores produce a list of
ten distinct categories.

Choose among 30-day, 90-day, and all-prior-history windows using July
mean NDCG at ten. Their validation scores are 0.5666, 0.5661, and
0.5501, respectively, among 297 evaluated users. The 30-day rule wins by
a very small margin over 90 days. The code preserves the specified tie
order and carries the chosen window to September. A further claim that
30 days is generally superior would need repeated or independent
validation.

\begin{table}[H]
\centering\small
\begin{tabular}{@{}lrrr@{}}
\toprule
September Rule & Evaluated Users & Recall at 10 & NDCG at 10 \\
\midrule

Global Popularity & 299 & 0.2799 & 0.3270 \\
Personal 30-Day Frequency & 299 & 0.4759 & 0.5862 \\
\bottomrule
\end{tabular}
\end{table}

\begin{figure}[htbp]
\centering
\begin{tikzpicture}[x=1cm,y=1cm]
\draw[edge] (0,0)--(13,0);
\foreach \x/\label in {0/January 1,6/July 1,8/September 1,12/December 31}{\draw[IlliniBlue] (\x,.12)--(\x,-.12);\node[anchor=north,font=\footnotesize] at (\x,-.2) {\label};}
\draw[IlliniBlue,very thick] (0,1)--(5.9,1);
\node[anchor=south,font=\footnotesize] at (3,1.1) {Prior History; Freeze Candidates};
\node[orangebox,text width=3.4cm] at (6,3.0) {July Validation\\297 Target Baskets\\Choose History Window};
\draw[accent] (6,2.1)--(6,.2);
\node[box,text width=3.6cm] at (10,3.0) {September Test\\299 Target Baskets\\Report Frozen Rule};
\draw[edge] (9.4,2.1)--(8,.2);
\node[align=center,text width=12cm] at (6,-1.5) {All features precede each origin.\\Each target is the first observed basket within that origin's month.};
\end{tikzpicture}

\caption{The candidate universe is fixed before validation. The September rule can use updated pre-September histories, while its chosen history-window length remains fixed from July.}\label{fig:M07_mhealth_timing}
\end{figure}

There are 362 eligible users at the September origin, but only 299 have
a target basket in that month. Reporting both counts exposes the
conditioning introduced by the target definition. All observed target
categories in this evaluation are covered by the frozen candidate set.
The implementation nevertheless retains all target categories in the
recall denominator and ideal DCG, so an unsupported category would count
against performance.

The personal rule recovers more recorded categories than the popularity
baseline. This supports personalization for predicting the observed
basket under the historical environment. It does not show that
displaying those suggestions changes choices, improves nutrition, or
generates incremental app value. Logging may omit foods, and the basket
can occur several days after the fixed origin. A same-day or event-time
recommender would need a corresponding evaluation design.

\hypertarget{M07-histories-candidate-sets-and-ranking-metrics}{%
\subsection{Histories, Candidate Sets, and Ranking
Metrics}\label{M07-histories-candidate-sets-and-ranking-metrics}}

\noindent \texttt{ranking.py} defines the common evaluation so that M8 changes the
ranking objective while holding the target baskets fixed. Recall divides
the number of recovered categories by the target basket size. NDCG
discounts relevant categories by position and divides by the best
possible ranking for that target. Each evaluated user contributes one
metric value, and the final result averages users equally.

\FloatBarrier
\Needspace{8\baselineskip}
Python source:
\href{../03_Data_and_Code/mHealth/recurring_examples/ranking.py}{ranking.py}.
From its directory, run \texttt{python\ ranking.py}.

\begin{Verbatim}[breaklines,breakanywhere,fontsize=\small]
"""Whole-basket chronological evaluation with a frozen category universe."""

import numpy as np
import pandas as pd


def candidate_categories(baskets):
    prior = baskets.loc[baskets.date.lt("2016-07-01")]
    return sorted({category for basket in prior.basket for category in basket})


def evaluation_cases(baskets, cutoff, candidates):
    cutoff = pd.Timestamp(cutoff)
    prior = baskets.loc[baskets.date.lt(cutoff)]
    counts = prior.groupby("Userid").size()
    eligible = counts[counts >= 30].index
    end = cutoff + pd.offsets.MonthBegin(1)
    target = baskets.loc[baskets.date.ge(cutoff) & baskets.date.lt(end)]
    target = target.loc[target.Userid.isin(eligible)].sort_values(["Userid", "date"])
    target = target.drop_duplicates("Userid")
    histories = dict(tuple(prior.groupby("Userid")))
    assert prior.date.max() < cutoff
    global_counts = (
        prior.basket.explode().value_counts().reindex(candidates, fill_value=0)
    )
    popularity = global_counts.to_numpy(dtype=float) / len(prior)
    cases = [(histories[row.Userid], set(row.basket)) for row in target.itertuples()]
    return cases, popularity, len(eligible)


def personal_frequency(history, cutoff, candidates, window):
    if window is not None:
        start = pd.Timestamp(cutoff) - pd.Timedelta(days=window)
        history = history.loc[history.date.ge(start)]
    counts = history.basket.explode().value_counts().reindex(candidates, fill_value=0)
    return counts.to_numpy(dtype=float) / max(1, len(history))


def rank_categories(scores, candidates, k=10):
    # Stable sorting resolves exact score ties using alphabetically sorted candidates.
    positions = np.argsort(-scores, kind="stable")[:k]
    return [candidates[index] for index in positions]


def ranking_metrics(ranking, target):
    relevance = np.array([category in target for category in ranking], dtype=float)
    discounts = 1 / np.log2(np.arange(2, len(ranking) + 2))
    ideal = discounts[: min(len(target), len(ranking))].sum()
    return float(relevance.sum() / len(target)), float(
        (relevance * discounts).sum() / ideal
    )


def evaluate(
    baskets, cutoff, candidates, window=None, popularity_only=False, novelty_weight=0
):
    cases, popularity, eligible = evaluation_cases(baskets, cutoff, candidates)
    rows = []
    for history, target in cases:
        frequency = personal_frequency(history, cutoff, candidates, window)
        recent = personal_frequency(history, cutoff, candidates, 30)
        unseen = recent == 0
        scores = popularity if popularity_only else frequency + 1e-6 * popularity
        ranking = rank_categories(scores + novelty_weight * unseen, candidates)
        recall, ndcg = ranking_metrics(ranking, target)
        recent_set = {candidates[j] for j in np.flatnonzero(recent)}
        novelty = sum(category not in recent_set for category in ranking) / len(ranking)
        coverage = len(target & set(candidates)) / len(target)
        rows.append([recall, ndcg, novelty, coverage])
    means = np.mean(rows, axis=0)
    return {
        "eligible_users": eligible,
        "evaluated_users": len(cases),
        "recall_at_10": float(means[0]),
        "ndcg_at_10": float(means[1]),
        "novelty_at_10": float(means[2]),
        "target_candidate_coverage": float(means[3]),
    }
\end{Verbatim}

\hypertarget{M07-validation-before-testing}{%
\subsection{Validation Before
Testing}\label{M07-validation-before-testing}}

\noindent The execution order matters. Fit the history-window choice using July,
freeze it, and report September performance without changing the rule
after seeing the test. The small metric counterexample confirms that
placing both relevant categories first yields NDCG of one, while placing
an irrelevant category first lowers it. A separate test supplies a
category that occurs only in the future and verifies that it never
enters personal history.

\FloatBarrier
\Needspace{8\baselineskip}
Python source:
\href{../03_Data_and_Code/mHealth/recurring_examples/m07.py}{m07.py}.
From its directory, run \texttt{python\ m07.py}.

\begin{Verbatim}[breaklines,breakanywhere,fontsize=\small]
"""Choose history length in July and evaluate the frozen rule in September."""

import json
from common import load_inputs, save_result
from ranking import candidate_categories, evaluate, ranking_metrics

food, daily, baskets = load_inputs()
candidates = candidate_categories(baskets)
windows = {"30_days": 30, "90_days": 90, "all_prior_days": None}
validation = {
    name: evaluate(baskets, "2016-07-01", candidates, window)
    for name, window in windows.items()
}
# Ties retain the first listed specification; September plays no role in selection.
chosen = max(validation, key=lambda name: validation[name]["ndcg_at_10"])
test = {
    "popularity": evaluate(baskets, "2016-09-01", candidates, popularity_only=True),
    "personal_frequency": evaluate(baskets, "2016-09-01", candidates, windows[chosen]),
}
assert ranking_metrics(["B", "C", "A"], {"B", "C"}) == (1.0, 1.0)
assert ranking_metrics(["A", "B", "C"], {"B", "C"})[1] < 1
result = {
    "candidate_categories": len(candidates),
    "validation": validation,
    "chosen_window": chosen,
    "test": test,
}
save_result("M07", result)
print(json.dumps(result, indent=2))
\end{Verbatim}

\hypertarget{M07-python-calculations-for-the-synthetic-cases}{%
\section{Python Calculations for the Synthetic
Cases}\label{M07-python-calculations-for-the-synthetic-cases}}

\noindent The inputs below are imposed teaching values. Running the program
creates the labeled synthetic CSV files and prints the corresponding
calculations. They are separate from the mHealth estimates. Keep the
script beside \texttt{example\_helpers.py} in
\texttt{03\_Data\_and\_Code/Lecture\_Examples}.

\FloatBarrier
\Needspace{8\baselineskip}
Python source:
\href{../03_Data_and_Code/Lecture_Examples/m07_synthetic.py}{m07\_synthetic.py}.
From its directory, run \texttt{python\ m07\_synthetic.py}.

\begin{Verbatim}[breaklines,breakanywhere,fontsize=\small]
"""Synthetic numerical examples for M7; run from any working directory."""

import itertools
import json
import math
from example_helpers import write_csv, ols_slope, shares

r = {}
r["M07"] = {
    "precision": 30 / 100,
    "recall": 30 / 50,
    "accuracy": (30 + 880) / 1000,
    "all_negative_accuracy": 950 / 1000,
    "cost_threshold": 1 / (1 + 9),
    "coupon_A": 40 * (0.8 - 0.78) - 5 * 0.8 - 2,
    "coupon_B": 40 * (0.3 - 0.15) - 5 * 0.3 - 2,
}
write_csv(
    "M07_coupon_synthetic.csv",
    [
        "customer",
        "purchase_with_coupon",
        "purchase_without_coupon",
        "margin_dollars",
        "discount_dollars",
        "contact_cost_dollars",
    ],
    [("A", 0.8, 0.78, 40, 5, 2), ("B", 0.3, 0.15, 40, 5, 2)],
)

rankings = {
    "ABC": [("A", 0), ("B", 1), ("C", 1)],
    "BCA": [("B", 1), ("C", 1), ("A", 0)],
}


def dcg(labels):
    return sum(
        (2**label - 1) / math.log2(rank + 1)
        for rank, label in enumerate(labels, start=1)
    )


ideal = dcg([1, 1, 0])
r["M07"]["rankings"] = {
    name: {
        "precision_at_3": sum(v for _, v in vals) / 3,
        "recall_at_3": sum(v for _, v in vals) / 2,
        "dcg_at_3": dcg([v for _, v in vals]),
        "ndcg_at_3": dcg([v for _, v in vals]) / ideal,
    }
    for name, vals in rankings.items()
}
write_csv(
    "M07_rankings_synthetic.csv",
    ["ranking", "position", "item", "binary_relevance"],
    [
        (name, pos, item, rel)
        for name, vals in rankings.items()
        for pos, (item, rel) in enumerate(vals, start=1)
    ],
)
assert math.isclose(r["M07"]["rankings"]["BCA"]["ndcg_at_3"], 1)
assert 0.693 < r["M07"]["rankings"]["ABC"]["ndcg_at_3"] < 0.694

goal_items = {"A": (0.9, 0.1), "B": (0.75, 0.6), "C": (0.6, 0.9)}
goal_score = (
    lambda item, alpha: (1 - alpha) * goal_items[item][0] + alpha * goal_items[item][1]
)
r["M07"]["goal_blend"] = {
    "scores": {
        str(alpha): {item: goal_score(item, alpha) for item in goal_items}
        for alpha in [0, 0.25, 0.5, 1]
    },
    "optimal_switches": [3 / 13, 1 / 3],
}
write_csv(
    "M07_goals_synthetic.csv",
    ["item", "preference_score", "goal_score"],
    [(item, *scores) for item, scores in goal_items.items()],
)
assert math.isclose(goal_score("A", 3 / 13), goal_score("B", 3 / 13))
assert math.isclose(goal_score("B", 1 / 3), goal_score("C", 1 / 3))
for i in range(1001):
    alpha = i / 1000
    expected = "A" if alpha < 3 / 13 else ("B" if alpha < 1 / 3 else "C")
    assert max(goal_items, key=lambda j: goal_score(j, alpha)) == expected
print(json.dumps(r, indent=2))
\end{Verbatim}

\hypertarget{M07-mhealth-basket-associations-and-future-validation}{%
\section{mHealth: Basket Associations and Future
Validation}\label{M07-mhealth-basket-associations-and-future-validation}}

\noindent A recommendation manager wants to identify categories that commonly
appear together in a recorded day. An association rule can generate
candidates for a basket-completion feature. Its commercial value still
depends on when a suggestion is shown, whether the customer has already
recorded the item, and whether the suggestion changes behavior. A high
co-occurrence rate alone does not establish incremental demand.

Treat each observed user-day with at least one classified category as a
basket. Deduplicate categories within the day, so repeated servings
contribute one presence indicator. For categories \(A\) and \(B\),
define

\[
\operatorname{support}(A,B)=P(A\cap B),\qquad
\operatorname{confidence}(A\to B)=P(B\mid A),\qquad
\operatorname{lift}(A\to B)=\frac{P(B\mid A)}{P(B)}.
\]

Probabilities are empirical fractions of recorded baskets. Confidence is
directional, whereas lift is symmetric for a two-category pair. A
popular category can have high confidence without an unusually strong
association. Lift compares the conditional rate with the marginal rate
and therefore supplies the necessary benchmark. The label ``confidence''
here is a conditional frequency, separate from a statistical confidence
interval.

Use January--June data to select the 20 most frequent categories. Among
their pairs, require at least 100 joint baskets and 30 distinct
supporting users. Rank eligible pairs by training lift and fix the
largest before examining July--August. These prespecified restrictions
avoid selecting a pair supported by only a few repeated entries, but
they do not make users' repeated baskets independent. The training
sample has 40,727 baskets from 358 users. The selected pair is oil and
tomato, supported jointly by 192 users.

\begin{table}[H]
\centering\small
\begin{tabular}{@{}lrr@{}}
\toprule
Metric & January--June & July--August \\
\midrule

Recorded Baskets & 40,727 & 12,956 \\
Baskets With Oil & 7,283 & 2,174 \\
Baskets With Tomato & 4,708 & 1,930 \\
Baskets With Both & 1,985 & 658 \\
Tomato Given Oil & 27.26\% & 30.27\% \\
Marginal Tomato Rate & 11.56\% & 14.90\% \\
Lift & 2.358 & 2.032 \\
\bottomrule
\end{tabular}
\end{table}

The later evaluation includes observed baskets from the previously known
users; 334 of them contribute baskets in July--August. Users with no
later basket supply no basket-level outcome. In the table, conditional
tomato frequency rises while lift falls because tomato also becomes more
common overall. Judging only the conditional frequency would miss that
change in the comparison benchmark.

\begin{figure}[htbp]
\centering
\begin{tikzpicture}
\begin{axis}[xbar,bar width=10pt,width=11.5cm,height=5.7cm,xlabel={Recorded Baskets With Tomato (Percent)},xmin=0,xmax=43,ymin=.4,ymax=2.6,ytick={1,2},yticklabels={July--August,January--June},xtick={0,10,20,30,40},xmajorgrids=true,legend style={at={(.5,1.17)},anchor=south,legend columns=2}]
\addplot[IlliniBlue,fill=IlliniBlue,bar shift=7pt] coordinates {(11.559899,2)(14.896573,1)};
\addlegendentry{All Recorded Baskets}
\addplot[IlliniOrange,fill=IlliniOrange,bar shift=-7pt] coordinates {(27.255252,2)(30.266789,1)};
\addlegendentry{Baskets With Oil}
\node[align=center,font=\footnotesize] at (axis cs:36.5,2.05) {Lift\\2.358};
\node[align=center,font=\footnotesize] at (axis cs:37.5,1.05) {Lift\\2.032};
\end{axis}\end{tikzpicture}

\caption{Conditional tomato frequency rises in the later period, but the marginal tomato rate rises proportionally more, reducing lift.}\label{fig:M07_basket_lift}
\end{figure}

A 500-draw bootstrap resamples the 358 training users, carrying each
user's complete later basket history into the draw, including zero later
baskets. The percentile interval for later lift is {[}1.762, 2.411{]}.
It describes sampling variation for the already selected pair, with
dependence within each user preserved. It does not address common
calendar shocks, uncertainty over the category-selection procedure, or
the effect of displaying a recommendation. The recorded category names
also do not establish nutritional complementarity.

The
\href{https://rasbt.github.io/mlxtend/user_guide/frequent_patterns/association_rules/}{mlxtend
association-rule guide} defines support, confidence, and lift and
provides scalable itemset-mining tools. The code below enumerates pairs
from a bounded candidate set so every numerator is inspectable. Its
results can be compared with that library after the same basket and
candidate definitions are imposed.

The
\href{https://github.com/benfred/implicit/tree/main/examples}{implicit
recommendation examples} and
\href{https://github.com/NicolasHug/Surprise/tree/master/examples}{Surprise
examples} address different feedback structures. Food logging supplies
observed occurrences; missing entries do not provide explicit dislike
ratings. For next-basket prediction, retain the chronological evaluation
and known candidate set from Section 7.15. For within-basket completion,
the source would additionally need an information rule specifying which
recorded categories are visible when a suggestion is made.

\FloatBarrier
\Needspace{8\baselineskip}
Python source:
\href{../03_Data_and_Code/mHealth/recurring_examples/association_rules.py}{association\_rules.py}.
From its directory, run \texttt{python\ association\_rules.py}.

\begin{Verbatim}[breaklines,breakanywhere,fontsize=\small]
"""Select a category pair in January-June; evaluate it in July-August."""

from itertools import combinations
import json
import numpy as np
import pandas as pd
from common import SEED, load_food, make_baskets, save_result


def metrics(n, a, b, both):
    if min(n, a, b) <= 0:
        raise ValueError("Positive basket and marginal counts are required")
    return {
        "baskets": int(n),
        "a_baskets": int(a),
        "b_baskets": int(b),
        "joint_baskets": int(both),
        "support": both / n,
        "confidence_a_to_b": both / a,
        "baseline_b": b / n,
        "lift": both * n / (a * b),
    }


def presence(baskets, categories):
    return np.array(
        [[c in basket for c in categories] for basket in baskets.basket], dtype=np.int64
    )


if __name__ == "__main__":
    baskets = make_baskets(load_food())
    train = baskets.loc[baskets.date < "2016-07-01"]
    frequency = train.explode("basket").groupby("basket").size()
    categories = sorted(
        frequency.sort_values(ascending=False, kind="stable").head(20).index
    )
    matrix = presence(train, categories)
    joint, marginal = matrix.T @ matrix, matrix.sum(axis=0)
    rules = []
    for a, b in combinations(range(len(categories)), 2):
        mask = (matrix[:, a] * matrix[:, b]).astype(bool)
        supporting_users = train.loc[mask, "Userid"].nunique()
        if joint[a, b] >= 100 and supporting_users >= 30:
            rule = metrics(len(train), marginal[a], marginal[b], joint[a, b])
            rules.append(
                {
                    "a": categories[a],
                    "b": categories[b],
                    "supporting_users": int(supporting_users),
                    **rule,
                }
            )
    rules.sort(key=lambda r: (-r["lift"], r["a"], r["b"]))
    selected = rules[0]
    users = sorted(train.Userid.unique())
    future = baskets.loc[
        baskets.Userid.isin(users) & baskets.date.between("2016-07-01", "2016-08-31")
    ].copy()
    binary = presence(future, [selected["a"], selected["b"]])
    future["a"], future["b"] = binary[:, 0], binary[:, 1]
    future["both"], future["n"] = binary.prod(axis=1), 1
    counts = (
        future.groupby("Userid")[["n", "a", "b", "both"]]
        .sum()
        .reindex(users, fill_value=0)
        .to_numpy()
    )
    evaluation = metrics(*counts.sum(axis=0))
    rng = np.random.default_rng(SEED)
    lifts = []
    for _ in range(500):
        total = counts[rng.integers(0, len(users), len(users))].sum(axis=0)
        if min(total[:3]) > 0:
            lifts.append(metrics(*total)["lift"])
    evaluation["user_bootstrap_ci95"] = np.quantile(lifts, [0.025, 0.975]).tolist()
    evaluation["bootstrap_valid_draws"] = len(lifts)
    result = {
        "candidate_categories": len(categories),
        "eligible_pairs": len(rules),
        "training_users": len(users),
        "evaluation_users": int(future.Userid.nunique()),
        "selected_training_pair": selected,
        "july_august": evaluation,
    }
    save_result("association_rules", result)
    print(json.dumps(result, indent=2))
\end{Verbatim}

\hypertarget{M08-m8.-current-topics-and-integrated-marketing-decisions}{%
\chapter{Current Topics and Integrated Marketing
Decisions}\label{M08-m8.-current-topics-and-integrated-marketing-decisions}}

\noindent A marketing team can now generate campaign copy, build a prediction
model, and summarize customer feedback within a short working session.
The remaining difficulty is deciding whether the resulting system
improves a business decision. A recommendation engine can reproduce past
exposure, a media model can attribute seasonal demand to advertising,
and an AI-generated analysis can calculate an invalid comparison
precisely. Evaluating these systems requires a clear account of their
information, incentives, assumptions, and consequences.

\hypertarget{M08-recommendation-systems-learn-from-selected-exposure}{%
\section{Recommendation Systems Learn From Selected
Exposure}\label{M08-recommendation-systems-learn-from-selected-exposure}}

\noindent A retailer wants to recommend products that customers will find useful.
A common starting point is past interaction: purchases, clicks, ratings,
or time spent. These observations combine preferences with the
opportunities customers received. A product displayed prominently has
more chances to be clicked than an otherwise similar product that was
rarely shown. An absent interaction can therefore mean dislike, lack of
exposure, lack of need, or an unavailable product.

Content-based methods use product attributes and customer histories to
identify similar items. Collaborative methods use patterns across users
and items, including matrix factorization, which approximates an
interaction matrix with lower-dimensional user and item representations.
Hybrid systems combine these sources. Each method needs a strategy for
new customers and new products that have little interaction history.
Popularity baselines and simple category rules provide useful
comparisons.

Recommendation systems also change the data used to train the next
model. A product that receives more exposure can accumulate more
interactions, become easier to estimate, and receive still more
exposure. This process can amplify initial differences without
establishing that the repeatedly displayed product would be preferred
under equal exposure.

\href{https://arxiv.org/abs/1710.11214}{Chaney, Stewart, and Engelhardt
(2018)} study this feedback through simulations in which recommendations
shape the next training sample. Their results show that consumption can
become more homogeneous without corresponding utility improvements. The
result depends on the simulated preferences, recommendation rules, and
updating process. It identifies a mechanism to examine in deployed
systems rather than a universal effect of personalization.

\begin{figure}[htbp]
\centering
% Original synthetic teaching figure; see corresponding module.
\begin{tikzpicture}
\node[box,text width=3cm] (m) at (0,0) {Model and ranking};
\node[box,text width=3cm] (e) at (5,0) {Product exposure};
\node[box,text width=3cm] (y) at (10,0) {Observed behavior};
\node[orangebox,text width=3cm] (p) at (10,2.4) {Customer preferences};
\node[box,text width=3cm] (a) at (5,2.4) {Availability and position};
\draw[edge] (m)--(e);\draw[edge] (e)--(y);\draw[edge] (p)--(y);\draw[edge] (a)--(e);
\draw[accent] (y.south)--(10,-1.5)--node[below] {New training data} (0,-1.5)--(m.south);
\node[align=center,text width=11cm] at (5,-2.7) {Interaction data reflect both preference and opportunity.\\A feedback loop can reinforce initially favored items.};
\end{tikzpicture}

\caption{A recommender changes the exposure that generates its next training sample.}\label{fig:M08_feedback}
\end{figure}

Rankings determine exposure, exposure affects observed behavior, and
that behavior enters subsequent training data. Customer preferences also
affect behavior, while availability and position affect the opportunity
to respond. Logging impressions, positions, eligible alternatives, and
the policy that selected them helps distinguish these components.
Purchase logs alone often omit the alternatives a customer could
actually have encountered.

Evaluate recommendation quality at the unit relevant to the decision.
Precision among the top recommended items and ranking metrics assess
aspects of relevance. Coverage, novelty, diversity, repeat engagement,
returns, and contribution describe other consequences. A system can
improve one metric while reducing another. Offline evaluation on
historical logs is particularly sensitive to the policy that generated
those logs. An experiment that changes exposure can supply information
that a more flexible model of the same selected data cannot supply.

\hypertarget{M08-case-study-a-temporary-campaign-can-change-the-learned-policy}{%
\subsection{Case Study: A Temporary Campaign Can Change the Learned
Policy}\label{M08-case-study-a-temporary-campaign-can-change-the-learned-policy}}

\noindent A platform may stop boosting a category while continuing to use the
behavior generated during that boost.
\href{https://arxiv.org/abs/2606.08265}{Luo, Yao, and Zhang's 2026
working paper} studies a celebrity sleep-reminder campaign on a
short-video platform. During the campaign, a ranking adjustment
increased the chance that treated users would see a reminder during the
late-night window. Exposure remained dependent on the ranking process.
The analysis uses a balanced panel of 200,000 users active on the
experiment's first day, with 100,000 drawn from each assignment group.

Table 6 reports assignment effects of 14.75\% higher late-night usage
and 2.18\% higher overall app usage relative to the respective control
means during the campaign. These are relative percentage changes in
usage, not percentage-point changes in exposure or measured effects on
sleep. The paper also reports persistent engagement changes and changes
in the content delivered afterward. The authors interpret the evidence
through discovery of previously underexposed preferences and subsequent
recommendation updates. The experiment estimates effects of the
assignment regime; attributing persistence specifically to learning
requires the additional mechanism evidence. The first-day activity
restriction also makes sample inclusion consequential when extending the
result to all assigned users.

\begin{figure}[htbp]
\centering
\begin{tikzpicture}
\node[orangebox,text width=2.5cm] (boost) at (0,0) {Temporary\\ranking boost};
\node[box,text width=2.5cm] (display) at (4.2,0) {Changed\\exposure};
\node[box,text width=2.5cm] (data) at (8.4,0) {New customer\\responses};
\draw[accent] (boost)--(display);
\draw[edge] (display)--(data);
\node[box,text width=2.5cm] (old) at (0,-3.2) {Existing\\model};
\node[orangebox,text width=2.5cm] (new) at (4.2,-3.2) {Updated\\model};
\node[box,text width=2.5cm] (future) at (8.4,-3.2) {Exposure after\\the boost ends};
\draw[edge] (old)--(new);
\draw[accent] (new)--(future);
\draw[accent] (data.south)--(8.4,-1.6)--node[above,font=\footnotesize] {Retrain on campaign data} (4.2,-1.6)--(new.north);
\draw[edge,dashed] (old.south)--(0,-4.7)--node[below,font=\footnotesize] {Freeze the model: bypass the update} (8.4,-4.7)--(future.south);
\node[anchor=west,font=\footnotesize,text=Storm] at (-1.45,1.05) {During the campaign};
\node[anchor=west,font=\footnotesize,text=Storm] at (-1.45,-2.2) {After the campaign};
\end{tikzpicture}

\caption{A temporary ranking intervention can affect later recommendations through the training data it creates.}\label{fig:M08_policy_horizon}
\end{figure}

The temporary boost ends, but its information can remain in the learned
model. Freezing the model after the campaign would interrupt that
particular channel while leaving any direct changes in customer behavior
possible. Updating it incorporates the new interactions into future
exposure. These policies can have different effects even when the
campaign content and immediate responses are identical. Separating their
contributions requires variation in the learning policy as well as the
initial intervention.

An evaluation should retain the assignment, exposure, model version,
retraining schedule, and outcomes after the intervention ends. A shared
model can transmit treatment-generated information to comparison
customers. Report the assignment unit and use uncertainty calculations
that account for repeated observations of the same unit. Measure the
intended customer outcome directly where possible, because increased
engagement cannot reveal whether a campaign achieved a goal such as less
late-night use or greater satisfaction.

\hypertarget{M08-case-study-netflix-personalization-and-incremental-engagement}{%
\subsection{Case Study: Netflix Personalization and Incremental
Engagement}\label{M08-case-study-netflix-personalization-and-incremental-engagement}}

\noindent A streaming service wants to distinguish the appeal of its catalog from
the value created by matching titles to viewers.
\href{https://www.ifo.de/en/cesifo/publications/2026/working-paper/value-personalized-recommendations-evidence-netflix}{Zielnicki,
Aridor, Bibaut, Tran, Chou, and Kallus}, in the supplied June 2026
paper, estimate a choice model using Netflix viewership and
recommendation exposure. The model incorporates heterogeneous tastes and
dependence on viewing history. Its identification treats exposure as
conditionally exogenous given the included user, title, and state
information. A separate randomized recommendation perturbation provides
diversion patterns used to assess the model.

Under that model, replacing the current recommender with matrix
factorization reduces engagement by approximately 4\%, and replacing it
with a popularity-based rule reduces engagement by approximately 12\%.
These are model-based counterfactual reductions relative to the existing
recommender, not direct A/B estimates of those complete replacements.
Their interpretation depends on the exposure condition, maintained
choice model, and counterfactual implementation.

For a subscription business, engagement can affect retention and hence
customer lifetime value, but the relationship requires evidence. A 12\%
engagement reduction does not imply a 12\% decline in revenue,
retention, or consumer welfare. Evaluate those outcomes separately,
account for substitution across titles, and compare the incremental
value with the cost of maintaining the recommendation system. The
model's experimental check is useful because it evaluates responses to
changed exposure, beyond fit to historical viewing alone.

\hypertarget{M08-case-study-relevance-variety-and-customer-goals}{%
\section{Case Study: Relevance, Variety, and Customer
Goals}\label{M08-case-study-relevance-variety-and-customer-goals}}

\noindent A customer may prefer a useful collection of products even when some
items in that collection have lower individual relevance scores.
Category repetition, past purchases, and the customer's stated goal
affect the value of the complete list. The food-recommendation studies
in M7, Section 7.12, motivate these distinctions. A simple list
objective makes the relevance-variety tradeoff explicit.

Suppose a retail customer requests three products and values variety
across categories. Assign four candidate items the relevance scores
below. The scores are comparable across these items and measure
relevance rather than purchase probability.

\begin{table}[H]
\centering\small
\begin{tabular}{@{}llr@{}}
\toprule
Item & Category & Relevance Score \\
\midrule

A & Coffee & 0.90 \\
B & Coffee & 0.85 \\
C & Tea & 0.75 \\
D & Fruit & 0.65 \\
\bottomrule
\end{tabular}
\end{table}

Choosing the three highest scores gives A, B, and C, with total
relevance 2.50 and two categories. Choosing A, C, and D gives relevance
2.30 and three categories. To make the tradeoff explicit, define a list
objective

\[
J(L)=\sum_{j\in L}s_j+\lambda K(L),
\]

where \(L\) is a three-item list, \(s_j\) is relevance, \(K(L)\) counts
distinct categories, and \(\lambda\) is the value placed on an
additional category in the same units as the relevance score. With
\(\lambda=0.30\), A, B, C has objective 3.10, while A, C, D has
objective 3.20. The latter list is preferred under this stated
objective. The threshold at which these two lists tie is
\(\lambda=0.20\).

The variety weight must come from an explicit objective, user
preference, or evaluation process. Category counts are also a coarse
measure: they treat all category differences equally and ignore which
products the customer already owns. A deployed system would need better
measurements and an evaluation of whether the resulting lists improve
the intended customer experience.

For an experiment, randomize eligible customers to the existing and
revised recommendation policies. Choose a primary outcome such as
contribution per assigned customer over a fixed horizon, alongside
measures of customer satisfaction and unwanted repetition. Analyze
assignment rather than only customers who clicked a recommendation.
Clicking occurs after the policy is assigned and can differ between
policies. Preserve the list shown and the alternatives eligible at each
impression so the mechanism can be examined.

\hypertarget{M08-engagement-considered-preferences-and-the-objective}{%
\subsection{Engagement, Considered Preferences, and the
Objective}\label{M08-engagement-considered-preferences-and-the-objective}}

\noindent A customer can click an item and later wish it had not been recommended.
\href{https://doi.org/10.1287/mnsc.2022.03683}{Kleinberg, Mullainathan,
and Raghavan} model inconsistent preferences to explain why behavior
alone can be insufficient for a platform seeking to improve user
utility. The problem can arise even when the platform wants to serve the
user: the observed action may reflect an immediate impulse that differs
from a considered preference. A useful evaluation therefore combines
behavior with explicit goals, satisfaction, and measures of regretted
use, with their measurement limitations stated.

\href{https://arxiv.org/abs/2107.10939v1}{Stray and coauthors (2021)}
discuss recommender design around objectives such as well-being,
diversity, and user control. For a retailer, a concrete formulation
could maximize expected contribution subject to customer-selected
spending limits, a maximum repetition rate, and a minimum satisfaction
criterion. These constraints make the acceptable tradeoffs explicit. A
weighted sum is another choice, but its conclusions depend on the scales
and weights. Estimates of the component outcomes remain necessary under
either formulation.

A list with more categories need not better serve every customer.
Someone replenishing a familiar weekly basket may value reliability more
than novelty. Elicit the relevant goal, preserve a way to revise it, and
evaluate whether the system responds appropriately. A policy that
increases variety on average can still repeatedly ignore the preferences
of a particular group.

\hypertarget{M08-case-study-content-exposure-and-political-attitudes-are-different-outcomes}{%
\subsection{Case Study: Content Exposure and Political Attitudes Are
Different
Outcomes}\label{M08-case-study-content-exposure-and-political-attitudes-are-different-outcomes}}

\noindent \href{https://doi.org/10.1093/pnasnexus/pgaf062}{Milli and coauthors
(2025)} conduct a preregistered audit with 806 Twitter users, comparing
posts selected by engagement ranking with recent posts from accounts
each user follows. They collect item-level stated preferences and
perceptions. Engagement ranking selects more emotionally charged and
out-group-hostile material, and the political posts selected are not
preferred by users on the stated-preference measure. The outcomes
describe selected content and users' evaluations within the audit's
exposure window.

\href{https://doi.org/10.1073/pnas.2318127122}{Liu and coauthors (2025)}
conduct four experiments with nearly 9,000 participants on a
YouTube-like interface. Manipulating recommendation sequences produces
limited short-term polarization effects in the studied settings. The
intervention, exposure duration, and attitude outcomes differ from the
Twitter audit. The studies therefore address different links between
recommendation, consumption, perception, and opinion. A change in the
composition of displayed content does not determine the magnitude of an
attitude effect.

For a managerial comparison, state the population, assignment, exposure
window, outcome, and uncertainty before combining findings across
settings. A nonsignificant attitude effect supplies no general assurance
of zero long-run effects, just as a difference in displayed content
supplies no estimate of long-run polarization. The relevant evidence is
the estimated effect and its uncertainty for the intervention actually
evaluated.

\hypertarget{M08-platform-incentives-and-preference-changes}{%
\subsection{Platform Incentives and Preference
Changes}\label{M08-platform-incentives-and-preference-changes}}

\noindent The platform's return from a recommendation can differ from the
customer's value. \href{https://doi.org/10.1111/1756-2171.12481}{Chen
and Tsai (2024)} examine Amazon's dual role as retailer and intermediary
using product-level recommendation data and variation from Amazon
stockouts. The study finds that Amazon's seller status affects the
receipt of ``Frequently Bought Together'' recommendations. The design
uses within-product changes, with accompanying controls and checks, to
distinguish seller identity from stable product appeal. Its causal
interpretation depends on the treatment of other changes accompanying
stockouts.

An analyst evaluating a marketplace should record whether an item is
sold by the platform, sponsored, or commission-bearing, alongside
availability and customer relevance. Contribution should use the
appropriate margin or fee for each item. A common ranking score can
conceal those different incentives. A policy change may improve platform
profit while reallocating sales among sellers, so report effects at the
customer, seller, and platform levels relevant to the decision.

\href{https://arxiv.org/abs/2210.05662v2}{Zhu and coauthors} develop a
benchmark for distinguishing better matching from changes in users'
preferences under a recommendation process. Their experiments use
synthetic and real-world datasets within the proposed framework. They
motivate measuring preference change separately from click-through
performance. They do not establish that every online performance gain
represents manipulation. A customer may learn to appreciate an
unfamiliar product; whether that change serves the customer requires a
stated criterion and evidence beyond a higher click count.

\hypertarget{M08-marketing-mix-models-combine-carryover-and-saturation}{%
\section{Marketing-Mix Models Combine Carryover and
Saturation}\label{M08-marketing-mix-models-combine-carryover-and-saturation}}

\noindent A marketing-mix model relates aggregate outcomes to media activity and
other determinants across time or markets. Its appeal is practical: many
firms observe weekly sales and channel spending even when individual
customer paths are unavailable. The analytical challenge is that media
decisions often respond to expected demand. Advertising, promotions,
holidays, distribution, and competitor actions can move together.

A simple carryover variable is

\[
a_{ct}=s_{ct}+\rho_c a_{c,t-1},\qquad 0\leq\rho_c<1,
\]

where \(s_{ct}\) is spending in channel \(c\) during period \(t\),
\(a_{ct}\) is its accumulated exposure proxy, and \(\rho_c\) controls
geometric decay. This unnormalized definition gives the same units as
spending, but its long-run level under constant spending is
\(s_c/(1-\rho_c)\). Some software normalizes the weights. Coefficients
and interpretations must follow the selected convention.

If one spending pulse of 100 occurs at period zero and \(\rho=0.5\), the
subsequent accumulated values are 100, 50, 25, and 12.5. A pulse can
therefore coincide with later sales even when current spending is zero.
A seasonal demand pattern can also create persistence. The model needs
enough information and a credible design to distinguish these
explanations.

Saturation describes declining additional response at higher exposure. A
saturating response curve is \(R(a)=A[1-\exp(-a/b)]\), where \(A\) is
the limiting response and \(b>0\) controls how quickly the curve
approaches that limit. A basic outcome model may combine channel
response curves, controls, trend, seasonality, and an error term. The
fitted components are model-dependent; their causal interpretation
requires assumptions about how spending was chosen and which demand
determinants were measured.

The official
\href{https://developers.google.com/meridian/docs/basics/meridian-introduction}{Meridian
documentation} describes a modern framework for marketing-mix modeling.
The
\href{https://www.pymc-marketing.io/en/latest/notebooks/mmm/mmm_example.html}{PyMC
Marketing examples} provide another implementation route. Use those
materials to inspect definitions, priors, and diagnostics. A software
package cannot supply a missing source of identifying variation.
Experiments and external evidence can help calibrate a model, subject to
whether their populations, interventions, and time horizons match the
intended use.

\hypertarget{M08-case-study-allocate-the-next-dollar-using-marginal-returns}{%
\section{Case Study: Allocate the Next Dollar Using Marginal
Returns}\label{M08-case-study-allocate-the-next-dollar-using-marginal-returns}}

\noindent Consider two channels with imposed contribution response curves

\[
R_A(s_A)=100[1-\exp(-s_A/20)],\qquad
R_B(s_B)=80[1-\exp(-s_B/40)].
\]

Both spending and contribution are measured in thousands of dollars over
the same horizon. These response functions omit cross-channel
interactions and are assumed known. At a total budget of 60, suppose the
current allocation is 40 to A and 20 to B. The question is whether
moving a small amount between channels increases contribution while
holding the total budget fixed.

Average contribution per unit spent is \(R_c(s_c)/s_c\). At the current
allocation, it is approximately 2.16 for A and 1.57 for B. Marginal
contribution is the derivative of the response curve. It is
approximately 0.68 for A and 1.21 for B. A has the higher average
return, while B offers the greater return on the next small unit of
spending.

\begin{figure}[htbp]
\centering
% Original synthetic teaching figure; see corresponding module.
\begin{tikzpicture}
\begin{axis}[xmin=0,xmax=65,ymin=0,ymax=110,xlabel={Channel spending (thousands of dollars)},ylabel={Contribution (thousands of dollars)},legend style={at={(.5,1.08)},anchor=south,legend columns=2},ymajorgrids=true]
\addplot[IlliniBlue,very thick,domain=0:60,samples=120] {100*(1-exp(-x/20))};
\addplot[IlliniOrange,very thick,domain=0:60,samples=120] {80*(1-exp(-x/40))};
\legend{Channel A,Channel B}
\addplot[IlliniBlue,dashed,domain=30:50] {86.4665+.676676*(x-40)};
\addplot[IlliniOrange,dashed,domain=10:30] {31.4775+1.21306*(x-20)};
\addplot[only marks,mark=*,IlliniBlue,mark size=3pt] coordinates {(40,86.4665)};
\addplot[only marks,mark=square*,IlliniOrange,mark size=3pt] coordinates {(20,31.4775)};
\node[anchor=south east,font=\scriptsize] at (axis cs:40,93) {At 40: slope 0.68};
\node[anchor=north west,font=\scriptsize] at (axis cs:22,28) {At 20: slope 1.21};
\end{axis}
\end{tikzpicture}

\caption{Average returns can favor one channel while marginal returns favor the other.}\label{fig:M08_marginal_returns}
\end{figure}

At the current allocations, B has the steeper response curve. Moving a
small amount from A to B therefore raises total contribution. The
comparison depends on the slopes at the current allocation, whereas
average returns summarize all spending up to that point. The budget is
held fixed, so the transfer does not change total spending expense.

For a differentiable, concave response with an interior optimum and no
other constraints, marginal returns are equal across channels. Here the
condition is

\[
5\exp(-s_A/20)=2\exp(-s_B/40),\qquad s_A+s_B=60.
\]

Solving gives \(s_A\approx32.22\) and \(s_B\approx27.78\). Total
contribution rises from approximately 117.94 to 120.09, or about 2.14
thousand dollars. If total budget were also a choice, marginal
contribution would be compared with the marginal cost of spending;
exhausting a fixed budget and choosing its size are different
optimization problems.

A real decision should also examine parameter uncertainty, operational
minimums, saturation outside observed spending ranges, delayed effects,
and interactions across channels. Report how the recommended allocation
changes under plausible response curves. A sharply stated optimum based
on weakly identified slopes can conceal substantial decision
uncertainty. Prefer an allocation that is supported by observed
variation and specify an experiment that would resolve the most
consequential uncertainty.

\hypertarget{M08-generative-ai-changes-the-cost-of-producing-and-checking-work}{%
\section{Generative AI Changes the Cost of Producing and Checking
Work}\label{M08-generative-ai-changes-the-cost-of-producing-and-checking-work}}

\noindent Generative AI can assist with code, text classification, documentation,
and communication. The business value depends on the complete workflow.
Faster generation can be offset by additional review, integration
failures, or errors that reach downstream decisions. Measure the output
that matters: a correct, reproducible analysis completed with an
appropriate understanding of its limitations.

Separate several uses. A coding assistant produces or edits a program
that can be inspected and executed. An annotation model produces labels
that require a measurement evaluation. A retrieval system selects source
passages and may generate a summary that needs factual support. A
content system produces customer-facing material whose effect should be
evaluated in its deployment context. Each use has a different unit of
quality and a different failure mechanism.

For coding, begin with an analytical specification and a small example
with a known result. Ask the assistant to implement a limited step,
explain the data transformations, and identify assumptions. Inspect the
code before running it. Then compare intermediate outputs with
independent checks. A program that executes successfully can still use
an incorrect denominator, join on the wrong key, or include future
information. Those errors require substantive checks rather than syntax
repair.

For retrieval and summarization, maintain the connection between a claim
and its source. A citation should support the particular statement next
to it. Ask the model to distinguish an observed result from a proposed
explanation, and check that distinction against the original text. If a
passage or file is unavailable, state that limitation. Fluent
reconstruction is not a substitute for access to the evidence.

\hypertarget{M08-case-study-evaluate-an-ai-assisted-analytics-workflow}{%
\section{Case Study: Evaluate an AI-Assisted Analytics
Workflow}\label{M08-case-study-evaluate-an-ai-assisted-analytics-workflow}}

\noindent Suppose a firm wants to know whether access to a coding assistant
improves analysts' work on customer-retention tasks. Define treatment as
access under a specified protocol and version, and define an outcome
that incorporates both completion and validity. A useful primary outcome
could be the fraction of assigned tasks meeting a prespecified
analytical rubric within a fixed time budget. Secondary outcomes could
include completion time, independent error severity, and a later
explanation task completed without assistance.

Randomize analysts or teams according to the likely spillovers. Repeated
tasks by the same analyst require attention to learning, carryover, and
dependence. If a crossover design is used, randomize order and consider
whether experience in the first condition changes performance in the
second. A task-level comparison with no accounting for analyst ability
or task difficulty can confound the treatment effect with who received
which work.

Quality review should use a common rubric and, where feasible, reviewers
unaware of treatment assignment. The rubric can check observation units,
joins, leakage, estimation, uncertainty, and interpretation. Record the
assigned task, the eligible inputs, the tool configuration, and the
evaluation procedure. Avoid interpreting a favorable result for one task
type as a universal effect across all analytical work.

Suppose a manual workflow uses 90 minutes to produce an analysis and 20
minutes to review it, totaling 110 minutes. An assisted workflow uses 45
minutes to produce a draft and 40 minutes to review it, totaling 85
minutes. The reduction in total time is 25 minutes, or about 22.7\%,
while the drafting-time reduction is 50\%. The total-time comparison is
informative only if the completed outputs meet comparable quality
requirements.

\begin{figure}[htbp]
\centering
% Original synthetic teaching figure; see corresponding module.
\begin{tikzpicture}[x=.095cm,y=1cm]
\node[anchor=east] at (-2,1.5) {Manual};\node[anchor=east] at (-2,0) {AI-assisted};
\fill[IlliniBlue] (0,1.1) rectangle (90,1.9);\fill[IlliniOrange] (90,1.1) rectangle (110,1.9);
\fill[IlliniBlue] (0,-.4) rectangle (45,.4);\fill[IlliniOrange] (45,-.4) rectangle (85,.4);
\node[text=white] at (45,1.5) {Produce: 90};\node[text=IlliniBlue,font=\scriptsize,align=center] at (100,1.5) {Review:\\20};
\node[text=white] at (22.5,0) {Produce: 45};\node[text=IlliniBlue] at (65,0) {Review: 40};
\node[anchor=west] at (112,1.5) {110};\node[anchor=west] at (87,0) {85};
\draw[edge] (0,-1)--(120,-1) node[right] {Minutes};
\foreach \x in {0,30,60,90} {\draw[IlliniBlue] (\x,-1.05)--(\x,-.95);\node[anchor=north,font=\scriptsize] at (\x,-1.1) {\x};}
\node[align=center,text width=12cm] at (58,-2.2) {Drafting time falls 50\%; total time falls 22.7\%.\\Task and completed-work quality held fixed.};
\end{tikzpicture}

\caption{The relevant time saving includes generation, review, and correction.}\label{fig:M08_workflow_time}
\end{figure}

This comparison holds the task and required quality fixed. Review is
part of production because an undetected error can alter a marketing
decision. In empirical work, report the completion and quality outcomes
alongside time so that faster abandonment or lower-quality work is not
counted as an equivalent improvement.

\hypertarget{M08-integrated-case-a-retail-subscription-launch}{%
\section{Integrated Case: A Retail Subscription
Launch}\label{M08-integrated-case-a-retail-subscription-launch}}

\noindent A retailer is considering a subscription service with recommendations
and delivery benefits. Begin with the decision: which offer to launch,
to which eligible customers, at what price, with what evidence of
incremental contribution? Define the evaluation horizon, capacity
constraints, and the alternatives to launch.

Construct the customer and order tables with explicit keys and dates.
Describe customer activity and purchasing cohorts, keeping unavailable
future retention separate from observed inactivity. Use regression to
summarize conditional patterns while examining confounding. Design a
conjoint exercise to learn stated tradeoffs between delivery,
assortment, and price, then assess whether its predictions generalize to
an actual offer.

Use segmentation to organize meaningful differences in needs and
behavior, and calculate customer value from contribution rather than
revenue. Analyze reviews to identify the service attributes customers
mention, checking annotation quality against human labels. Develop a
prediction model using only information available at the targeting date.
Its forecasts can help plan capacity, while a randomized offer evaluates
incremental take-up, contribution, and retention.

Finally, evaluate the recommendation component and the acquisition
budget. Log exposure to recommended items, distinguish average from
marginal media returns, and use AI assistance within a documented
verification process. The management recommendation should identify the
supported action, the uncertainty that could change it, and the next
observation or experiment with the greatest decision value. A
complicated collection of models is useful only if these pieces support
a coherent choice.

\hypertarget{M08-practice-and-further-reading}{%
\section{Practice and Further
Reading}\label{M08-practice-and-further-reading}}

\noindent \textbf{Recommendation exercise.} Compare all four three-item lists in
the variety example at \(\lambda=0\), 0.20, and 0.30. At zero, A, B, C
is optimal. At 0.20, it ties A, C, D. At 0.30, A, C, D is optimal.
Explain how a different category definition could change that conclusion
without changing any relevance score.

\textbf{Dynamic-policy exercise.} Design a comparison between a
temporary recommendation boost with a frozen model and the same boost
followed by retraining. Specify how training data are separated across
conditions, when outcomes are measured, and how shared content or
customers could transmit effects. Explain which comparison identifies
the effect of allowing the model to learn from campaign-period behavior.

\textbf{Evidence exercise.} Compare the MovieLens field experiment, the
Netflix model-based replacement counterfactuals, and the short-term
YouTube experiments. State the intervention and estimand for each.
Explain why their reported outcomes cannot be combined into a single
percentage measuring the value or harm of recommendation systems.

\textbf{Budget exercise.} Calculate the gain from moving one thousand
dollars from channel A to channel B at the original allocation. Compare
the exact change with the difference between the two marginal returns.
The derivative is a local approximation; the exact finite transfer also
reflects curvature. Then explain why an unobserved seasonal demand shock
could invalidate a causal interpretation of an estimated response curve.

\textbf{AI evaluation exercise.} Draft a rubric for an
assistant-generated retention analysis. Include at least one check that
detects a substantive error in code that executes successfully, such as
an order-line denominator used for average order value. Define a
hand-checkable example that exposes the error and explain how a reviewer
would determine whether it occurred.

Read Nielsen and coauthors,
\href{https://doi.org/10.1086/720453}{``Consumption Variety in Food
Recommendation''}, and Ling and coauthors,
\href{https://doi.org/10.1145/3485447.3514193}{``Following Good
Examples''}, for applications connecting recommendation objectives with
customer behavior. The official marketing-mix documentation linked above
develops implementation details. Consult the
\href{../01_Resources/10_AI_Coding_Concepts_and_Evidence/README.md}{AI-coding
evidence} and
\href{../01_Resources/11_AI_Coding_Tools_and_Workflows/README.md}{workflow
resources} for further study.

\hypertarget{M08-mhealth-familiarity-novelty-and-the-decision-objective}{%
\section{mHealth: Familiarity, Novelty, and the Decision
Objective}\label{M08-mhealth-familiarity-novelty-and-the-decision-objective}}

\noindent A manager may value introducing unfamiliar categories even when the most
accurate prediction repeats a customer's routine. That preference
changes the ranking objective. It should be stated as a decision choice
before interpreting any performance comparison. A category absent from
recent logs is novel relative to recorded history; it may be familiar to
the person outside the app.

Retain the September users, target baskets, candidate set, and 30-day
frequency rule from M7. Add a bonus \(\lambda\) to a candidate category
if it has not appeared in the user's prior 30 calendar days of logs:

\[
S_{ic}(\lambda)=f_{ic}+10^{-6}p_c+
\lambda\mathbf{1}\{c\text{ absent from the prior 30-day log}\}.
\]

Here \(f_{ic}\) is the fraction of observed history days containing
category \(c\), and \(p_c\) is global popularity used for tie
resolution. The bonus is measured in the same score units as frequency.
It is an imposed objective weight, not an estimated willingness to try a
food or a nutritional benefit. Rank all candidates by this score and
count the fraction of the ten recommendations absent from recent
history.

\begin{table}[H]
\centering\small
\begin{tabular}{@{}lrrr@{}}
\toprule
Novelty Bonus & Recall at 10 & NDCG at 10 & Novel Categories in the
List \\
\midrule

0.00 & 0.4759 & 0.5862 & 4.78\% \\
0.05 & 0.4759 & 0.5862 & 4.82\% \\
0.10 & 0.4759 & 0.5863 & 5.28\% \\
0.20 & 0.4661 & 0.5809 & 9.83\% \\
\bottomrule
\end{tabular}
\end{table}

The grid is a specified sensitivity comparison on the same 299 test
users. It is not used to select a winning policy and then claim an
independent test of that winner. At a bonus of 0.20, novelty roughly
doubles while recall falls by about 0.98 percentage points. At smaller
bonuses, changes in the list can occur among categories that are all
absent from the target, leaving recall unchanged. NDCG can move without
a recall change when relevant categories exchange positions.

\begin{figure}[htbp]
\centering
\begin{tikzpicture}
\begin{axis}[xlabel={Novel Categories in Top Ten (Percent)},ylabel={Mean Recall at Ten},xmin=4,xmax=11,ymin=.462,ymax=.482,ytick={.465,.470,.475,.480},yticklabel style={/pgf/number format/fixed,/pgf/number format/precision=3,/pgf/number format/fixed zerofill},grid=major]
\addplot[IlliniBlue,mark=*,mark size=3pt] coordinates {(4.78261,.4758957)(4.81605,.4758957)(5.28428,.4758957)(9.83278,.4660707)};
\node[anchor=south west,font=\footnotesize] at (axis cs:4.8,.477) {Bonuses 0, 0.05, 0.10};
\node[anchor=south east,font=\footnotesize] at (axis cs:9.8,.467) {Bonus 0.20};
\end{axis}
\end{tikzpicture}

\caption{A novelty bonus changes list composition before it substantially changes basket recovery.}\label{fig:M08_mhealth_tradeoff}
\end{figure}

The first three points lie close together in recall even though novelty
rises. The fourth trades some recovery of the recorded basket for
greater novelty. This is an offline relationship between a ranking rule
and historical labels. It does not measure whether customers prefer the
new list, accept the suggestions, or change diet.

A nutrition objective would need a validated nutrient database,
appropriate portions, restrictions supplied by the customer, and an
outcome definition. Source categories and calorie-share variety alone do
not establish healthy choices. The data-cleaning distinctions in M2
therefore affect what the manager can optimize. Dropping low-calorie
days might improve an apparent fit by selecting more complete logs,
while also excluding the customers for whom the interface is hardest to
use.

A concrete next decision is whether to run the M6 user-randomized
comparison, using a prespecified novelty setting and separate measures
of logging, acceptance, and any validated nutrition outcome. The
experiment should preserve the assignment population and record
exposure. An AI assistant can implement the frozen ranking rule and
audit category descriptions, while unsupported health claims remain
outside the factual explanation supplied with the recommendation.

\FloatBarrier
\Needspace{8\baselineskip}
Python source:
\href{../03_Data_and_Code/mHealth/recurring_examples/m08.py}{m08.py}.
From its directory, run \texttt{python\ m08.py}.

\begin{Verbatim}[breaklines,breakanywhere,fontsize=\small]
"""Measure an imposed novelty-relevance tradeoff on the same September baskets."""

import json
from common import OUTPUT, load_inputs, save_result
from ranking import candidate_categories, evaluate

food, daily, baskets = load_inputs()
selection = json.loads((OUTPUT / "M07.json").read_text())
windows = {"30_days": 30, "90_days": 90, "all_prior_days": None}
window = windows[selection["chosen_window"]]
candidates = candidate_categories(baskets)
# This grid is a specified sensitivity analysis, not a test-set model search.
weights = [0.0, 0.05, 0.10, 0.20]
tradeoff = [
    {
        "novelty_weight": weight,
        **evaluate(baskets, "2016-09-01", candidates, window, novelty_weight=weight),
    }
    for weight in weights
]
assert (
    abs(
        tradeoff[0]["recall_at_10"]
        - selection["test"]["personal_frequency"]["recall_at_10"]
    )
    < 1e-12
)
result = {"history_window": selection["chosen_window"], "tradeoff": tradeoff}
save_result("M08", result)
print(json.dumps(result, indent=2))
\end{Verbatim}

\hypertarget{M08-python-calculations-for-the-synthetic-cases}{%
\section{Python Calculations for the Synthetic
Cases}\label{M08-python-calculations-for-the-synthetic-cases}}

\noindent The inputs below are imposed teaching values. Running the program
creates the labeled synthetic CSV files and prints the corresponding
calculations. They are separate from the mHealth estimates. Keep the
script beside \texttt{example\_helpers.py} in
\texttt{03\_Data\_and\_Code/Lecture\_Examples}.

\FloatBarrier
\Needspace{8\baselineskip}
Python source:
\href{../03_Data_and_Code/Lecture_Examples/m08_synthetic.py}{m08\_synthetic.py}.
From its directory, run \texttt{python\ m08\_synthetic.py}.

\begin{Verbatim}[breaklines,breakanywhere,fontsize=\small]
"""Synthetic numerical examples for M8; run from any working directory."""

import itertools
import json
import math
from example_helpers import write_csv, ols_slope, shares

r = {}
items = {
    "A": ("Coffee", 0.90),
    "B": ("Coffee", 0.85),
    "C": ("Tea", 0.75),
    "D": ("Fruit", 0.65),
}
list_values = {}
for weight in [0, 0.2, 0.3]:
    list_values[str(weight)] = {
        "".join(choice): sum(items[j][1] for j in choice)
        + weight * len({items[j][0] for j in choice})
        for choice in itertools.combinations(items, 3)
    }
ra = lambda s: 100 * (-math.expm1(-s / 20))
rb = lambda s: 80 * (-math.expm1(-s / 40))
sa = (40 * math.log(2.5) + 60) / 3
sb = 60 - sa
r["M08"] = {
    "recommendation_objectives": list_values,
    "carryover": [100 * 0.5**t for t in range(4)],
    "average_A": ra(40) / 40,
    "average_B": rb(20) / 20,
    "marginal_A": 5 * math.exp(-40 / 20),
    "marginal_B": 2 * math.exp(-20 / 40),
    "current_total_contribution": ra(40) + rb(20),
    "optimal_spend_A": sa,
    "optimal_spend_B": sb,
    "optimal_total_contribution": ra(sa) + rb(sb),
    "contribution_gain": ra(sa) + rb(sb) - ra(40) - rb(20),
    "one_unit_transfer_gain": ra(39) + rb(21) - ra(40) - rb(20),
    "manual_total_minutes": 110,
    "assisted_total_minutes": 85,
    "time_reduction_fraction": 25 / 110,
}
assert math.isclose(5 * math.exp(-sa / 20), 2 * math.exp(-sb / 40))
# Independent grid check of the constrained optimum, to the grid's precision.
grid_best = max((ra(i / 100) + rb(60 - i / 100), i / 100) for i in range(6001))
assert abs(grid_best[1] - sa) < 0.01
write_csv(
    "M08_recommendation_synthetic.csv",
    ["item", "category", "relevance_score"],
    [(key, *value) for key, value in items.items()],
)
print(json.dumps(r, indent=2))
\end{Verbatim}

\hypertarget{M08-evidence-checks-for-generated-marketing-explanations}{%
\section{Evidence Checks for Generated Marketing
Explanations}\label{M08-evidence-checks-for-generated-marketing-explanations}}

\noindent An analytics manager may accept a fluent business review whose numbers
and interpretation describe different populations. Structured output
helps software locate a metric, period, value, unit, and denominator.
Each field also needs evidence. A correctly formatted statement can
contain the wrong number or describe a causal effect that the underlying
analysis never estimated.

Use the fixed-cohort mHealth review from Section 1.10 as the evidence
source. June contains 6,331 observed user-days out of 10,200 eligible
calendar user-days, giving a logging fraction of 0.620686. A structured
claim must identify that month, denominator, unit, and descriptive
interpretation. Four constructed candidates retain the same field
structure: one reproduces the supported result, one adds 0.10 to its
value, one substitutes the number of logging users for the denominator,
and one labels the result a causal health benefit.

\begin{table}[H]
\centering\small
\begin{tabular}{@{}lll@{}}
\toprule
Constructed Candidate & Format & Evidence Check \\
\midrule

Supported Result & Valid Fields & Accepted \\
Changed Number & Valid Fields & Value Mismatch \\
Changed Denominator & Valid Fields & Denominator Mismatch \\
Causal Health Claim & Valid Fields & Unsupported Interpretation \\
\bottomrule
\end{tabular}
\end{table}

All four candidates pass the same basic structural requirement. Only the
supported claim passes the evidence checks. They are deliberate teaching
inputs, not observed outputs from a particular LLM. The code makes no
external model request and sends no customer data to an AI service. Its
checks concern the declared fields. Free-form prose, omitted
qualifications, and advice require a separate substantive review; an
``interpretation'' field cannot validate an unrestricted paragraph.

\begin{figure}[htbp]
\centering
\begin{tikzpicture}
\node[box,text width=2.8cm,minimum height=1.4cm] (s) at (0,0) {Required Fields\\4 Candidates};
\node[box,text width=2.8cm,minimum height=1.4cm] (e) at (4.4,0) {Matching Numbers\\2 Candidates};
\node[box,text width=2.8cm,minimum height=1.4cm] (i) at (8.8,0) {Supported Scope\\1 Candidate};
\draw[edge] (s) -- (e);
\draw[edge] (e) -- (i);
\node[orangebox,text width=3.1cm] (badnum) at (3.7,-2.3) {Changed Value or Denominator};
\node[orangebox,text width=3.1cm] (badscope) at (8.8,-2.3) {Causal Health Claim};
\draw[accent] (e.south) -- (badnum.north);
\draw[accent] (i.south) -- (badscope.north);
\node[font=\footnotesize,align=center,text width=3cm] at (0,-2.3) {All Four Inputs Have Valid Field Structure};
\end{tikzpicture}

\caption{A generated statement passes through separate checks for structure, numerical evidence, and substantive interpretation before informing a decision.}\label{fig:M08_evidence_checks}
\end{figure}

\href{https://github.com/google-marketing-solutions/copycat}{Copycat}
demonstrates generation of advertising text from brand examples and
evaluates similarity and memorization. Those diagnostics concern the
generated content. Incremental sales or engagement requires an outcome
comparison under an appropriate exposure design, such as a randomized
test of the generated and existing alternatives. The Gemini and Google
Cloud marketing notebooks in the resource collection supply further
examples of generation and structured extraction; model credentials,
costs, and data permissions are separate requirements for executing
those reference notebooks.

Synthetic consumer panels raise a related distinction. The
\href{https://github.com/pymc-labs/semantic-similarity-rating}{semantic-similarity-rating
package} maps generated text to distributions over response anchors. The
mapping can standardize how generated answers are summarized. Validation
against a target population still requires human outcomes, a sampling
definition, and comparisons held out from prompt or anchor selection.
Generating more answers under the same misspecified prompt reduces
simulation noise around that prompt's distribution while preserving its
systematic difference from consumers. Agreement among generated
respondents is therefore a property to examine alongside external
validity.

A useful AI coding request is: ``Compare every numerical field with its
computed evidence; reject unknown metrics, periods, units, and
denominators; retain the analysis's descriptive interpretation; and test
a correctly formatted false claim.'' The checking rule should be written
before inspecting a favored generated answer. Report the factual error
rate and the time spent checking, alongside any reduction in drafting
time.

\FloatBarrier
\Needspace{8\baselineskip}
Python source:
\href{../03_Data_and_Code/mHealth/recurring_examples/claim_checks.py}{claim\_checks.py}.
From its directory, run \texttt{python\ claim\_checks.py}.

\begin{Verbatim}[breaklines,breakanywhere,fontsize=\small]
"""Check constructed structured claims against computed mHealth evidence."""

import json
import math
from copy import deepcopy
from common import load_food, save_result
from reporting import monthly_review


def check_claim(claim, evidence):
    required = {"metric", "period", "value", "unit", "denominator", "interpretation"}
    if not isinstance(claim, dict) or set(claim) != required:
        return ["schema"]
    rows = [row for row in evidence["periods"] if row["period"] == claim["period"]]
    if len(rows) != 1 or claim["metric"] != evidence["metric"]:
        return ["unknown metric or period"]
    row = rows[0]
    errors = []
    if type(claim["value"]) not in (int, float) or not math.isfinite(claim["value"]):
        errors.append("numeric value")
    elif not math.isclose(
        claim["value"], row["logging_fraction"], rel_tol=0, abs_tol=1e-10
    ):
        errors.append("value disagrees with evidence")
    if (
        type(claim["denominator"]) is not int
        or claim["denominator"] != row["calendar_user_days"]
    ):
        errors.append("denominator")
    if claim["unit"] != evidence["unit"]:
        errors.append("unit")
    if claim["interpretation"] != evidence["interpretation"]:
        errors.append("unsupported interpretation")
    return errors


if __name__ == "__main__":
    evidence = monthly_review(load_food())
    june = evidence["periods"][1]
    valid = {
        "metric": "logging_fraction",
        "period": "2016-06",
        "value": june["logging_fraction"],
        "unit": evidence["unit"],
        "denominator": june["calendar_user_days"],
        "interpretation": "descriptive",
    }
    candidates = {"supported": valid}
    changes = {
        "wrong_number": ("value", valid["value"] + 0.1),
        "wrong_denominator": ("denominator", june["users_with_logs"]),
        "causal_claim": ("interpretation", "causal health benefit"),
    }
    for name, (field, value) in changes.items():
        candidates[name] = deepcopy(valid)
        candidates[name][field] = value
    result = {
        "candidate_origin": "constructed teaching examples; no LLM request",
        "checks": {
            name: check_claim(claim, evidence) for name, claim in candidates.items()
        },
    }
    assert result["checks"]["supported"] == []
    assert all(result["checks"][name] for name in changes)
    save_result("claim_checks", result)
    print(json.dumps(result, indent=2))
\end{Verbatim}

\end{document}
